% !TeX program = xelatex
\documentclass[11pt,a4paper,fontset=fandol]{ctexart}
\usepackage[margin=18mm]{geometry}
\usepackage{amsmath,amssymb,bm,enumitem,hyperref}
\hypersetup{hidelinks,pdftitle={工科数学分析逐题训练·上册}}
\setlist[enumerate]{label=（\arabic*）,leftmargin=2.6em,itemsep=0.25em}
\newcommand{\dd}{\,\mathrm{d}}
\newenvironment{exercise}[2]{%
  \par\medskip\noindent\textbf{第#1题}\quad\textit{#2}\par\smallskip
}{\par\medskip}
\allowdisplaybreaks
\setcounter{secnumdepth}{-1}
\sloppy
\setlength{\emergencystretch}{4em}
\title{工科数学分析逐题训练·上册\\[0.4em]\large 计算、论证与条件辨析}
\author{}
\date{2026年10月20日}
\begin{document}
\maketitle
\section*{使用说明}
本册覆盖函数与极限、导数与微分、中值定理与导数应用、不定积分、定积分及其应用、微分方程，难度按工科数学分析期末考试控制。题目结合计算、论证、反例、条件辨析与参数讨论，使用定理时应说明关键条件。

全册共830题，按阅读顺序统一编号1至830。40节各增加两道递进题，纳入连续编号；每节先完成基本训练，再完成条件与综合训练。具体考试范围以授课要求为准，微分方程部分按课程实际教学选用。

常规版与iPad横屏版共用同一份题面，后者每页一道题，下方纯空白用于草稿。本次交付LaTeX源文件，按要求未编译，未附答案。
\tableofcontents
\clearpage
\section{第1章 函数与极限}
\subsection{函数}
\begin{exercise}{1}{定义域与参数约束}
求 $f(x)=\arccos(2x+1)+\log_2(3x+2)$ 的定义域。再确定实数 $a$ 的取值范围，使区间 $[a,a+\frac14]$ 全部包含于该定义域。
\end{exercise}
\begin{exercise}{2}{复合函数的定义域}
设 $f(x)=\frac{x-1}{x+1}$。求 $f(-2)$、$f(f(x))$，并写出复合函数的完整定义域；判断化简后的表达式能否在其自然定义域上与原复合函数视为同一函数。
\end{exercise}
\begin{exercise}{3}{换元与函数确定范围}
已知对一切 $x\ne0$，$f(1/x)=4/x+2x^2+1$。求 $f(t)$ 的已知表达式及其适用范围，再求 $f(x^2+1)$；说明上述条件是否确定了 $f(0)$。
\end{exercise}
\begin{exercise}{4}{反三角函数与值域}
设 $f(x)=\arcsin\frac4x$。求其定义域及 $f(-4)$、$f(8/\sqrt2)$、$f(8/\sqrt3)$，并确定其值域及各端点是否能够取到。
\end{exercise}
\begin{exercise}{5}{复合函数值域}
求 $f(x)=\sin\frac{\pi x}{2(1+x^2)}$ 的值域及取到最大值、最小值的全部自变量，并说明为何不能直接把正弦函数的值域作为答案。
\end{exercise}
\begin{exercise}{6}{反三角函数的主值}
设 $f(x)=\arccos2x$、$g(x)=\arctan x$、$h(x)=\operatorname{arccot}(x/2)$，其中 $\operatorname{arccot}$ 取值于 $(0,\pi)$。求 $f(\sqrt3/4)g(-1)h(2)$；分别写出三个函数的定义域和值域。
\end{exercise}
\begin{exercise}{7}{非单射换元与定义域}
已知对一切非零实数 $x$，$f(x+1/x)=\frac{4x^2}{x^4+1}$。求条件所能确定的 $f(t)$ 及 $t$ 的范围，并验证 $x$ 与 $1/x$ 给出的函数值一致；说明能否据此确定 $f(0)$。
\end{exercise}
\begin{exercise}{8}{分段函数与奇偶性}
设 $f(x)=\begin{cases}\frac{x}{\pi}|\sin x|,&|x|<\pi/3,\\0,&|x|\ge\pi/3.\end{cases}$ 求 $f(\pi/6)$、$f(-\pi/4)$、$f(2)$；判断其奇偶性，并讨论 $x=\pm\pi/3$ 两点处能否通过重新赋值使函数连续。
\end{exercise}
\begin{exercise}{9}{反解复合函数的分支限制}
设 $f(t)=\tan2t$，$f(g(x))=x^2-1$，且 $|g(x)|\le\pi/8$。求 $g(x)$ 及其最大定义域，说明范围限制在确定反函数分支和定义域时各起什么作用。
\end{exercise}
\begin{exercise}{10}{最大值函数的分段表达}
设 $\varphi(x)=\max\{x,x^2,1\}$。求定义域、分段表达式以及 $\varphi(-2)$、$\varphi(-1)$、$\varphi(0)$；确定该函数的最小值与全部取值点。
\end{exercise}
\begin{exercise}{11}{复合次序与分段边界}
设 $f(x)=\begin{cases}-1,&x<0,\\1,&x\ge0,\end{cases}$、$g(x)=2x-1$。分别写出 $f(g(x))$、$g(f(x))$ 的分段表达式，并求两者相等的全部 $x$。
\end{exercise}
\begin{exercise}{12}{复合映射与值域约束}
设 $f(x)=\begin{cases}2-x,&x<0,\\2+x,&x\ge0,\end{cases}$、$g(x)=\begin{cases}x,&x<0,\\x^2,&x\ge0.\end{cases}$ 求 $f(g(x))$、$g(f(x))$ 的分段表达式和值域，并判断它们是否为偶函数。
\end{exercise}
\begin{exercise}{13}{分段函数的四则运算}
设 $f(x)=\begin{cases}2x+1,&x<0,\\x^2+1,&x\ge0,\end{cases}$、$g(x)=\begin{cases}\sin x,&x<\pi,\\\cos x,&x\ge\pi.\end{cases}$ 写出 $f+g$、$fg$ 的分段表达式，并分别讨论二者在 $0$ 与 $\pi$ 处的连续性。
\end{exercise}
\begin{exercise}{14}{反函数及端点范围}
设 $f(x)=\frac{1-\sqrt{1+x}}{1+\sqrt{1+x}}$。确定定义域和值域，证明该函数可逆并求反函数；写出反函数的定义域，解释原函数值域中未包含的端点。
\end{exercise}
\begin{exercise}{15}{奇偶性的定义域条件}
设 $a>0$。分别判断下列函数的奇偶性，写明定义域，并对非奇非偶的函数给出依据：$ (x^2+1)\sin^3x$，$\ln(x+\sqrt{x^2+a^2})$，$\frac{2^{-x}-2^x}{2^{-x}+2^x}$，$x^2\sec x+\csc^2x$。讨论第二个函数在何种 $a$ 下成为奇函数。
\end{exercise}
\begin{exercise}{16}{抽象函数复合的奇偶性}
设 $f,g:\mathbb R\to\mathbb R$，$f$ 为奇函数，$g$ 为偶函数。判断 $f\circ f$、$f\circ g$、$g\circ f$、$g\circ g$ 必然具有的奇偶性并证明。说明“偶函数一定不是奇函数”是否成立。
\end{exercise}
\begin{exercise}{17}{单调性的直接证明}
不用导数，证明 $f(x)=x^3+2x+3$ 在 $\mathbb R$ 上严格单调；进一步确定参数 $a$ 的全部取值，使 $x^3+ax+3$ 仍在 $\mathbb R$ 上严格递增。
\end{exercise}
\begin{exercise}{18}{区间上的有界性}
分别判断 $f(x)=\frac1{x-2}$ 在 $(2,3)$ 和 $(3,+\infty)$ 上是否有界；有界时求上确界与下确界，无界时按定义构造使函数值任意大的点。
\end{exercise}
\begin{exercise}{19}{最小正周期}
求 $f(x)=2+\cos^2(2x-1)$ 的最小正周期，并证明不存在更小的正周期。讨论常数项与相位平移是否影响该结论。
\end{exercise}
\begin{exercise}{20}{取整函数与周期}
记 $\lfloor x\rfloor$ 为不超过 $x$ 的最大整数。对 $f(x)=x-\lfloor x\rfloor$，确定其值域、最小正周期，并判断有界性、单调性和奇偶性；写出 $f(-x)$ 与 $f(x)$ 在非整数点的关系。
\end{exercise}
\begin{exercise}{21}{函数关系与反函数}
设 $f:(0,+\infty)\to\mathbb R$ 满足 $f(x)+2f(1/x)=x$。求 $f(x)$；选取包含 $1$ 的最大区间，使 $f$ 在该区间上一一对应，并求其反函数及定义域。
\end{exercise}
\begin{exercise}{22}{周期函数的有界条件}
设 $f:\mathbb R\to\mathbb R$ 以 $T>0$ 为周期。证明 $f$ 在 $\mathbb R$ 上有界当且仅当它在 $[0,T)$ 上有界；构造一个无界周期函数，并说明“周期函数必有界”遗漏了什么条件。
\end{exercise}
\subsection{数列的极限}
\begin{exercise}{23}{数列极限定义}
用 $\varepsilon$--$N$ 定义证明 $\lim_{n\to\infty}\frac{n}{2n+1}=\frac12$，给出可取的 $N(\varepsilon)$；再求误差小于 $10^{-3}$ 时所需的最小正整数 $n$。
\end{exercise}
\begin{exercise}{24}{幂数列的参数分类}
对全部实数 $x$，讨论 $x^n$ 的有限极限、正负无穷趋势或振荡情形。对有限极限不存在的情形，用子列或绝对值增长给出理由。
\end{exercise}
\begin{exercise}{25}{实指数数列}
对全部实数 $x$，讨论 $n^x$ 在 $n\to\infty$ 时的变化。对 $x<0$ 的情形，用极限定义给出 $N(\varepsilon,x)$；说明固定 $x$ 的分类结论能否直接用于 $x_n=-1/\ln(n+1)$。
\end{exercise}
\begin{exercise}{26}{交错符号数列的定义证明}
设 $x_{2k-1}=-1/(2k)$、$x_{2k}=2^{-k}$（$k\ge1$）。判断数列的极限是否存在并用定义证明；比较两种符号的子列各自趋零的速度，说明交错符号与不同衰减速度是否阻碍收敛。
\end{exercise}
\begin{exercise}{27}{奇偶子列与收敛}
设 $x_n=\begin{cases}(n+1)/(n+2),&n\text{ 为偶数},\\1,&n\text{ 为奇数}.\end{cases}$ 求其极限，并用统一的误差估计验证；若将奇数项改为常数 $a$，求数列收敛的充要条件。
\end{exercise}
\begin{exercise}{28}{交错排列的子列判据}
设 $x_n\to-1$、$y_n\to1$，把它们交错排列成 $z_{2n-1}=x_n$、$z_{2n}=y_n$。判断 $z_n$ 是否收敛，并求它的全部子列极限；进一步证明一般情形 $x_n\to a,y_n\to b$ 下，$z_n$ 收敛当且仅当 $a=b$。
\end{exercise}
\begin{exercise}{29}{收敛性质的条件辨析}
判断并证明或举反例：$x_n\to a$ 则 $|x_n|\to|a|$；$|x_n|\to|a|$ 则 $x_n\to a$；$x_n\to a$ 则 $x_{n+1}\to a$；$x_n\to a$ 则 $1/x_n\to1/a$。对最后一个断言补充使其成立的必要限制。
\end{exercise}
\begin{exercise}{30}{周期振荡数列的子列}
证明 $x_n=(1+1/n)\sin(n\pi/4)$ 不收敛；求其全部子列极限，并说明因子 $1+1/n\to1$ 为何不能消除周期振荡。
\end{exercise}
\begin{exercise}{31}{收敛与连续变换}
设实数列 $x_n$ 满足 $x_n^2\to1$。证明 $|x_n|\to1$；给出 $x_n$ 不收敛的例子，并证明增加条件“$x_n$ 从某项起同号”后，$x_n$ 必收敛。
\end{exercise}
\begin{exercise}{32}{差分趋零的适用范围}
若 $x_n\to a$，证明 $x_{n+1}-x_n\to0$。分别构造无界数列和有界数列，说明该结论的逆命题均不成立；有界例子须给出两个极限不同的子列。
\end{exercise}
\subsection{函数的极限}
\begin{exercise}{33}{函数极限的定义证明}
用 $\varepsilon$--$\delta$ 定义证明 $\lim_{x\to1/2}\frac{1-4x^2}{2x-1}=-2$，给出明确的 $\delta(\varepsilon)$。若在 $x=1/2$ 任意补定义一个值，说明该极限是否改变。
\end{exercise}
\begin{exercise}{34}{无穷远左右方向}
分别求 $e^{2x-1}$ 与 $\operatorname{arccot}(1-x)$ 在 $x\to-\infty$、$x\to+\infty$ 时的极限或无穷趋势，其中 $\operatorname{arccot}\in(0,\pi)$；判断各函数是否存在 $x\to\infty$（即 $|x|\to\infty$）时的有限极限。
\end{exercise}
\begin{exercise}{35}{分段点的左右极限}
设 $f(x)=\begin{cases}(x+1)\sin x,&x<0,\\1-\cos x,&x\ge0.\end{cases}$ 求 $x=0$ 处的左右极限及双侧极限；若把左支改为 $(x+1)\sin x+a$，确定双侧极限存在的全部 $a$。
\end{exercise}
\begin{exercise}{36}{分段极限与定义域方向}
设 $f(x)=\begin{cases}\ln(1-x),&x<1,\\1/\sqrt{2x+1},&x\ge1.\end{cases}$ 求 $x\to1$、$x\to0$、$x\to2$ 及 $|x|\to\infty$ 时的极限或变化趋势，讨论时须区分各方向。
\end{exercise}
\begin{exercise}{37}{跳跃、可去点与无穷远极限}
设 $f(x)=\begin{cases}x+1,&x<0,\\3,&x=0,\\2-x,&0<x<1,\\0,&x=1,\\1+1/x,&x>1.\end{cases}$ 求 $x\to0^-$、$x\to0^+$、$x\to1$、$x\to2$ 及 $x\to+\infty$ 时的极限，并指出哪些结论不能用直接代入函数值获得。
\end{exercise}
\begin{exercise}{38}{局部保号性与有界性}
设 $\lim_{x\to1}f(x)/x=-1$。证明存在 $1$ 的去心邻域使 $f(x)<0$ 且有界；举例说明“任意去心邻域内均处处 $f(x)<0$”与“任意去心邻域内均有界”不由条件保证。
\end{exercise}
\begin{exercise}{39}{极限不等式与量词}
设 $f(x)\to A$、$g(x)\to B$（$x\to x_0$）。证明：若某去心邻域内 $f(x)<g(x)$，则 $A\le B$；若 $A<B$，则某去心邻域内 $f(x)<g(x)$。分别举例说明第一结论不能加强为 $A<B$，第二结论不能扩大到所有 $x$；说明极限存在是否要求 $f(x_0)$ 有定义。
\end{exercise}
\begin{exercise}{40}{无穷远振荡与海涅定理}
用两列趋于 $+\infty$ 的点证明 $x\cos x$ 没有有限极限，并证明它也不趋于绝对值无穷大；再判断 $\cos x/x$ 在同一变化过程中的极限。
\end{exercise}
\begin{exercise}{41}{局部有界性的准确范围}
用极限定义证明：若 $\lim_{x\to x_0}f(x)=A\in\mathbb R$，则 $f$ 在 $x_0$ 的某去心邻域内有界。若 $f(x_0)$ 也有定义，证明可把去心邻域改为邻域；举例说明局部有界不保证极限存在。
\end{exercise}
\begin{exercise}{42}{复合极限的条件辨析}
设 $\lim_{u\to0}f(u)=2$、$\lim_{x\to0}g(x)=0$。构造例子说明未必有 $f(g(x))\to2$；证明若 $g(x)$ 在 $0$ 的某去心邻域内处处不等于 $0$，该复合极限成立，并讨论把这一条件换成 $f(0)=2$ 的情形。
\end{exercise}
\begin{exercise}{43}{海涅定理与局部振荡}
定义 $f(x)=\sin(1/x)+ax$（$x\ne0$），其中 $a\in\mathbb R$。用两列趋于 $0$ 的非零点证明 $f$ 在 $0$ 处无极限；求 $xf(x)$ 在 $0$ 处的极限，并给出相应的 $\varepsilon$--$\delta$ 估计。
\end{exercise}
\subsection{无穷小与无穷大}
\begin{exercise}{44}{有界振荡与无穷小}
证明 $2x\sin(1/x)$ 在 $x\to0$ 时为无穷小。对参数 $p\in\mathbb R$，确定 $|x|^p\sin(1/x)$ 为无穷小的充要条件，并对不成立的参数给出趋零点列。
\end{exercise}
\begin{exercise}{45}{无界与无穷大的区别}
判断 $x\cos x$ 在 $\mathbb R$ 上是否无界，以及在 $x\to+\infty$ 时是否为无穷大（按绝对值趋于无穷解释）；分别构造点列证明，并说明这两个概念的量词差异。
\end{exercise}
\begin{exercise}{46}{变化过程与单侧趋势}
对 $f(x)=e^{1/x}$ 与 $g(x)=x\arctan x$，分别讨论 $x\to0^-$、$x\to0^+$、$x\to-\infty$、$x\to+\infty$ 时的极限或无穷趋势，指出哪些过程中它们是无穷小或无穷大。
\end{exercise}
\begin{exercise}{47}{无穷大与局部有界}
设 $f(x)=\frac{2x}{x-2}$、$g(x)=\ln(x+1)$。分别讨论 $x\to0$、$x\to2$、$x\to-1^+$、$x\to+\infty$ 时两函数的无穷小、无穷大或局部有界性质；对 $x\to2$ 须区分有符号极限与绝对值趋于无穷。
\end{exercise}
\begin{exercise}{48}{乘积趋零的充分条件}
设 $x_ny_n\to0$。判断下列断言并证明或举反例：$x_n$ 发散则 $y_n$ 收敛；$x_n$ 无界则 $y_n$ 有界；$x_n$ 有界则 $y_n\to0$；若 $x_n$ 最终非零且 $1/x_n\to0$，则 $y_n\to0$。
\end{exercise}
\begin{exercise}{49}{有界因子估计}
求 $\lim_{x\to0}\sin x(2+\cos(3/x))$ 与 $\lim_{|x|\to\infty}\frac{\arctan|2x+1|}{2x+1}$，分别写出能保证结论的绝对值估计；讨论第一个式子中的 $2$ 换成任意固定实数后结论是否改变。
\end{exercise}
\begin{exercise}{50}{无穷小的判定}
判定下列各量是否为无穷小，并给出相应极限或增长依据：$\frac{x^2}{\sqrt{x^3+1}}$（$x\to+\infty$），$\frac{x^3-1}{2x^4+x^2}(3-2\sin x)$（$|x|\to\infty$），$\frac{\sin2x}{x\sqrt x}$（$x\to0^+$），$\frac{\ln(1+x)}x$（$x\to0$）。
\end{exercise}
\begin{exercise}{51}{非负数列的极限推断}
设非负数列 $x_n\to0$、$y_n\to1$、$z_n\to+\infty$。证明 $x_n<y_n<z_n$ 最终成立，说明为何不能要求对所有 $n$ 成立；判断 $x_nz_n$ 与 $y_nz_n$ 的变化是否由条件唯一确定，未确定时给出两种不同结果。
\end{exercise}
\begin{exercise}{52}{由组合极限反推函数}
已知 $\lim_{x\to0}\left(1+2x+\frac{f(x)}x\right)=4$。求 $\lim_{x\to0}f(x)$ 及 $\lim_{x\to0}f(x)/x$，并用小量形式写出 $f(x)$；说明此条件能否确定 $f(0)$。
\end{exercise}
\begin{exercise}{53}{无界乘子与反例}
设 $\alpha(x)\to0$（$x\to0$）。证明有界函数 $b(x)$ 满足 $\alpha(x)b(x)\to0$；若仅知 $b$ 在任意去心邻域内无界，分别构造乘积趋零、趋于非零常数和无极限的例子。
\end{exercise}
\begin{exercise}{54}{倒数关系的定义证明}
设 $f$ 在 $x_0$ 的某去心邻域内非零。用定义证明 $|f(x)|\to\infty$ 当且仅当 $1/f(x)\to0$。再构造一个满足上述结论却既不趋于 $+\infty$ 也不趋于 $-\infty$ 的函数。
\end{exercise}
\subsection{极限的运算法则}
\begin{exercise}{55}{无穷大运算的条件}
设 $|f(x)|\to\infty$、$|g(x)|\to\infty$（$x\to x_0$）。判断 $f+g$ 必为无穷大、$f-g\to0$、$kf$ 必为无穷大（$k\ne0$）、$1/(f+g)\to0$ 四个断言，分别证明或举反例；再说明增加 $f,g\to+\infty$ 后哪些结论改变。
\end{exercise}
\begin{exercise}{56}{有限乘积的极限与误差}
求 $\lim_{|x|\to\infty}(1-1/x)(2+2/x)(3-3/x)$；记乘积为 $P(x)$，再求 $\lim_{|x|\to\infty}x(P(x)-6)$，解释两次求极限所需信息的差异。
\end{exercise}
\begin{exercise}{57}{直接代入的合法性}
求 $\lim_{x\to1/2}\frac{2\sin x+2\pi/x}{\cos2x}$，写明分母非零的依据；进一步讨论该表达式在 $x=0$ 处的左右变化趋势。
\end{exercise}
\begin{exercise}{58}{差式极限与有理化}
求 $\lim_{x\to1}\left(\frac1{x+1}-\frac3{x^3+1}\right)$ 与 $\lim_{n\to\infty}(\sqrt{n+2\sqrt n}-\sqrt{n-\sqrt n})$。将第二式中的 $2,-1$ 换成实数 $a,b$，求一般极限。
\end{exercise}
\begin{exercise}{59}{零比零型的消因子}
求 $\lim_{x\to2}\frac{x^2+x-6}{x^2-x-2}$ 与 $\lim_{x\to0}\frac{x^2+x}{\sqrt{1+x}-1}$；分别补定义函数在相应点的值使其连续，并说明约分、乘共轭式时应保留的定义域限制。
\end{exercise}
\begin{exercise}{60}{振荡项与最高次幂}
求 $\lim_{|x|\to\infty}\frac{2x^2+x\cos x}{x^3-x\sin x+1}$，并证明分母在 $|x|$ 充分大时非零；再求原商乘以 $x$ 后的极限。
\end{exercise}
\begin{exercise}{61}{有限和与几何级数}
求 $\lim_{n\to\infty}\sum_{k=1}^n\frac{k}{n^2+1}$ 与 $\lim_{n\to\infty}\sum_{k=1}^na^{-k}$（$a>1$）；再分类讨论第二个极限在 $a\in\mathbb R\setminus\{0\}$ 下存在有限值的全部情形。
\end{exercise}
\begin{exercise}{62}{望远镜求和与差式}
求 $\lim_{n\to\infty}\sum_{k=1}^n\ln\frac{k+1}{k+2}$ 及 $\lim_{n\to\infty}\left(\sqrt{\sum_{k=1}^{n+1}k}-\sqrt{\sum_{k=1}^nk}\right)$，区分有限极限与无穷趋势，并说明能否把第一个式子的极限移入不断增加的各项中。
\end{exercise}
\begin{exercise}{63}{乘积的对数化}
设 $f(x)=2^x$。求 $\lim_{n\to\infty}\frac1{n^2}\ln\prod_{k=1}^nf(k)$；再将 $f(x)$ 改为 $2^{x+c}$，证明该极限与固定常数 $c$ 无关。
\end{exercise}
\begin{exercise}{64}{乘积恒等式与参数范围}
设 $|a|<1$，求 $\lim_{n\to\infty}\prod_{k=0}^n(1+a^{2^k})$。先证明有限乘积恒等式，再讨论 $a=1$、$a=-1$ 两个边界参数，说明不能直接沿用内部参数的化简结论。
\end{exercise}
\begin{exercise}{65}{分式约分后的单侧极限}
讨论 $f(x)=\frac{x^2-1}{x+1}e^{1/(x+1)}$ 在 $x\to-1^-$、$x\to-1^+$ 时的变化，判断双侧极限是否存在；说明消去公因子是否消除了函数在 $-1$ 附近的全部奇异性。
\end{exercise}
\begin{exercise}{66}{根式与无穷远方向}
分别求 $\sqrt{x^2+x+3}-\sqrt{x^2-x+3}$ 在 $x\to+\infty$、$x\to-\infty$ 时的极限，判断 $|x|\to\infty$ 时是否有极限，并指出提取 $\sqrt{x^2}$ 时的注意事项。
\end{exercise}
\begin{exercise}{67}{幂指函数的连续运算}
求 $\lim_{x\to e}(1+\sin2x)^{\ln\ln x}$，验证实数幂在 $e$ 的某邻域内有定义；再把指数改为 $c+\ln\ln x$，求极限关于实数 $c$ 的表达式。
\end{exercise}
\begin{exercise}{68}{底数趋于非单位常数的幂指极限}
求 $\left(\frac{x^2+x\cos x-1}{2x^2-x\sin x}\right)^{x^2+1}$ 与 $\left(\frac{3x+2}{2x-1}\right)^{x^2}$ 在 $|x|\to\infty$ 时的极限或无穷趋势。证明底数最终为正，并用固定常数夹住底数来论证。
\end{exercise}
\begin{exercise}{69}{有限极限反求参数}
已知 $\lim_{x\to2}\frac{x^2+ax+b}{x^2-4}=3$，求 $a,b$。先证明分子在 $2$ 处必须为零，再求参数；并写出使该商在 $2$ 处连续的补定义。
\end{exercise}
\begin{exercise}{70}{对数分母的极限反问题}
设 $a>0$，求全部实数 $a,b$，使 $\lim_{x\to0}\frac{\sin2x}{\ln(a+x)}(\cos x-b)=4$。分别讨论 $a=1$ 与 $a\ne1$，验证所得参数使原商在充分小去心邻域内有定义。
\end{exercise}
\begin{exercise}{71}{无穷远根式抵消的参数}
设 $a>0$ 且 $\lim_{x\to-\infty}(4x+\sqrt{ax^2-bx-1})=1$。求 $a,b$，先比较一次增长项，再用有理化求剩余常数；检查改为 $x\to+\infty$ 时能否有相同结论。
\end{exercise}
\begin{exercise}{72}{根式渐近展开反求系数}
求常数 $a,b$，使 $\lim_{|x|\to\infty}(\sqrt{x^4+2x^2+3}-ax^2-b)=0$；再求所得差式乘以 $x^2$ 后的极限，说明前两个系数确定后尚余何阶误差。
\end{exercise}
\begin{exercise}{73}{运算法则的逆向条件}
设 $f(x)+g(x)\to3$、$f(x)g(x)\to2$（$x\to0$）。证明 $f,g$ 在某去心邻域内有界；构造两者均无极限的例子。若再知该邻域内 $f(x)\le g(x)$，求两者极限。
\end{exercise}
\begin{exercise}{74}{根式抵消的参数分类}
对实数 $a,b$，分类讨论 $\lim_{x\to+\infty}(\sqrt{x^2+ax+b}-x)$，并求 $\lim_{x\to+\infty}x(\sqrt{x^2+ax+b}-x-a/2)$。证明表达式在充分大的 $x$ 上有实数意义。
\end{exercise}
\subsection{极限存在准则与两个重要极限}
\begin{exercise}{75}{阶乘增长与夹逼}
求 $\lim_{n\to\infty}\frac{3^n}{n!}\sin\frac1n$；不用斯特林公式，证明存在常数 $C>0$，使充分大的 $n$ 满足 $3^n/n!\le C(1/2)^n$。
\end{exercise}
\begin{exercise}{76}{变上限和的统一估计}
证明 $\lim_{n\to\infty}\sum_{k=1}^n\frac{k}{n^2+2n+k}=\frac12$；给出上下界，并由此估计该和与 $1/2$ 的误差为 $O(1/n)$。
\end{exercise}
\begin{exercise}{77}{逐渐缩小相位差的夹逼}
求 $\lim_{x\to+\infty}\bigl(\cos\sqrt{x+1}-\cos\sqrt x\bigr)$。用三角恒等式和根式差估计证明，不得分别给两个振荡项取极限；再讨论将 $1$ 换成固定实数 $a$ 的情形。
\end{exercise}
\begin{exercise}{78}{递推数列的收敛证明}
设 $x_1>0$、$x_{n+1}=\frac12(x_n+4/x_n)$。证明从第二项起 $x_n\ge2$ 且单调，求其极限；再导出 $x_{n+1}-2$ 与 $x_n-2$ 的精确关系，说明初值小于 $2$ 时第一步的变化。
\end{exercise}
\begin{exercise}{79}{指数递推与不等式}
设 $x_1>0$、$x_ne^{x_{n+1}}=e^{x_n}-1$。证明递推式始终确定正数 $x_{n+1}$，且数列严格递减、收敛；求其极限。可使用 $e^t>1+t$（$t\ne0$），不得先假设极限存在。
\end{exercise}
\begin{exercise}{80}{积分数列的递推与比值极限}
设 $a_n=\int_0^1x^n\sqrt{1-x^2}\,\dd x$（$n\ge0$）。证明 $a_n$ 严格递减且 $a_n=\frac{n-1}{n+2}a_{n-2}$（$n\ge2$）；利用递推式和单调性夹逼求 $\lim_{n\to\infty}a_n/a_{n-1}$，写明正性保证的除法条件。
\end{exercise}
\begin{exercise}{81}{重要极限与无穷远估计}
设 $f(x)=\sin2x/(5x)$（$x\ne0$）。分别求 $x\to0$ 与 $|x|\to\infty$ 时的极限，并给出所用依据；说明为何同一分式在两个变化过程中需要不同的估计。
\end{exercise}
\begin{exercise}{82}{重要极限与参数}
求 $\lim_{x\to0}\frac{\tan3x}{4x}$；推广求 $\lim_{x\to0}\frac{\tan(ax)}{bx}$（$a\in\mathbb R,b\ne0$），并说明 $a=0$ 时结论是否仍适用。
\end{exercise}
\begin{exercise}{83}{反三角函数的换元极限}
不用洛必达法则，求 $\lim_{x\to0}\frac{\arcsin2x}{x}$ 和 $\lim_{x\to0}\frac{\arctan x}{2x}$；分别说明换元后自变量趋于零的依据及反三角函数的取值范围。
\end{exercise}
\begin{exercise}{84}{同一表达式的不同变化过程}
分别求 $\frac{2x-\sin3x}{2x+\sin4x}$ 在 $x\to0$ 与 $|x|\to\infty$ 时的极限。再确定实数 $a$ 的全部取值，使 $\lim_{x\to0}\frac{ax-\sin3x}{ax+\sin4x}$ 有有限值，并求该值。
\end{exercise}
\begin{exercise}{85}{复合零点的因式分解}
求 $\lim_{x\to3}\frac{\sin(x^2-2x-3)}{x-3}$；再讨论 $x\to-1$ 时同一分式的极限。将正弦的自变量因式分解，说明两个零点处的分母是否同时趋零。
\end{exercise}
\begin{exercise}{86}{余弦差与阶数}
求 $\lim_{x\to0}\frac{1-\cos2x}{x^2}$；进一步分类讨论 $\lim_{x\to0}\frac{1-\cos2x}{|x|^p}$（$p\in\mathbb R$），说明何时为非零有限值。
\end{exercise}
\begin{exercise}{87}{固定参数与变动参数}
对固定实数 $x$，求 $\lim_{n\to\infty}n^2\sin\frac{x^2+1}{n^2}$。再令 $x=n$，判断能否把固定参数的结论直接用于此时的表达式，并求该表达式除以 $n^2$ 后的极限。
\end{exercise}
\begin{exercise}{88}{幂指型极限}
求 $\lim_{|x|\to\infty}(1-2/x)^{2x+100}$，证明底数最终为正；再将指数中的 $100$ 换为 $c x$，求极限关于实数 $c$ 的表达式。
\end{exercise}
\begin{exercise}{89}{商底数的幂指极限}
求 $\lim_{|x|\to\infty}\left(\frac{x+1}{x-1}\right)^{2x}$；推广求 $\lim_{|x|\to\infty}\left(\frac{x+a}{x+b}\right)^{cx}$，其中 $a,b,c$ 为固定实数，并说明实数幂的定义条件最终成立。
\end{exercise}
\begin{exercise}{90}{二次尺度的重要极限}
求 $\lim_{|x|\to\infty}\left(\frac{2x^2-1}{2x^2+1}\right)^{x^2}$。若把指数 $x^2$ 改为 $|x|^p$，对实数 $p$ 分类讨论极限。
\end{exercise}
\begin{exercise}{91}{对数等价量与参数}
设 $k\ne0$，求 $\lim_{x\to0}\frac{\sin3x}{\ln(1+kx)}$；证明对数在 $0$ 的某邻域内有定义。再讨论当 $k>0$、$x\to+\infty$ 时该商的极限，区分局部等价量的适用过程。
\end{exercise}
\begin{exercise}{92}{取对数与符号控制}
求 $\lim_{x\to0}(1+\tan2x)^{-3/x}$，并验证底数在 $0$ 的某邻域内为正；推广求 $(1+\tan ax)^{b/x}$ 的极限，其中 $a,b\in\mathbb R$。
\end{exercise}
\begin{exercise}{93}{指数与底数的共同变化}
求 $\lim_{x\to0}(\cos x)^{-1/\sin x}$ 和 $\lim_{x\to0}\left(\frac{\tan x}{x}\right)^{1/x}$。将两式的指数分别改为 $-1/\sin^2x$ 和 $1/x^2$，再求极限，比较一阶与二阶尺度。
\end{exercise}
\begin{exercise}{94}{幂平均的局部极限}
设 $a,b>0$，求 $\lim_{x\to0}\left(\frac{a^x+b^x}{2}\right)^{1/x}$；推广到 $\left(\lambda a^x+(1-\lambda)b^x\right)^{1/x}$（$0\le\lambda\le1$），说明端点权重情形。
\end{exercise}
\begin{exercise}{95}{对数与幂的括号范围}
设 $a\ne\frac12$，求 $\lim_{n\to\infty}\ln\left[\left(\frac{n-2an-1}{n(1-2a)}\right)^n\right]$，说明对数为何最终有定义；再求幂本身的极限。
\end{exercise}
\begin{exercise}{96}{幂指极限反求参数}
求全部实数对 $(a,b)$，使 $\lim_{|x|\to\infty}\left(\frac{x^2}{(x-a)(x+b)}\right)^x=e^2$。证明所得条件充分且必要，并检查底数最终为正。
\end{exercise}
\begin{exercise}{97}{递推极限与初值}
设 $x_1\ge0$、$x_{n+1}=\sqrt{2+x_n}$。分别讨论 $0\le x_1<2$、$x_1=2$、$x_1>2$ 的单调性，证明数列收敛并求极限；当 $x_1\ne2$ 时给出 $|x_{n+1}-2|/|x_n-2|$ 的极限。
\end{exercise}
\begin{exercise}{98}{重要极限中的高阶修正}
对实数 $a,b$，求 $\lim_{x\to0}(1+ax+bx^2)^{1/x}$；当 $a=0$ 时，再求 $\lim_{x\to0}(1+bx^2)^{1/x^2}$。各情形均须说明底数在充分小邻域内的正性。
\end{exercise}
\subsection{无穷小的比较}
\begin{exercise}{99}{以根式为基准比较阶数}
在 $x\to0^+$ 时，分别求 $e^{\sqrt x}-1$、$\ln(1+\sqrt x)$、$1-\cos\sqrt x$、$\arctan\sqrt x$ 与 $\sqrt x$ 的比值极限；对不等价者给出其关于 $\sqrt x$ 的阶数与首项。
\end{exercise}
\begin{exercise}{100}{不同阶无穷小的比较}
当 $x\to0$ 时，分别求 $\arctan2x$、$\sqrt{1+x}-1$、$1-\cos2x$、$\tan x-\sin x$ 的首个非零等价项，按阶数排序并指出最高阶者；说明同阶是否必然等价。
\end{exercise}
\begin{exercise}{101}{无穷远的等价量}
求 $\lim_{|x|\to\infty}\frac{2x-3}{x^2-x+1}$；进一步求 $x\frac{2x-3}{x^2-x+1}$ 的极限，并写出原商相对于 $1/x$ 的等价关系。
\end{exercise}
\begin{exercise}{102}{指数差的消去阶数}
当 $x\to0$ 时，确定使 $e^{x\cos x^2}-e^x$ 与 $x^n$ 同阶的正整数 $n$，并求两者比值的极限；论证时须保留指数差中首先不为零的项。
\end{exercise}
\begin{exercise}{103}{三阶等价无穷小}
证明 $\tan x-x\sim x^3/3$（$x\to0$）；利用此结论求 $\lim_{x\to0}\frac{\tan(ax)-a\tan x}{x^3}$，并讨论实数 $a$ 使极限为零的全部情形。
\end{exercise}
\begin{exercise}{104}{高阶无穷小与首项}
证明 $(2^{x^2}-1)\arctan2x=o(x^2)$（$x\to0$）；再求其与 $x^3$ 的比值极限，说明“高于二阶”是否意味着“高于三阶”。
\end{exercise}
\begin{exercise}{105}{等价关系能确定的参数信息}
求全部实数 $a,b,c$，使 $\frac1{ax^2+bx+c}\sim\frac1{2x^2+1}$（$|x|\to\infty$）。若进一步要求两式之差为 $o(x^{-4})$，再求参数；说明增强前后分别约束了哪些系数。
\end{exercise}
\begin{exercise}{106}{含实参数幂的阶数比较}
当 $x\to0^+$ 时，求全部实数 $k$，使 $(\sqrt{1+x^2}-1)\ln(1+x^2)=o\bigl(x(\arcsin x)^k\bigr)$；再讨论临界参数时的比值极限。
\end{exercise}
\begin{exercise}{107}{乘积中的高阶消去}
求 $\lim_{x\to0}\frac{(\sin x-\tan x)(e-e^{\cos x})}{(2x+1)\sin^5x}$。分别写出分子两个因子的等价首项，并说明为何不能把 $\sin x$ 与 $\tan x$ 都代换成 $x$ 后相减。
\end{exercise}
\begin{exercise}{108}{等价关系反求两个参数}
设 $ab\ne0$，当 $x\to0$ 时 $x-\sin(ax)\sim x^2\ln(1-bx)$。求 $a,b$ 并验证；说明为何必须先排除分子的一阶项，再比较三阶项。
\end{exercise}
\begin{exercise}{109}{对数极限反推原函数}
设 $a>0$、$a\ne1$，且 $\lim_{x\to0}\frac{\ln(1+f(x)/\sin^2x)}{a^2-1}=A\in\mathbb R$，原式在充分小去心邻域内有定义。求 $\lim_{x\to0}f(x)/x^2$；分别讨论 $A=0$ 与 $A\ne0$ 时能否直接使用 $\ln(1+u)\sim u$。
\end{exercise}
\begin{exercise}{110}{已知高阶余项反求极限}
已知 $\lim_{x\to0}\frac{\sin6x+xf(x)}{x^3}=0$，求 $\lim_{x\to0}\frac{6+f(x)}{x^2}$；再写出 $f(x)$ 到 $x^2$ 阶的渐近表达式，并说明是否能据此确定 $f(0)$。
\end{exercise}
\begin{exercise}{111}{振荡余项的阶估计}
求 $\lim_{x\to0}\frac{4\sin x+3x^2\cos(1/x^2)}{(1+\cos x)\ln(1+2x)}$。若把分子中的 $x^2$ 改为 $|x|^p$，确定极限仍等于原值的全部实数 $p$，对边界情形说明理由。
\end{exercise}
\begin{exercise}{112}{根式差与三阶比较}
求 $\lim_{x\to0}\frac{\sqrt{1+\tan x}-\sqrt{1+\sin x}}{x\ln(1+x)-x^2}$；指出两个差式各自的首个非零阶，并验证根式和对数在充分小邻域内有定义。
\end{exercise}
\begin{exercise}{113}{指数相近时的四阶消去}
求 $\lim_{x\to0}\frac{e^{x^2}-e^{2-2\cos x}}{x^3\sin x}$；写出两个指数之差的首个非零项，并说明分别只保留指数到二阶会丢失什么信息。
\end{exercise}
\begin{exercise}{114}{对数差的等价转化}
求 $\lim_{x\to0}\frac{\ln(\sin^2x+e^{2x})-2x}{\ln(x^2+e^{3x})-3x}$。先把每个对数差改写为 $\ln(1+u)$ 型，再说明相减结构下允许等价代换的具体位置。
\end{exercise}
\begin{exercise}{115}{无穷远复合等价量}
求 $\lim_{|x|\to\infty}x\left[\sin\ln(1+4/x)-\sin\ln(1+1/x)\right]$；将 $4,1$ 换为实数 $a,b$，求一般结果，并证明两个对数最终有定义。
\end{exercise}
\begin{exercise}{116}{积分无穷小的阶数比较}
在 $x\to0^+$ 时，求 $\int_0^x(e^t-1)\,\dd t$、$\int_0^x\ln(1+\sqrt{t^2})\,\dd t$、$\int_0^x\sin(t^2)\,\dd t$、$\int_0^x\sqrt{\sin^2t}\,\dd t$ 的等价首项，按阶数排序并找出最高阶者。说明两个根式在正半轴上如何化简。
\end{exercise}
\begin{exercise}{117}{极限减去常数后的阶数}
设 $f(x)=\frac{1+x}{\sin x}-\frac1x$，$a=\lim_{x\to0}f(x)$。求 $a$；确定使 $f(x)-a$ 与 $x^k$ 同阶的正整数 $k$，并求比值极限；说明分子分母主项相消后应保留到哪一阶。
\end{exercise}
\begin{exercise}{118}{积分分母的非零有限极限}
求全部常数 $a,b,c$（$c\ne0$），使 $\lim_{x\to0}\frac{ax-\sin(bx)}{\int_0^x\ln(1+t^2)\,\dd t}=c$。先比较阶数再确定系数关系，并验证所求关系足以保证极限存在。
\end{exercise}
\begin{exercise}{119}{加减式中等价代换的限度}
已知 $\alpha\sim\beta$、$\gamma\sim\delta$（$x\to0$），说明是否必有 $\alpha-\gamma\sim\beta-\delta$，并给出反例。再求 $\lim_{x\to0}\frac{\sin x-x\cos x}{\tan x-x}$，用首项说明相消时需要保留到哪一阶。
\end{exercise}
\begin{exercise}{120}{参数控制的阶数跃迁}
设 $F_a(x)=\ln(1+x)-x+ax^2$。对所有实数 $a$，求 $F_a(x)$ 在 $x\to0$ 时的首个非零等价项，并确定其阶数；求全部 $(a,c)$，使 $F_a(x)\sim cx^3$ 且 $c\ne0$。
\end{exercise}
\subsection{函数的连续性与间断点}
\begin{exercise}{121}{差商极限与连续性}
已知 $f$ 在 $1$ 处连续，且 $\lim_{x\to1}\frac{f(x)-4}{x-1}=A\in\mathbb R$。求 $f(1)$，证明 $f$ 在 $1$ 处可导并求其导数；去掉连续性条件后，说明极限条件还可以推出哪些结论。
\end{exercise}
\begin{exercise}{122}{连续性的定义证明}
用 $\varepsilon$--$\delta$ 定义证明 $f(x)=2x^2-x+1$ 在任意实数 $x_0$ 处连续，并给出依赖 $\varepsilon,x_0$ 的明确 $\delta$；说明该估计是否可在固定闭区间上统一选取。
\end{exercise}
\begin{exercise}{123}{连续延拓及原式的定义范围}
设 $f(x)=\begin{cases}(e^{2x}-1)/\sin4x,&x\ne0,\\k,&x=0.\end{cases}$ 在 $0$ 的充分小邻域内定义。求使其在 $0$ 处连续的 $k$；列出非零分母零点，判断仅确定 $k$ 能否使原表达式在整个 $\mathbb R$ 上有定义且连续。
\end{exercise}
\begin{exercise}{124}{分段点的连续性与单侧修复}
设 $f(x)=\begin{cases}\frac{\sin2x}{\tan x}(1+x^2),&x>0,\\2,&x=0,\\\ln\sqrt{1-2x},&x<0.\end{cases}$，各支在 $0$ 的充分小邻域内考虑。判断 $f$ 在 $0$ 处的连续性；若只允许给左支加常数 $c$，求能使函数连续的全部 $c$。
\end{exercise}
\begin{exercise}{125}{振荡因子的连续延拓}
讨论 $f(x)=\begin{cases}(x-1)^2\cos\frac1{x-1},&x<1,\\(x+1)\ln x,&x\ge1\end{cases}$ 在定义域内的连续性；若把左支中的 $(x-1)^2$ 改为 $|x-1|^p$，确定在 $1$ 处连续的全部实数 $p$。
\end{exercise}
\begin{exercise}{126}{根式与反三角函数的延拓}
为 $f(x)=\frac{\sqrt{1+x^2}-\sqrt{1-x^2}}{x\arcsin x}$ 补定义 $f(0)$，使其在 $0$ 处连续；写出延拓后函数的定义域，并判断它在定义域的端点是否连续。
\end{exercise}
\begin{exercise}{127}{由数列极限定义的函数}
对 $x>0$，设 $f(x)=\lim_{n\to\infty}\frac1n\ln(e^{nx}+x^2)$。求 $f(x)$ 并讨论其连续性；若按同一极限式将定义范围扩大到 $\mathbb R$，求扩展后的分段表达式并判断 $0$ 处是否连续。
\end{exercise}
\begin{exercise}{128}{连续函数值的反推}
设 $f$ 在 $0$ 的某邻域内连续，在充分小的 $x>0$ 时原商有定义，且 $\lim_{x\to0^+}\frac{1-\cos(\sqrt x\,f(x))}{(e^{2x}-1)f(x)}=\frac12$。求 $f(0)$；证明时须同时考虑 $f(0)=0$ 与 $f(0)\ne0$，不能未经论证就约去可能趋零的量。
\end{exercise}
\begin{exercise}{129}{有理函数的间断点}
求 $f(x)=\frac{x^2-1}{x^2-3x+2}$ 的全部间断点并分类；对可去间断点作连续延拓，对其余间断点写出左右变化趋势，说明二者不能采用同一补定义方式处理。
\end{exercise}
\begin{exercise}{130}{指数分式的间断点分类}
求 $f(x)=\frac1{e^{x/(x-1)}-1}$ 的全部间断点，计算各自的左右极限或无穷趋势并分类；判断每一点能否通过补定义成为连续点。
\end{exercise}
\begin{exercise}{131}{可去间断点的参数条件}
设 $f(x)=\begin{cases}\frac{x^2\sin(1/(2x))}{2^x-1},&x<0,\\b,&x=0,\\a+\frac{\ln(1+3x)}x,&x>0.\end{cases}$ 求使 $0$ 为可去间断点的全部参数对 $(a,b)$；再求使其在 $0$ 连续的参数对，说明可去间断点与连续点对 $b$ 的要求。
\end{exercise}
\begin{exercise}{132}{点态极限与间断点}
设 $f(x)=\lim_{n\to\infty}\frac{1+x}{1+x^{2n}}$。求其分段表达式、全部连续区间及间断点类型；说明为何不能从各个近似函数连续直接推断 $f$ 连续。
\end{exercise}
\begin{exercise}{133}{幂平均极限与最大值函数}
设 $a_1,\ldots,a_m>0$，证明 $\lim_{n\to\infty}(a_1^n+\cdots+a_m^n)^{1/n}=\max_{1\le i\le m}a_i$。据此求 $f(x)=\lim_{n\to\infty}\bigl(1+x^n+(x^2/2)^n\bigr)^{1/n}$（$x\ge0$）的分段表达式，并讨论其连续性。
\end{exercise}
\begin{exercise}{134}{连续与可去间断的参数分类}
设 $f(x)=\begin{cases}\frac{\ln(1+ax^3)}{x-\arcsin x},&x<0,\\6,&x=0,\\\frac{e^{ax}+x^2-ax-1}{x\tan(x/4)},&x>0,\end{cases}$ 各支在 $0$ 的充分小邻域内考虑。求使 $f$ 在 $0$ 连续的全部 $a$，以及使 $0$ 为可去间断点的全部 $a$；给出各情形的左右极限。
\end{exercise}
\begin{exercise}{135}{参数极限生成函数}
对 $\sin x\ne0$ 的实数 $x$，设 $f(x)=\lim_{t\to x}\left(\frac{\sin t}{\sin x}\right)^{x/(\sin t-\sin x)}$，极限沿表达式有定义的 $t$ 取。求 $f(x)$，找出其定义域边界上的间断点并分类；区分 $x=0$ 与其他 $\pi$ 的整数倍。
\end{exercise}
\begin{exercise}{136}{稠密分段与连续点}
设 $f(x)=\begin{cases}x^2,&x\in\mathbb Q,\\ax+b,&x\notin\mathbb Q.\end{cases}$ 求其全部连续点关于参数 $a,b$ 的分类；给出恰有两个连续点的参数，并证明该函数不可能在任何非退化区间上连续。
\end{exercise}
\begin{exercise}{137}{振荡函数的参数延拓}
设 $f_{p,c}(x)=\begin{cases}|x|^p\sin(1/x),&x\ne0,\\c,&x=0.\end{cases}$ 对全部实数 $p,c$，讨论 $0$ 处的连续性与间断点类型；对 $p<0$，说明为什么函数在每个邻域内无界仍不构成无穷间断点的充分条件。
\end{exercise}
\subsection{连续函数的运算与初等函数的连续性}
\begin{exercise}{138}{连续运算的分母条件}
设 $f(x)=\frac{\lg(100+x^2)}{a^2+\arcsin2x}$。分别讨论 $a\ne0$ 与 $a=0$ 时 $x\to0$ 的双侧极限或单侧趋势，说明何时可直接使用连续函数商的法则。
\end{exercise}
\begin{exercise}{139}{间断性在运算中的保留}
设 $f$ 在 $\mathbb R$ 上连续且处处非零，$\varphi:\mathbb R\to\mathbb R$ 有间断点。判断 $\varphi\circ f$、$f\circ\varphi$、$\varphi^2$、$\varphi/f$ 是否必有间断点；必然成立者证明，其余给出满足全部条件的反例。
\end{exercise}
\begin{exercise}{140}{反函数连续性的条件}
设 $f$ 在闭区间 $[a,b]$ 上连续且严格递增。证明 $f^{-1}$ 在 $[f(a),f(b)]$ 上连续，端点按单侧连续理解；举例说明仅有严格递增而无连续性时，函数值域可能不是区间。
\end{exercise}
\begin{exercise}{141}{连续复合与零点集合}
设 $f$ 在 $[a,b]$ 上连续。证明 $g(x)=\max\{f(x),0\}$ 与 $h(x)=|f(x)|$ 连续；再证明若 $f$ 处处非零，则 $1/f$ 在该区间上连续且有界，并指出闭区间条件在有界性结论中的作用。
\end{exercise}
\subsection{闭区间上连续函数的性质}
\begin{exercise}{142}{三次方程的实根个数}
证明方程 $x^3-9x-2=0$ 恰有三个不同实根，并将每个根定位在一个长度为 $1$ 且互不相交的开区间内；分别说明“至少三个”与“至多三个”的依据。
\end{exercise}
\begin{exercise}{143}{零点存在与唯一性}
证明方程 $x=\frac12\cos x$ 在 $\mathbb R$ 上恰有一个实根，并给出长度不超过 $1/2$ 的含根闭区间；存在性与唯一性须分别论证。
\end{exercise}
\begin{exercise}{144}{平移差函数与零点}
设 $f$ 在 $[0,1]$ 上连续且 $f(0)=f(1)$。证明存在 $\xi\in[0,1/2]$，使 $f(\xi+1/2)=f(\xi)$；说明为何平移差函数只能先定义在 $[0,1/2]$ 上。
\end{exercise}
\begin{exercise}{145}{半无限区间的有界性}
设 $f$ 在 $[a,+\infty)$ 上连续，且 $\lim_{x\to+\infty}f(x)=L\in\mathbb R$。证明 $f$ 有界；再分别举例说明它未必取得最大值、未必取得最小值。
\end{exercise}
\begin{exercise}{146}{加权平均与介值性}
设 $f$ 在 $[a,b]$ 上连续，$a<x_1<x_2<x_3<b$。证明存在 $\xi\in(a,b)$ 使 $f(x_1)+2f(x_2)+3f(x_3)=6f(\xi)$；若三点函数值不全相同，证明可取 $\xi\in(x_1,x_3)$。
\end{exercise}
\begin{exercise}{147}{开区间端点与有界性}
设 $f(x)=\frac{x\sin(x-2)}{|x|(x-1)(x-2)^2}$。分别判断它在 $(-1,0)$、$(0,1)$、$(1,2)$、$(2,3)$ 内是否有界；有界情形用连续延拓或明确估计证明，无界情形给出导致无界的端点趋势。
\end{exercise}
\begin{exercise}{148}{等分平移与望远镜求和}
设 $f$ 在 $[0,1]$ 上连续且 $f(0)=f(1)$。证明对每个正整数 $n$，存在 $x_0\in[0,1-1/n]$ 使 $f(x_0)=f(x_0+1/n)$。可考虑平移差函数在等分点上的和，并说明 $n=1$ 的情形。
\end{exercise}
\begin{exercise}{149}{闭区间自映射的不动点}
设连续函数 $f:[a,b]\to[a,b]$，其中 $a<b$。证明至少存在一点 $\xi\in[a,b]$ 使 $f(\xi)=\xi$；若另有 $|f(x)-f(y)|\le q|x-y|$（$0\le q<1$），证明不动点唯一；举例说明仅连续时唯一性可能不成立。
\end{exercise}
\begin{exercise}{150}{严格不等式的统一余量}
设 $f,g$ 在 $[a,b]$ 上连续且处处 $f(x)<g(x)$。证明存在 $\eta>0$，使所有 $x\in[a,b]$ 均满足 $f(x)+\eta\le g(x)$；再给出开区间 $(a,b)$ 上的反例，说明逐点严格不等式未必给出统一正余量。
\end{exercise}

\clearpage
\section{第2章 导数与微分}
\subsection{导数的基本概念}
\begin{exercise}{151}{定义求导与线性主部}
用导数定义求 $f(x)=1+2x+3x^2$ 的 $f'(1)$，并把 $f(1+h)-f(1)$ 写成线性主部与余项之和，验证余项为 $o(h)$。
\end{exercise}
\begin{exercise}{152}{零因子处的导数}
设 $f(x)=x(2x+1)(3x+2)(4x+3)$。不展开多项式，求 $f'(0)$；再求 $f'(-1/2)$，说明在单零点处如何利用其余因子的连续性。
\end{exercise}
\begin{exercise}{153}{定义求导与零点因子}
设 $f(x)=x(x-2)\sqrt{\frac{3x^2+1}{2x^2+1}}$。用导数定义求 $f'(0)$、$f'(2)$，并分别写出曲线在这两点的切线方程。
\end{exercise}
\begin{exercise}{154}{补定义函数的导数}
设 $f(x)=\begin{cases}\sin^2x/(2x),&x\ne0,\\0,&x=0.\end{cases}$ 用定义求 $f'(0)$；若把 $f(0)$ 改为参数 $c$，确定在 $0$ 处连续、可导各需要什么条件。
\end{exercise}
\begin{exercise}{155}{连续因子与乘积可导}
设 $f(x)=(x^2-1)\varphi(x-1)$，$\varphi$ 在 $0$ 处连续。证明 $f$ 在 $1$ 处可导并求 $f'(1)$；构造 $\varphi$ 在 $0$ 处不可导而该结论仍成立的例子。
\end{exercise}
\begin{exercise}{156}{平移函数关系的导数}
设 $f$ 在 $0$ 处可导，$f'(0)=1/2$，且对所有实数 $x$ 有 $f(2+x)=2f(x)$。用定义求 $f'(2)$；进一步证明 $f$ 在所有偶整数点可导，并求 $f'(2n)$（$n\in\mathbb Z$）。
\end{exercise}
\begin{exercise}{157}{差商极限与点值}
设 $f$ 在 $0$ 处连续，且 $\lim_{x\to0}\frac{f(x)+3}{x}=2$。判断 $f$ 在 $0$ 处是否可导并求导数；若去掉连续性，说明能否保留可导结论，并给出反例。
\end{exercise}
\begin{exercise}{158}{缩放增量的差商}
设 $f$ 在 $x_0$ 处可导。求 $\lim_{h\to0}\frac{f(x_0+5h)-f(x_0)}h$ 与 $\lim_{h\to0}\frac{f(x_0)-f(x_0-3h)}h$；用同一方法写出一般系数 $a,b$ 下 $\lim_{h\to0}\frac{f(x_0+ah)-f(x_0+bh)}h$。
\end{exercise}
\begin{exercise}{159}{双移位差商与逆命题}
设 $f'(x_0)$ 存在，求 $\lim_{h\to0}\frac{f(x_0+4h)-f(x_0+2h)}h$。若不要求 $f$ 在 $x_0$ 连续，构造一个在该点不可导但上述极限仍存在的函数，说明差分式中没有出现点值的影响。
\end{exercise}
\begin{exercise}{160}{二次增量下的单侧信息}
设 $f$ 在 $3$ 处可导，且 $\lim_{x\to0}\frac{f(3)-f(3-3x^2)}{\ln(1+x^2)}=3$。求 $f'(3)$；说明若仅假设 $f$ 在 $3$ 处连续，给定极限能否保证双侧导数存在。
\end{exercise}
\begin{exercise}{161}{可导函数的渐近形式}
设 $f(0)=0$、$f'(0)=1$。求 $\lim_{x\to0}\frac{f(3x)-f(-5x)}x$ 与 $\lim_{x\to0}\frac{f(2x^2)-xf(x)}{x^2}$；分别写明使用 $f(t)=t+o(t)$ 时余项的变化。
\end{exercise}
\begin{exercise}{162}{可导充分条件的辨析}
设 $f(0)=0$，分别判断下列极限存在是否足以保证 $f'(0)$ 存在，证明或举反例：$\lim_{h\to0}\frac{f(3h)-f(h)}{2h}$，$\lim_{h\to0}\frac{f(h)-f(-h)}{2h}$，$\lim_{h\to0}\frac{f(h^3)}{h^3}$，$\lim_{h\to0}\frac{f(h^2)}{h^2}$。指出奇次与偶次换元提供的自变量方向有何不同。
\end{exercise}
\begin{exercise}{163}{二次内增量与已知导数}
设 $f$ 在 $1$ 处可导，$f(1)=0$、$f'(1)=2$。求 $\lim_{x\to0}\frac{\ln(1+\sin^2(3x))-f(\cos x)}{x^2}$，写出 $f(\cos x)$ 的首个非零渐近项；说明为何只知道 $f'(1)$ 已足以计算该极限。
\end{exercise}
\begin{exercise}{164}{内点导数与端点约束}
分别判断 $|x|$、$\sqrt x$、$x^2$、$x\ln x$ 在 $0$ 处是否具有通常的双侧导数。对原式在 $0$ 未定义的函数，先判断能否连续补定义，再讨论其右差商是否有有限极限。
\end{exercise}
\begin{exercise}{165}{左右导数的拼接}
设 $f(x)=\begin{cases}\cos x,&x>0,\\x^2+1,&x\le0.\end{cases}$ 求 $f'_-(0)$、$f'_+(0)$ 并判断可导性；再比较两侧二阶导数，判断 $f''(0)$ 是否存在。
\end{exercise}
\begin{exercise}{166}{指数分段函数的单侧导数}
设 $f(x)=\begin{cases}x(1+e^{1/x}),&x\ne0,\\0,&x=0.\end{cases}$ 分别讨论其在 $0$ 处的左连续性、右连续性与左右导数，区分有限单侧导数与无穷差商；说明一侧导数存在能否推出另一侧连续。
\end{exercise}
\begin{exercise}{167}{绝对值造成的尖点}
讨论 $f(x)=x^2|x-1|$ 在 $1$ 处的左右导数与可导性；若将因子 $x^2$ 换为 $(x-1)^2$，重新判断并解释零点因子的作用。
\end{exercise}
\begin{exercise}{168}{零点消除绝对值尖点}
讨论 $f(x)=(x^3-1)|x-1|$ 在 $1$ 处的可导性并求导数；进一步判断 $f'$ 在 $1$ 处是否连续、$f''(1)$ 是否存在。
\end{exercise}
\begin{exercise}{169}{可导的充要条件}
设 $\varphi$ 在 $1$ 处连续，$f(x)=|x-1|\varphi(x)$。证明 $f$ 在 $1$ 处可导当且仅当 $\varphi(1)=0$，并求此时的导数；说明证明中连续性使用在何处。
\end{exercise}
\begin{exercise}{170}{绝对值复合的全部不可导点}
求 $f(x)=(x^2+x-2)|x^2-4x|$ 的全部不可导点及其左右导数；说明绝对值内部多项式的零点是否一定造成不可导。
\end{exercise}
\begin{exercise}{171}{切线与法线}
求曲线 $y=2x+\cos x$ 在 $(0,1)$ 处的切线、法线方程；再确定该曲线是否存在水平切线，并给出依据。
\end{exercise}
\begin{exercise}{172}{切线垂直的可行性}
在相同横坐标 $x_0$ 处，曲线 $y=x^2$ 与 $y=x^3$ 的切线互相垂直。求所有 $x_0$，并检查用斜率乘积为 $-1$ 的判据时是否遗漏竖直切线情形。
\end{exercise}
\begin{exercise}{173}{指定方向的切线}
直线 $L$ 与 $2x-6y+1=0$ 垂直，且与曲线 $y=x^3+3x^2-5$ 相切。求全部可能的切点及 $L$ 的方程；说明斜率方程的重根是否表示存在两条不同切线。
\end{exercise}
\begin{exercise}{174}{切线截距与面积不变量}
证明双曲线 $xy=1$ 上任一点处的切线与坐标轴围成的三角形面积恒为 $2$；写明两截距的符号，并说明点位于第三象限时面积结论为何仍成立。
\end{exercise}
\begin{exercise}{175}{公共切线的两个条件}
求 $a,b$，使 $y=1+\ln x$ 与 $y=x^2+ax+b$ 在 $x=1$ 处具有同一条切线；分别列出点值条件和斜率条件，并求公共切线方程。
\end{exercise}
\begin{exercise}{176}{偶函数导数与倒数差商}
设 $f$ 为可导偶函数，且 $\lim_{h\to0}\frac{h}{f(3+h)-f(3)}=-2$。求曲线在 $(-3,f(-3))$ 处的切线斜率；从奇偶性定义证明所使用的导数关系。
\end{exercise}
\begin{exercise}{177}{可导蕴含连续}
设 $f$ 在 $2$ 处可导，且 $\lim_{x\to2}f(x)=5$。求 $f(2)$；证明所使用的“可导蕴含连续”结论，并构造函数说明连续不足以保证可导。
\end{exercise}
\begin{exercise}{178}{振荡函数的连续与可导}
设 $f(x)=\begin{cases}x\cos(1/x^2),&x\ne0,\\0,&x=0.\end{cases}$ 讨论 $0$ 处的连续性与可导性；若将 $x$ 改为 $x^2$，重新讨论，并用点列解释差商无极限的情形。
\end{exercise}
\begin{exercise}{179}{连续可导的分段拼接}
设 $f(x)=\begin{cases}xe^x,&x\ge0,\\x^3+x,&x<0.\end{cases}$ 判断其在 $0$ 处的连续性、可导性和导函数连续性；进一步判断 $f''(0)$ 是否存在。
\end{exercise}
\begin{exercise}{180}{先连续再比较导数}
设 $f(x)=\begin{cases}x^3+\ln x,&x\ge1,\\x^3\ln x,&0<x<1.\end{cases}$ 判断 $1$ 处的连续性与可导性；若只允许给右支加一个常数 $c$，确定能使函数连续的 $c$，并判断此时是否可导。
\end{exercise}
\begin{exercise}{181}{可导拼接反求参数}
求 $a,b$，使 $f(x)=\begin{cases}x\sin\pi x,&x>1,\\ax^2+b,&x\le1\end{cases}$ 在 $1$ 处可导；再判断此时 $f'$ 在 $1$ 处是否连续及 $f''(1)$ 是否存在。
\end{exercise}
\begin{exercise}{182}{有界因子能提供的局部信息}
设 $\varphi$ 在 $0$ 的某邻域内有界，$f(x)=\begin{cases}\ln(1+2x^2)/x^2,&x>0,\\x^2\varphi(x),&x\le0.\end{cases}$ 求 $0$ 处左右极限并判断连续性、可导性；说明给 $f(0)$ 重新赋值能否修复该处的间断。
\end{exercise}
\begin{exercise}{183}{对称差商的条件辨析}
设 $f$ 在 $0$ 处连续，且 $\lim_{h\to0}\frac{f(h)-f(-h)}{2h}=L$。举例说明 $f'(0)$ 未必存在；若再有 $\lim_{h\to0}\frac{f(h)+f(-h)-2f(0)}h=0$，证明 $f'(0)=L$。
\end{exercise}
\begin{exercise}{184}{绝对值复合的可导判据}
设 $g$ 在 $x_0$ 处可导且 $g(x_0)=0$。证明 $|g|$ 在 $x_0$ 处可导当且仅当 $g'(x_0)=0$；再求 $|x^3-3x+2|$ 的全部不可导点，并解释其中重根与单根的区别。
\end{exercise}
\subsection{函数的求导法则}
\begin{exercise}{185}{基本导数与定义域}
设 $f(x)=3x^2+2e^x-5\ln x+\sin1$。求 $f'(x)$、$f'(1)$ 及 $x=1$ 处的切线方程，说明常数项在点值和斜率中的不同作用。
\end{exercise}
\begin{exercise}{186}{乘积法则与定点导数}
设 $f(x)=x\cos x+3x^2$。求 $f'(x)$、$f'(\pi/2)$；判断表达式 $(x\cos x)'=\cos x\,(x)'$ 是否成立，并写明遗漏项。
\end{exercise}
\begin{exercise}{187}{指数三角乘积的求导}
设 $y=e^x(\cos x-2\sin x)$。求 $y'$，并求 $x\in[0,2\pi)$ 中全部水平切线所对应的横坐标。
\end{exercise}
\begin{exercise}{188}{对数因子的定义范围}
设 $y=(1+x^2)\arctan x-x\ln x$。求 $y'$；将函数在 $0$ 处连续补定义后，判断其右导数是否为有限值。
\end{exercise}
\begin{exercise}{189}{商式化简后求导}
设 $y=\frac{\sqrt x+xe^x-3x^3\sin x}{x^2}$。写出定义域并求 $y'$；比较直接使用商法则与先拆成三项时的表达式，验证结果一致。
\end{exercise}
\begin{exercise}{190}{商法则与分母零点}
设 $y=\frac{\sin x}{x-\cos x}$。求 $y'$ 并写明适用条件；证明方程 $x-\cos x=0$ 只有一个实根，判断该点能否连续补定义。
\end{exercise}
\begin{exercise}{191}{对数有理式的求导}
设 $y=\frac{2\ln x+x^3}{3\ln x+x^2}$。求 $y'$、$y'(1)$，并说明导数表达式使用时对 $x$ 及分母的限制；证明正半轴上恰有一个分母零点。
\end{exercise}
\begin{exercise}{192}{指数有理式的乘商求导}
设 $y=\frac{e^x(x^2+1)}{x+1}$。求 $y'$ 并化简，写出定义域；进一步求 $\lim_{x\to-1}(x+1)^2y'$，说明奇点处不能直接代入求导公式。
\end{exercise}
\begin{exercise}{193}{换元确定函数及导数}
已知对一切 $x\ne0$，$f(x-1/x)=\frac{x^4+1}{x^4+x^2+1}$。求 $f(t)$、$f'(t)$，说明换元后 $t$ 的范围；验证不同原像给出的函数值一致。
\end{exercise}
\begin{exercise}{194}{分段函数的导函数}
设 $f(x)=\begin{cases}x/(2x-1),&x<0,\\e^x\sin x,&x\ge0.\end{cases}$ 求所有可导点处的 $f'(x)$，单独判断 $0$ 处是否可导，并写出该处左右导数。
\end{exercise}
\begin{exercise}{195}{反函数导数的取点}
设可导函数 $f$ 的反函数为 $\varphi$，已知 $f(1)=2$、$f(2)=3$、$f'(1)=4$、$f'(2)=5$。求 $\varphi'(2)$、$\varphi'(3)$，说明为何不能将 $\varphi'(2)$ 直接写为 $1/f'(2)$。
\end{exercise}
\begin{exercise}{196}{反函数存在性与导数}
设 $f(x)=3x^3+2x+1$。证明其在 $\mathbb R$ 上有反函数 $\varphi$，求 $\varphi'(1)$；进一步求 $\varphi''(1)$，并说明所用分母为何非零。
\end{exercise}
\begin{exercise}{197}{多层复合的链式法则}
设 $y=(1-x+2x^3)^4$。求 $y'(0)$、$y'(1)$，并写出导函数；比较内函数在两点处的值及其导数各如何进入结果。
\end{exercise}
\begin{exercise}{198}{对数三角指数复合}
设 $y=\ln(\cos(e^{-x}))$。求 $y'$，写出函数的完整定义域；确定包含 $x=0$ 的最大定义区间，并说明该区间端点处表达式的变化。
\end{exercise}
\begin{exercise}{199}{根式反三角函数的求导范围}
设 $f(x)=\arcsin\sqrt{\frac{1-x}{1+x}}$。求定义域、内部各点的 $f'(x)$ 及 $f'(1/2)$；判断定义域两端点的单侧差商是否有有限极限。
\end{exercise}
\begin{exercise}{200}{共轭式化简与端点导数}
设 $y=\frac{\sqrt{x+1}-\sqrt{x-1}}{\sqrt{x+1}+\sqrt{x-1}}$。求定义域与内部导数，给出一种先化简再求导的方法；判断 $x=1$ 处右差商是否有有限极限。
\end{exercise}
\begin{exercise}{201}{多种指数结构的求导}
设 $y=3^x+x^3+x^{2x}+3^{x^2}$（$x>0$）。求 $y'$、$y'(1)$；分别指出常底指数、常指数幂、幂指函数及复合指数项所用的求导规则。
\end{exercise}
\begin{exercise}{202}{三角函数乘积与复合}
设 $y=\cos3x\,\sin^3x$。求 $y'$ 与 $y'(\pi/6)$；将 $y'$ 在 $\sin x=0$ 和 $\cos3x=0$ 的点分别代入，说明不能通过直接除去这些因子求全部驻点。
\end{exercise}
\begin{exercise}{203}{反三角函数与根式导数}
设 $y=x\sqrt{1-x^2}-\arccos x$。求内部导数并尽量化简，写出适用区间；用差商判断两个端点的单侧导数是否有限。
\end{exercise}
\begin{exercise}{204}{根式与对数的协同化简}
设 $y=x^2+x\sqrt{x^2+1}+\ln(x+\sqrt{x^2+1})$。证明对数对所有实数 $x$ 有定义，求 $y'$ 并化简；据此判断该函数在 $\mathbb R$ 上的单调性。
\end{exercise}
\begin{exercise}{205}{对数相位的复合导数}
设 $y=x[\cos(\ln x)+\sin(\ln x)]$。求 $y'$；将 $y$ 在 $0$ 处连续延拓后，判断右导数是否存在，并给出两列差商证明。
\end{exercise}
\begin{exercise}{206}{反余切指数复合求导}
设 $y=-x+\frac12\ln(1+e^{2x})-e^{-x}\operatorname{arccot}e^x$，其中 $\operatorname{arccot}\in(0,\pi)$。求 $y'$ 并化简，再判断该函数在 $\mathbb R$ 上的单调性。
\end{exercise}
\begin{exercise}{207}{先取参数极限再求导}
对每个固定实数 $x$，设 $f(x)=\lim_{t\to+\infty}x\left(\frac{t-x}{t+2x}\right)^t$，实数幂在充分大的 $t$ 上考虑。求 $f(x)$ 及 $f'(x)$；先计算定义函数的极限，不得未经论证交换求导与取极限。
\end{exercise}
\begin{exercise}{208}{抽象函数的链式法则}
设 $f$ 可导。分别求 $y=f(3^x+x^2)$ 与 $y=f(x\ln x)$ 的导数，写明第二式的定义范围；若 $f'(1)=2$，再求第一式在 $x=0$ 处的导数。
\end{exercise}
\begin{exercise}{209}{复合求导的取点对应}
已知 $f$ 可导且 $f'(1)=-2$，$g(x)=f(2x^3-3x)$。求 $g'(-1)$；再求满足 $2x^3-3x=1$ 的全部实数 $x$，指出仅用已知的 $f'(1)$ 还可确定哪些点的 $g'$。
\end{exercise}
\begin{exercise}{210}{反函数与复合函数求导}
设 $f(x)=x+e^x$，其反函数为 $\varphi$。证明反函数在 $\mathbb R$ 上存在且可导，求 $\varphi'(1)$；再设 $G(x)=\varphi(e^x)$，求 $G'(0)$ 与 $G''(0)$。
\end{exercise}
\begin{exercise}{211}{函数关系的求导约束}
设 $f:\mathbb R\to\mathbb R$ 可导，满足 $f(x+y)=f(x)+f(y)+2xy$。先求 $f(0)$，再用固定 $x$、令增量趋零的方法证明 $f'(x)=f'(0)+2x$；若 $f'(0)=3$，确定 $f(x)$ 并验证。
\end{exercise}
\subsection{高阶导数}
\begin{exercise}{212}{二阶导数与定义域}
分别求 $y=\frac{\sin x-\cos x}{\sin x+\cos x}$ 与 $y=2x^3\arctan x+\ln(1+x^2)$ 的二阶导数；写明前者的定义范围，并求后者的 $y''(0)$。
\end{exercise}
\begin{exercise}{213}{指数三角函数的二阶导数}
设 $f(x)=e^{3x}\sin2x$。求 $f''(\pi/2)$；进一步找出常数 $a,b$，使 $f''+af'+bf\equiv0$，并直接验证。
\end{exercise}
\begin{exercise}{214}{对数根式的二阶求导}
设 $y=\ln\sqrt{\frac{1+e^{2x}}{e^{2x}}}$。求 $y'(0)$、$y''(0)$，再写出 $y$ 在 $0$ 处到二阶的展开式，并说明余项形式。
\end{exercise}
\begin{exercise}{215}{反正弦函数的二阶导数}
设 $y=\arcsin(x/3)$。求 $y''(\sqrt5)$ 及 $y''$ 的定义区间；根据 $y''$ 的符号判断其在负、正半区间上的凹凸性。
\end{exercise}
\begin{exercise}{216}{对数乘积的高阶导数}
设 $y=x\ln x$（$x>0$）。求 $y^{(10)}$，再给出 $y^{(n)}$（$n\ge2$）的通式并用数学归纳法证明，明确符号与阶乘的下标。
\end{exercise}
\begin{exercise}{217}{三角函数导数的周期性}
设 $y=(\sin(x/2)+\cos(x/2))^2$。先化简再求 $y^{(10)}$；给出 $y^{(n)}$（$n\ge1$）的统一公式，并核对前两阶。
\end{exercise}
\begin{exercise}{218}{莱布尼茨公式的简化}
设 $f(x)=xe^{-x}$。求 $f^{(n)}(x)$（$n\ge1$）并证明；再求 $f^{(n)}(n)$，说明乘积中仅有两项对高阶导数有贡献的原因。
\end{exercise}
\begin{exercise}{219}{多项式次数与最高阶导数}
设 $y=(x^2+1)^5(x^5+1)^2$。不完全展开，求 $y^{(20)}$、$y^{(21)}$；再求 $y^{(19)}(0)$，说明只需检查哪些幂次。
\end{exercise}
\begin{exercise}{220}{降幂后的任意阶导数}
设 $f(x)=\sin^4x+\cos^4x$。化简并求 $f^{(n)}(x)$（$n>1$）；再求满足 $f^{(n)}(0)=0$ 的全部正整数 $n$。
\end{exercise}
\begin{exercise}{221}{复合对数的高阶导数}
设 $f(x)=\ln(2-3x)$。求 $f^{(n)}(x)$（$n\ge1$）并用归纳法证明，写出适用区间；再求 $f^{(n)}(0)$。
\end{exercise}
\begin{exercise}{222}{导数函数关系的继续求导}
设 $f$ 在 $1$ 的某邻域内可导，满足 $f'=f^2$ 且 $f(1)=2$。证明条件保证 $f''$ 存在，求 $f''(1)$；再求 $f'''(1)$，说明每次求导的合法性。
\end{exercise}
\begin{exercise}{223}{高阶零点与莱布尼茨公式}
设 $g$ 在 $a$ 附近有 $n+1$ 阶导数，$F(x)=(x-a)^ng(x)$，其中 $n\ge1$。求 $F^{(k)}(a)$（$0\le k\le n+1$），并用莱布尼茨公式证明；据此求 $(x-1)^5\ln x$ 在 $1$ 处的六阶导数。
\end{exercise}
\begin{exercise}{224}{高阶导数的递推公式}
设 $f(x)=e^{ax}\sin bx$。证明存在数列 $A_n,B_n$，使 $f^{(n)}(x)=e^{ax}(A_n\sin bx+B_n\cos bx)$，写出初值与递推式；取 $a=b=1$，求 $f^{(8)}(0)$，并用相位表示式核对。
\end{exercise}
\subsection{隐函数及由参数方程所确定的函数的导数}
\begin{exercise}{225}{隐函数求导与正则点}
分别对 $\cos(xy)+\ln y-x=1$、$x^2+y^2=6xy$ 求 $\dd y/\dd x$，写出公式分母非零的条件；对第二条曲线单独讨论原点的直线分支与切线。
\end{exercise}
\begin{exercise}{226}{隐函数分支与定点导数}
方程 $xy-\sin(\pi y^2)=0$ 在 $(0,1)$ 附近确定分支 $y=y(x)$。求 $y'(0)$；若改取过 $(0,\sqrt2)$ 的分支，重新计算并说明只给 $x=0$ 是否足以唯一确定导数。
\end{exercise}
\begin{exercise}{227}{隐曲线的切线与法线}
求曲线 $\tan(x+y+\pi/4)=y$ 在 $(-1,1)$ 处的切线与法线方程；检查该点满足原方程，并说明求导中 $y$ 对 $x$ 的依赖如何进入。
\end{exercise}
\begin{exercise}{228}{同一横坐标的多条隐分支}
求曲线 $y^2+2e^{xy}=3y$ 在横坐标 $x=0$ 对应的全部点及各点的切线方程；检查每一点处隐函数求导式的分母是否非零。
\end{exercise}
\begin{exercise}{229}{抽象隐函数的定点导数}
设 $xf(y)=e^y\ln y$ 确定过 $(0,1)$ 的可导分支，$f$ 可导且 $f(1)=1$。求 $y'(0)$；若再知 $f$ 在 $1$ 附近连续可导且 $f'(1)=c$，求 $y''(0)$，写明所需的连续性与可导性来自何处。
\end{exercise}
\begin{exercise}{230}{隐显曲线公共切线}
曲线 $y=x^2+ax+b$ 与 $2y=xy^2-3$ 在 $(1,-1)$ 处有公共切线。求全部参数 $a,b$ 及公共切线方程；分别写明通过指定点和切线斜率相同给出的条件。
\end{exercise}
\begin{exercise}{231}{隐函数的二阶导数}
在 $x\ne0$ 的区域，方程 $\ln(x^2+y^2)=2\arctan(y/x)+\pi/2+\ln2$ 确定局部可导分支 $y(x)$。求 $y'$、$y''$，写出导数公式的分母条件，并说明原方程的定义范围。
\end{exercise}
\begin{exercise}{232}{隐函数二阶导数的定点计算}
设 $e^y+xy=e+1$ 确定过 $(1,1)$ 的局部分支。求 $y'(1)$、$y''(1)$；写出该点处的一阶线性近似，并检查一阶求导式的分母非零。
\end{exercise}
\begin{exercise}{233}{含导数关系的再次求导}
设曲线 $y(x)$ 二阶可导，经过 $(1,-1)$，满足 $x^2+y^2+(y')^3=1$。求 $y'(1)$、$y''(1)$；说明为何能在该点附近由再次求导的方程解出 $y''$。
\end{exercise}
\begin{exercise}{234}{对数求导与正性}
设 $y=x^{\ln x}$（$x>0$）。用对数求导法求 $y'$，确定全部驻点及单调区间；说明取自然对数的合法性。
\end{exercise}
\begin{exercise}{235}{变底变指数的求导}
在 $1+\sin x>0$ 的区间内，设 $y=(1+\sin x)^{x^2}$。求 $y'$ 与 $y'(0)$；说明底数等于零的点为何不能直接代入对数求导公式。
\end{exercise}
\begin{exercise}{236}{幂指函数与根式乘积}
设 $y=x^{\sin x}\sqrt{e^x-1}$（$x>0$）。求 $y'$，再求 $\lim_{x\to0^+}y/\sqrt x$；据此判断将 $y(0)$ 定义为 $0$ 后右导数是否有限。
\end{exercise}
\begin{exercise}{237}{变指数的商函数}
设 $y=(1+2x)^{1/x^2}e^{-x}$，其中 $x>-1/2$、$x\ne0$。求 $y'$；分别讨论 $x\to0^+$、$x\to0^-$ 时 $y$ 的变化，判断能否在 $0$ 处连续补定义。
\end{exercise}
\begin{exercise}{238}{多因子对数求导}
设 $y=\frac{(x+1)^2\sqrt{2x+2}}{(x-1)^3\sqrt{3x-3}}$。确定实定义域，用对数求导法求 $y'$ 并化简；根据 $y'/y$ 的符号判断单调性。
\end{exercise}
\begin{exercise}{239}{幂指函数的一阶导数}
求 $y=x^x$（$x>0$）的 $y'$，确定其最小值及取值点；再求 $\lim_{x\to0^+}x^x$，判断连续延拓后的右导数是否有限。
\end{exercise}
\begin{exercise}{240}{幂指函数的二阶导数}
设 $y=x^{2x}$（$x>0$）。求 $y''$，并判断 $y''$ 在定义域内的符号；给出判断中使用的不等式或平方分解。
\end{exercise}
\begin{exercise}{241}{隐式幂指方程的分支}
在 $x,y>0$ 且 $1+\ln y\ne0$ 的分支上，方程 $x^{1/y}=y^{1/x}$ 确定 $y=y(x)$。求 $y'$、$y''$；说明取对数后方程如何化简，并核对分支 $y=x$ 上的结果。
\end{exercise}
\begin{exercise}{242}{参数方程求导与例外参数}
分别求参数曲线 $x=2t/(1+t^2),\ y=2t^2/(1+t^2)$ 和 $x=\cos^3t,\ y=\sin^3t$ 的 $\dd y/\dd x$。列出 $\dd x/\dd t=0$ 的参数，并对第一条曲线在这些参数处判断切线方向。
\end{exercise}
\begin{exercise}{243}{参数导数的定点取值}
设 $x=t^2e^{2t}$、$y=(t^2+1)e^t$。求 $\left.\dd y/\dd x\right|_{t=2}$；列出一阶参数求导公式不能直接使用的全部 $t$，并说明原因。
\end{exercise}
\begin{exercise}{244}{参数曲线切线与双分支}
求 $x=\sqrt{t^2+3}$、$y=\ln(t+\sqrt{t^2+3})$ 在 $t=1$ 对应点的切线方程；再比较 $t=-1$ 对应点的坐标及切线斜率，说明相同 $x$ 是否对应相同 $y$。
\end{exercise}
\begin{exercise}{245}{由坐标确定参数}
求曲线 $x=e^t\sin2t$、$y=e^t\cos t$ 在 $(0,1)$ 处的切线和法线方程；先确定该点对应的全部实数参数，检查是否存在多条切线。
\end{exercise}
\begin{exercise}{246}{参数曲线二阶导数}
设 $x=\cos t$、$y=t\cos t-\sin t$。求 $\dd y/\dd x$ 与 $\dd^2y/\dd x^2$，写出公式适用的参数范围；说明求二阶导数时为何还须除以 $\dd x/\dd t$。
\end{exercise}
\begin{exercise}{247}{反三角参数的二阶导数}
设 $x=\arctan t$、$y=3t+t^3$。求 $\left.\dd^2y/\dd x^2\right|_{t=1}$；消去参数写出显函数，并用显函数求导核对。
\end{exercise}
\begin{exercise}{248}{参数导数与全局单值性}
设 $x=e^t+t$、$y=\sin t$。求 $t=0$ 时的一阶、二阶 $x$ 导数；证明参数式在整个实轴上确定单值函数 $y(x)$，并写出其 $x$ 定义域。
\end{exercise}
\begin{exercise}{249}{隐函数正则点与退化点}
考虑曲线 $x^2+xy+y^2=3$。求过 $(1,1)$ 的局部分支的一阶、二阶导数及切线；找出曲线上 $\dd y/\dd x$ 公式分母为零的点，并判断这些点的切线方向。
\end{exercise}
\begin{exercise}{250}{参数退化点的导数定义}
设 $x=t^3$、$y=t^4$。消去参数并讨论原点处的可导性与切线，再讨论二阶导数是否存在；解释为什么 $t=0$ 时参数导数商为 $0/0$ 并不能直接判定不可导。
\end{exercise}
\subsection{函数的微分}
\begin{exercise}{251}{微分与增量的误差}
设 $y=x^3+x$。在 $x=3$ 处，分别取 $\Delta x=0.1,0.01$，计算 $\dd y$、$\Delta y$ 及二者之差；写出一般增量 $h$ 时的误差表达式，并说明误差相对于 $h$ 的阶。
\end{exercise}
\begin{exercise}{252}{可微定义与余项}
已知 $f$ 在每个 $x>-1$ 处的增量可写为 $\Delta y=\frac{\Delta x}{\sqrt{1+x^3}}+\alpha$，其中 $\alpha/\Delta x\to0$（$\Delta x\to0$）。求 $f'(2)$ 及 $\dd y|_{x=2}$；说明线性系数唯一的原因。
\end{exercise}
\begin{exercise}{253}{微分和增量的等价关系}
设 $f$ 在 $x_0$ 处可导且 $f'(x_0)=1$。比较 $\dd y$、$\Delta y$ 与 $\Delta x$ 在 $\Delta x\to0$ 时的等价关系并证明；若 $f'(x_0)=0$，说明还能保留哪些高阶关系。
\end{exercise}
\begin{exercise}{254}{复合函数的微分}
求 $y=e^{\sin^2(1/x)}$ 与 $y=\arcsin(x/2)+\sqrt{4-x^2}$ 的微分，并写出公式成立的范围；对第二个函数判断能否把内部微分公式直接用于端点。
\end{exercise}
\begin{exercise}{255}{定点微分的系数}
分别求 $y=\sqrt2+\sin x$ 在 $x=0$ 处和 $y=\ln(1+2^{-2x})$ 在 $x=1$ 处的微分；将 $\dd x=0.01$ 代入，写出相应的函数增量线性近似。
\end{exercise}
\begin{exercise}{256}{隐函数微分的唯一性}
设 $2y+\sin(x+y)=0$ 确定 $y=y(x)$。求 $\dd y$，证明导数公式的分母恒不为零；在过原点的分支上，求 $x=0$ 处的线性近似。
\end{exercise}
\begin{exercise}{257}{隐函数的定点微分}
设 $e^y+xy-e=0$ 在 $x=0$ 附近确定分支 $y(x)$。求该点对应的 $y$ 值及 $\dd y|_{x=0}$；据此估计 $x$ 增加 $0.01$ 时 $y$ 的增量，说明这是线性近似。
\end{exercise}
\begin{exercise}{258}{可去点处的微分}
设 $f(x)=\begin{cases}(\cos x-1)/x,&x\ne0,\\0,&x=0.\end{cases}$ 用定义求 $\dd f|_{x=0}$；说明为何不能直接把 $x=0$ 代入非零点求导所得的商式。
\end{exercise}
\begin{exercise}{259}{幂指函数的微分}
求 $y=(1+e^x)^x$ 的微分，逐项标明对底数与指数求导的贡献；再求 $\dd y|_{x=0}$，检查是否与原函数在 $0$ 处的线性近似一致。
\end{exercise}
\begin{exercise}{260}{由微分确定函数}
在各自定义区间内，求导数分别为 $x^2$、$\cos3x$、$1/(x\ln x)$、$e^x/(1+e^{2x})$ 的全部函数。写明任意常数与定义区间的关系，并直接求导核对四个结果。
\end{exercise}
\begin{exercise}{261}{夹逼条件与可微性}
设 $f$ 在 $0$ 的某邻域内有定义且 $|f(x)|\le x^2$。证明 $f$ 在 $0$ 处连续、可导并求 $f'(0)$；构造满足此条件而导函数在 $0$ 附近无界的函数，说明一点可导不能保证邻域内导数有界。
\end{exercise}
\begin{exercise}{262}{定点导数的结构化计算}
分别求：$f(x)=3x+\frac{x^2-1}{x^2+1}\arctan\frac{2x}{1+x}$ 的 $f'(0)$；$p(x)=x(x-1)\cdots(x-100)$ 的 $p'(100)$；$y=\frac{(x^2+1)\sqrt{3x+1}}{(x+1)^2e^{2x}}$ 的 $y'(1)$。每式选取合适方法，避免完全展开乘积。
\end{exercise}
\begin{exercise}{263}{等价差分给出的点值和导数}
设 $f$ 在 $x_0$ 处可导，且 $f(x_0+3h)-3f(x_0-h)\sim h$（$h\to0$）。求 $f(x_0)$、$f'(x_0)$，并用 $f(x_0+t)=f(x_0)+f'(x_0)t+o(t)$ 验证参数条件充分。
\end{exercise}
\begin{exercise}{264}{可导信息与幂指极限}
设 $f$ 在 $0$ 附近连续可导，$f(0)=1$、$f'(0)=4$。求 $\lim_{x\to0}[f(\sin x)]^{1/(2x)}$，证明底数最终为正；将指数改为 $1/(cx)$（$c\ne0$）时写出一般结果。
\end{exercise}
\begin{exercise}{265}{切线截距构造数列}
设曲线 $y=f(x)=1/(1+x^2)$ 在 $(n,f(n))$ 处的切线与 $x$ 轴相交于 $(\xi_n,0)$，其中 $n\ge1$。求 $\xi_n$，再求 $\lim_{n\to\infty}f(\xi_n)$ 及 $\lim_{n\to\infty}n^2f(\xi_n)$。
\end{exercise}
\begin{exercise}{266}{复合差商与法线}
设 $f$ 连续可导，且 $\lim_{x\to1}\frac{f(3x-1)+3}{x-1}=3$。求 $f(2)$、$f'(2)$，再求曲线在 $(2,f(2))$ 处的法线方程，说明复合函数内部系数的作用。
\end{exercise}
\begin{exercise}{267}{幂根极限与不可导点}
设 $f(x)=\lim_{n\to\infty}\frac{x^2}{(1+x^{2n})^{1/n}}$。求其分段表达式，讨论整个实轴上的连续性与可导性；对每个不可导点计算左右导数。
\end{exercise}
\begin{exercise}{268}{由连续延拓反求点值}
设 $f$ 连续可导且 $f(1)=0$、$f'(1)=1$。令 $g(x)=\begin{cases}\frac{f(e^{2x})+\ln(1-4x)}x,&x\ne0,\\b,&x=0,\end{cases}$ 在 $0$ 的充分小邻域内定义。求使 $g$ 在 $0$ 连续的 $b$，并说明现有条件是否足以保证 $g'(0)$ 存在。
\end{exercise}
\begin{exercise}{269}{极限生成分段函数的可导条件}
设 $f(x)=\lim_{n\to\infty}\frac{x^2e^{n(x-1)}+3ax+b}{e^{n(x-1)}+1}$。求分段表达式；确定使其在 $1$ 处可导的全部 $a,b$；在这些参数下求整个实轴上的 $f'$，并判断其连续性。
\end{exercise}
\begin{exercise}{270}{互补内函数的链式法则}
设 $f$ 可导，$y=f(\sin^2x)+f(\cos^2x)$。求 $y'$、$y'(0)$，再证明该函数以 $\pi/2$ 为周期；说明这是否保证最小正周期一定等于 $\pi/2$。
\end{exercise}
\begin{exercise}{271}{抽象函数的两层复合}
设 $f:\mathbb R\to\mathbb R$ 可导，$y=f(f(1+3x^2))$。求 $y'$ 与 $y'(0)$；若 $f$ 二阶可导，再求 $y''(0)$。
\end{exercise}
\begin{exercise}{272}{公共切线与复合平方}
曲线 $y=f(x)$ 与 $y=(3+\ln x)/x$ 在横坐标 $1$ 处的切线相同。设 $F(x)=[f(e^x)]^2$，求 $F'(0)$；写明公共切线同时给出的点值信息和导数信息。
\end{exercise}
\begin{exercise}{273}{二阶可导与极限反推}
设 $f$ 在 $0$ 处二阶可导，且 $\lim_{x\to0}\frac{e^{f(x)}-1}{1-\cos x}=2$。依次求 $f(0)$、$f'(0)$、$f''(0)$；说明为何不能直接假定 $f(x)$ 从二阶项开始。
\end{exercise}
\begin{exercise}{274}{复合二阶导数的正则性条件}
设 $f$ 在 $0$ 附近连续可导，$y=f(x^2)$。用导数定义求 $y''(0)$，结果用 $f'(0)$ 表示；若仅假设 $f$ 在邻域内可导，判断该结论是否必成立，并构造反例或给出证明。
\end{exercise}
\begin{exercise}{275}{导数方程与高阶通式}
设 $f$ 任意阶可导且 $f'=e^f$。求 $f^{(n)}$（$n\ge1$）关于 $f$ 的表达式，并用数学归纳法证明；若 $f(0)=0$，再求 $f^{(n)}(0)$。
\end{exercise}
\begin{exercise}{276}{高阶零点与定点导数}
设 $n$ 为正整数，$f(x)=(x+1)^ne^{-x^2}$。用莱布尼茨公式求 $f^{(n)}(-1)$；再求 $f^{(n+1)}(-1)$，说明绝大多数项消失的原因。
\end{exercise}
\begin{exercise}{277}{反函数的二阶导数}
设 $x=\varphi(y)$ 是 $y=x+\ln x$（$x>0$）的反函数。证明反函数在 $\mathbb R$ 上存在，求 $\dd^2x/\dd y^2$，并求其在 $y=1$ 处的值。
\end{exercise}
\begin{exercise}{278}{函数及导函数的逐层正则性}
设 $f(x)=\begin{cases}x^3\sin(1/x)+\sin x+1,&x\ne0,\\1,&x=0.\end{cases}$ 判断 $f$ 在 $0$ 处的连续性和可导性；求 $f'$，讨论其在 $0$ 处的连续性与可导性；求 $f''$ 的存在范围及表达式。
\end{exercise}
\begin{exercise}{279}{隐函数定点导数与无穷远极限}
方程 $\ln(x^2+y)=x^2y+\sin x$ 在 $(0,1)$ 附近确定分支 $y=f(x)$。求 $f'(0)$，再求 $\lim_{|x|\to\infty}x\left(1-f(2/x)\right)$；说明无穷远表达式只使用了 $f$ 在哪个点附近的信息。
\end{exercise}
\begin{exercise}{280}{隐函数与外层抽象函数的复合}
设 $f$ 二阶可导且 $f'(1)=1$，$xy-e^y+e^x=0$ 确定过原点的分支 $y=y(x)$。令 $z=f(e^{y(x)}-\sin x)$，求 $z''(0)$；计算中说明为何不需要知道 $f''(1)$ 的具体值。
\end{exercise}
\begin{exercise}{281}{同一横坐标对应的参数分支}
设 $x=2t^2+3t+1$、$y=e^t\sin t+2$。求 $x=0$ 对应的全部参数及各分支的 $\dd y/\dd x$；说明题目只给横坐标时，导数是否必为唯一数值。
\end{exercise}
\begin{exercise}{282}{参数函数的三阶导数}
设 $x=t-\ln\sqrt{1+e^{2t}}$、$y=\arctan e^t$。求 $\dd^3y/\dd x^3$，结果可用 $t$ 表示；证明 $\dd x/\dd t$ 恒非零，写出 $x$ 的值域。
\end{exercise}
\begin{exercise}{283}{极坐标曲线的切线与法线}
求阿基米德螺线 $\rho=\theta$ 在 $(\rho,\theta)=(\pi/2,\pi/2)$ 处的切线与法线方程，以直角坐标表示；先写出参数式并检查该点切线是否竖直。
\end{exercise}
\begin{exercise}{284}{线性主部反求复合导数}
设 $f$ 可微，$y=f(e^{2x})$ 在 $x=0$ 处，当 $\Delta x=-0.1$ 时，增量的线性主部也为 $-0.1$。求 $f'(1)$；说明这里给出线性主部是否意味着真实增量恰好等于 $-0.1$。
\end{exercise}
\begin{exercise}{285}{复合函数的定点微分}
设 $f'(x)=\arctan(1/x)$（$x\ne0$），$y=f\left(\frac{2x-1}{x+1}\right)$。求 $\dd y|_{x=2}$；列出复合求导公式需要排除的 $x$，说明分别来自哪一层函数。
\end{exercise}
\begin{exercise}{286}{参数函数的微分}
设 $f$ 可导且 $f'(0)\ne0$，$x=f(\ln t)$、$y=f(t^2-1)$（$t>0$）。求 $t=1$ 时的 $\dd y$，分别用 $\dd x$ 与 $\dd t$ 表示，并验证两种写法一致。
\end{exercise}
\begin{exercise}{287}{线性近似与误差控制}
用微分估算 $\sqrt{4.04}$，并利用函数在 $[4,4.04]$ 上的二阶导数给出误差上界；判断近似值偏大还是偏小。说明只有 $o(\Delta x)$ 信息时，能否直接得到给定步长下的数值误差上界。
\end{exercise}
\begin{exercise}{288}{隐函数微分与局部近似}
设方程 $ye^y=x$ 在 $(x,y)=(e,1)$ 附近确定函数 $y(x)$。求该点的 $\dd y$ 与 $y''$，写出 $x=e+h$ 时到二阶的近似式；说明为什么近似中出现的 $h^2$ 项不属于一阶微分。
\end{exercise}

\clearpage
\section{第3章 微分中值定理与导数的应用}
\subsection{中值定理}
\begin{exercise}{289}{罗尔定理的条件与结论}
分别判断函数 $f(x)=|x|$、$f(x)=x^2-1$、$f(x)=x\mathrm e^{-x}$、$f(x)=1/x$ 在区间 $[-1,1]$、$[-1,1]$、$[0,1]$、$[-1,1]$ 上是否满足罗尔定理的全部条件，逐项指出依据。对满足条件的函数求全部罗尔点；对不满足条件的函数，判断开区间内是否仍有导数为零的点。
\end{exercise}
\begin{exercise}{290}{罗尔点的存在性与唯一性}
验证 $f(x)=\ln\sin x$ 在 $[\pi/6,5\pi/6]$ 上满足罗尔定理，求出全部罗尔点。若将右端点改为 $b\in(\pi/6,\pi)$，求使两个端点函数值相等的全部 $b$，并说明相应罗尔点是否唯一。
\end{exercise}
\begin{exercise}{291}{辅助函数与导数方程}
设 $f\in C[0,1]$，在 $(0,1)$ 可导，$f(0)=0$、$f(1)=\pi/4$。证明存在 $\xi\in(0,1)$ 使 $(1+\xi^2)f'(\xi)=1$，写出所用辅助函数。将端点条件改为 $f(1)=c$ 时，求使同一辅助函数的两个端点值相等的 $c$，并解释该方法为何依赖此值。
\end{exercise}
\begin{exercise}{292}{多项式零点与原函数}
设实数 $a_0,\ldots,a_n$ 满足 $\sum_{k=0}^n a_k/(k+1)=0$。不用定积分，证明多项式 $P(x)=\sum_{k=0}^n a_kx^k$ 在 $(0,1)$ 内有零点。再说明若 $P$ 不恒为零，则它在 $(0,1)$ 内必同时取正值和负值。
\end{exercise}
\begin{exercise}{293}{极限条件与二阶罗尔定理}
设 $f$ 在包含 $[0,1]$ 的开区间内二阶可导，且 $\lim_{x\to0}f(x)/x=0$、$f(1)=0$。由极限条件确定 $f(0)$ 与 $f'(0)$，再证明存在 $\xi\in(0,1)$ 使 $f''(\xi)=0$。若仅保留 $f(0)=f(1)=0$，举例说明二阶导数为零的结论是否仍成立。
\end{exercise}
\begin{exercise}{294}{端点高阶零点与连续使用罗尔定理}
设 $f$ 在包含 $[0,1]$ 的开区间内三阶可导，$f(0)=f(1)=0$，令 $F(x)=x^3f(x)$。求 $F(0)$、$F'(0)$，证明在 $(0,1)$ 内至少存在一点使 $F''$ 为零。进一步证明 $F'''$ 在 $(0,1)$ 内也至少有一个零点，明确每次应用罗尔定理的区间。
\end{exercise}
\begin{exercise}{295}{端点匹配与指数辅助函数}
设 $f\in C[0,1]$，在 $(0,1)$ 可导，$f(0)=1$、$f(1)=\mathrm e^{-1}$。证明存在 $\xi\in(0,1)$ 使 $f'(\xi)+\mathrm e^{-\xi}=0$。若另知 $f''(x)>\mathrm e^{-x}$ 对所有 $x\in(0,1)$ 成立，证明该点唯一；此问假设 $f$ 在 $(0,1)$ 二阶可导。
\end{exercise}
\begin{exercise}{296}{乘积辅助函数与参数推广}
设 $a>0$，$f\in C[0,a]$，在 $(0,a)$ 可导，且 $f(a)=0$。证明存在 $\xi\in(0,a)$ 使 $f(\xi)+\xi f'(\xi)=0$。对任意给定 $p>0$，进一步证明存在 $\eta\in(0,a)$ 使 $pf(\eta)+\eta f'(\eta)=0$，并说明为何无需假设 $x^pf(x)$ 在 $0$ 处可导。
\end{exercise}
\begin{exercise}{297}{加权导数零点}
设 $f\in C[0,1]$，在 $(0,1)$ 可导，且 $f(1)=0$。证明存在 $\xi\in(0,1)$ 使 $f'(\xi)=-2f(\xi)/\xi$。讨论若把区间改为 $[a,1]$，其中 $0<a<1$，且仍只给出 $f(1)=0$，结论是否必然成立；给出证明或反例。
\end{exercise}
\begin{exercise}{298}{对称乘积与对数导数}
设 $f\in C[0,1]$，在 $(0,1)$ 可导，$f(0)=0$，且 $f(x)\ne0$ 对 $0<x<1$ 成立。证明存在 $\xi\in(0,1)$ 使
\[\frac{f'(\xi)}{f(\xi)}=\frac{f'(1-\xi)}{f(1-\xi)}.\]
分别用对称乘积构造法与取 $\xi=1/2$ 的方法论证，并说明上述条件中哪些对后一个论证并非必要。
\end{exercise}
\begin{exercise}{299}{平方辅助函数与非零分母}
设 $f\in C[0,1]$，在 $(0,1)$ 可导，$f(0)=f(1/2)=1$、$f(1)=2$。证明存在 $\xi\in(0,1)$ 使 $2f(\xi)f'(\xi)=\mathrm e^{\xi}$，从而 $f'(\xi)=\mathrm e^{\xi}/[2f(\xi)]$。证明中须说明辅助函数为何取得相同的两个函数值，以及所得点上 $f(\xi)\ne0$。
\end{exercise}
\begin{exercise}{300}{离散函数值约束与驻点}
设 $f\in C[0,3]$，在 $(0,3)$ 可导，且 $f(0)+f(1)+f(2)=3$、$f(3)=1$。证明存在 $\xi\in(0,3)$ 使 $f'(\xi)=0$。再判断能否保证这样的点位于 $(0,2)$，并用具体函数支持结论。
\end{exercise}
\begin{exercise}{301}{介值定理与指数辅助函数}
设 $f\in C[0,1]$，在 $(0,1)$ 可导，$f(0)=f(1)=0$、$f(1/2)=1$。证明存在 $\eta\in(1/2,1)$ 使 $f(\eta)=\eta$。对任意实数 $\lambda$，证明存在 $\xi\in(0,\eta)$ 使 $f'(\xi)-\lambda f(\xi)+\lambda\xi=1$，并写出与所选 $\eta$ 对应的辅助函数。
\end{exercise}
\begin{exercise}{302}{三零点与二阶导数}
设抛物线 $p(x)=-x^2+Bx+C$ 的两个零点为 $a<b$，$f$ 在包含 $[a,b]$ 的开区间内二阶可导，且 $f(a)=f(b)=0$。若存在 $c\in(a,b)$ 使 $f(c)=p(c)$，证明存在 $\xi\in(a,b)$ 使 $f''(\xi)+2=0$。说明若交点仅允许为端点，为何结论不再得到保证。
\end{exercise}
\begin{exercise}{303}{拉格朗日中值定理的适用条件}
分别判断 $\ln(1+x^2)$、$|x|$、$1/\sqrt{1-x^2}$、$\sin(1/x)$ 在 $[-1,1]$ 上是否满足拉格朗日中值定理的条件，指出每项缺失或满足的条件。对第一项求所有中值点；对第二项直接判断中值等式是否存在解。
\end{exercise}
\begin{exercise}{304}{端点连续性不可省略}
设 $a<b$，$f$ 在 $[a,b]$ 有定义且在 $(a,b)$ 可导。判断下列四个断言是否必然成立，错误断言须给出反例：若 $f(a)f(b)<0$，则内部有驻点；若 $f(a)=f(b)$，则内部有驻点；存在 $\xi\in(a,b)$ 使 $f(b)-f(a)=f'(\xi)(b-a)$；对任意 $\xi\in(a,b)$ 都有 $\lim_{x\to b^-}[f(x)-f(\xi)]=0$。说明加上 $f\in C[a,b]$ 后，哪些断言变为真。
\end{exercise}
\begin{exercise}{305}{中值定理证明三角不等式}
用拉格朗日中值定理证明 $\sin x/x>\cos x$ 对 $0<x<\pi$ 成立。再判断该不等式在 $-\pi<x<0$ 是否成立，说明奇偶性在推广中的作用。
\end{exercise}
\begin{exercise}{306}{对数增量的双侧估计}
设 $a>b>0$。用中值定理证明
\[\frac{a-b}{a}<\ln\frac ab<\frac{a-b}{b}.\]
固定 $b$，求 $a\to b^+$ 时 $\ln(a/b)$ 与这两个界的比值，说明严格不等号为何不妨碍一阶近似。
\end{exercise}
\begin{exercise}{307}{对数不等式与余项控制}
用中值定理证明 $x/(1+x)<\ln(1+x)<x$ 对 $x>0$ 成立。利用这一结论证明 $0<x-\ln(1+x)<x^2/(1+x)$，并判断所得上界能否直接用于确定 $[x-\ln(1+x)]/x^2$ 的精确极限。
\end{exercise}
\begin{exercise}{308}{指数函数的切线比较}
用中值定理证明 $\mathrm e^x>\mathrm e x$ 对 $x>1$ 成立。进一步确定全部实数 $x$ 上 $\mathrm e^x-\mathrm e x$ 的最小值及等号条件，并解释两部分论证的联系。
\end{exercise}
\begin{exercise}{309}{导数恒等式与端点延拓}
在 $(-1,1)$ 上用求导与中值定理证明 $\arcsin x+\arccos x=\pi/2$，再将恒等式延拓至 $[-1,1]$。说明在端点不能直接使用通常的反三角函数求导公式，仍可得到等式的依据。
\end{exercise}
\begin{exercise}{310}{反三角函数的分段恒等式}
证明 $2\arctan x+\arcsin[2x/(1+x^2)]=\pi$ 对 $x\ge1$ 成立。进一步求该表达式在 $0\le x\le1$ 上的简化式，说明反正弦主值范围为何导致分段。
\end{exercise}
\begin{exercise}{311}{乘积函数的中值等式}
设 $f\in C[a,b]$，在 $(a,b)$ 可导，$a<b$。证明存在 $\xi\in(a,b)$ 使
\[\frac{bf(b)-af(a)}{b-a}=f(\xi)+\xi f'(\xi).\]
取 $f(x)=x^2$、$[a,b]=[1,2]$ 求全部满足等式的点，并检查它们确实位于开区间内。
\end{exercise}
\begin{exercise}{312}{导数单调与次可加性}
设 $c>0$，$f\in C[0,c]$，在 $(0,c)$ 可导，$f'$ 单调不增，$f(0)=0$。对 $0\le a\le b\le a+b\le c$，用拉格朗日中值定理证明 $f(a+b)\le f(a)+f(b)$。若 $f'$ 严格递减，讨论 $a=0$ 与 $a>0$ 两种情况下等号能否成立。
\end{exercise}
\begin{exercise}{313}{柯西中值定理与导数非零条件}
对 $f(x)=\sin x$、$g(x)=1+\cos x$ 在 $[0,\pi/2]$ 上验证柯西中值定理，求全部满足
\[\frac{f(\pi/2)-f(0)}{g(\pi/2)-g(0)}=\frac{f'(\xi)}{g'(\xi)}\]
的 $\xi$。说明 $g'(0)=0$ 是否妨碍应用定理，以及分母为何在所求点处非零。
\end{exercise}
\begin{exercise}{314}{倒数变量与柯西中值定理}
设 $0<a<b$，$f\in C[a,b]$ 且在 $(a,b)$ 可导。用柯西中值定理证明存在 $\xi\in(a,b)$ 使
\[\frac{af(b)-bf(a)}{a-b}=f(\xi)-\xi f'(\xi).\]
明确选取的两个辅助函数，并指出 $a>0$ 在此证明中保证了什么。
\end{exercise}
\begin{exercise}{315}{幂函数参照与加权中值点}
设 $0<a<b$，$f\in C[a,b]$ 且在 $(a,b)$ 可导。证明存在 $\xi\in(a,b)$ 使
\[\frac{f(b)-f(a)}{b-a}=(a^2+ab+b^2)\frac{f'(\xi)}{3\xi^2}.\]
将区间改为 $[-1,1]$，判断上述含商等式是否仍对所有可导函数成立，并说明原证明中哪一条件失效。
\end{exercise}
\begin{exercise}{316}{对数复合函数的中值点}
用拉格朗日中值定理证明存在 $\xi\in(1,\mathrm e)$ 使 $\sin1=\cos(\ln\xi)$，写出中值定理所作用的区间与函数。利用余弦在 $(0,1)$ 上的单调性，证明这样的 $\xi$ 唯一。
\end{exercise}
\begin{exercise}{317}{带参数的中值恒等式}
设 $f\in C[0,1]$，在 $(0,1)$ 可导，$f(0)=0$、$f(1)=1$。证明对每个实数 $\lambda$ 都存在 $\xi\in(0,1)$ 使 $f'(\xi)+\lambda[f(\xi)-\xi]=1$。若要求同一个 $\xi$ 对所有 $\lambda$ 同时适用，现有条件是否足够？给出证明或反例。
\end{exercise}
\begin{exercise}{318}{线性插值误差与二阶中值定理}
设 $f\in C^2[a,b]$，$a<b$，$L(x)=[(b-x)f(a)+(x-a)f(b)]/(b-a)$。证明对每个 $x\in(a,b)$ 存在 $\xi\in(a,b)$ 使 $f(x)-L(x)=\tfrac12f''(\xi)(x-a)(x-b)$。由此在 $|f''|\le M$ 的条件下给出整个 $[a,b]$ 上的统一误差界，并说明界中常数能否改小。
\end{exercise}
\subsection{洛必达法则}
\begin{exercise}{319}{高阶分母与单侧无穷极限}
讨论 $\lim_{x\to1}(x^3-3x^2+2)/(x^3-x^2-x+1)$ 是否存在，并分别求左右极限。先判断不定式类型，再用因式分解或洛必达法则分析零点阶数，说明不能在分母求导后为零时直接把该步写成有限数值。
\end{exercise}
\begin{exercise}{320}{基本对数极限与二阶修正}
用洛必达法则求 $\lim_{x\to1}\ln x/(x-1)$，核对使用条件。再求 $\lim_{x\to1}[\ln x-(x-1)]/(x-1)^2$，说明再次求导前为何必须重新检查是否仍为不定式。
\end{exercise}
\begin{exercise}{321}{混合函数的一阶极限}
用洛必达法则求 $\lim_{x\to1}[x^3-1+\ln x]/(\mathrm e^x-\mathrm e)$。若将分子中的 $\ln x$ 改为 $\alpha\ln x$，求使极限为零的实参数 $\alpha$，再求该参数下整个商与 $x-1$ 的比值极限。
\end{exercise}
\begin{exercise}{322}{指数与多项式的增长比较}
求 $\lim_{x\to+\infty}(x^4+\mathrm e^x)/(x\mathrm e^x)$。进一步对实数 $p$ 讨论 $\lim_{x\to+\infty}(x^4+\mathrm e^x)/(x^p\mathrm e^x)$，用洛必达法则或等价增长比较给出分类，不能只凭“指数增长快”省略关键极限。
\end{exercise}
\begin{exercise}{323}{三阶零点与连续洛必达}
求 $\lim_{x\to0}(x-\sin x)/(x^2\sin x)$，每次使用洛必达法则前确认不定式。再求使 $x-\sin x-cx^3=o(x^3)$ 的 $c$，说明结果与所求极限之间的关系。
\end{exercise}
\begin{exercise}{324}{对数与任意正幂的比较}
设 $a>0$，求 $\lim_{x\to+\infty}\ln(x\ln x)/x^a$，说明此处 $x$ 应取在什么范围。进一步讨论 $a=0$、$a<0$ 时的极限，展示正参数限制对结论的作用。
\end{exercise}
\begin{exercise}{325}{根式零点与指数三角相消}
求 $\lim_{x\to0}(\mathrm e^x-\sin x-1)/(1-\sqrt{1-x^2})$，写出实定义邻域。再把分子中的 $\sin x$ 改为 $\alpha\sin x$，求使双侧有限极限存在的全部实数 $\alpha$，并求相应极限。
\end{exercise}
\begin{exercise}{326}{反正切差与对数三阶零点}
求 $\lim_{x\to0}(\arctan x-x)/\ln(1+2x^3)$，给出对数在零点附近有意义的依据。将分母改为 $\ln(1+2x^2)$ 后再求极限，并从零点阶数说明两者不同的原因。
\end{exercise}
\begin{exercise}{327}{多个无穷小因子的主阶}
求
\[\lim_{x\to0}\frac{\sin^2x-x^2\cos^2x}{\sin x(\mathrm e^{2x}-1)\ln(1+\tan^2x)}.\]
先分别确定分子与分母的最低非零阶，再选用洛必达法则或等价无穷小完成计算，说明为何分子中的两项不能直接分别替换为 $x^2$ 后相减。
\end{exercise}
\begin{exercise}{328}{参数导致的零点阶数变化}
设 $a^2+b^2\ne0$。求使
\[\lim_{x\to0}\frac{a\tan x+b(1-\cos x)}{c\ln(1-2x)+d(1-\mathrm e^{-x^2})}=2\]
成立的全部实参数 $(a,b,c,d)$，分别处理 $c\ne0$、$c=0,d\ne0$、$c=d=0$。据此判断 $b=4d$、$b=-4d$、$a=4c$、$a=-4c$ 中哪些是所有可行参数必满足的关系。
\end{exercise}
\begin{exercise}{329}{对数三角复合的二阶相消}
求 $\lim_{x\to0}(1-\cos x)[x-\ln(1+\tan x)]/\sin^4x$。先求 $[x-\ln(1+\tan x)]/x^2$ 的极限，再组合结果，检查对数真数在零点邻域的正性。
\end{exercise}
\begin{exercise}{330}{嵌套正弦的等价无穷小}
求 $\lim_{x\to0}[\sin x-\sin(\sin x)]\sin x/x^4$。明确内变量 $u=\sin x$ 的使用，并进一步求 $[\sin x-\sin(\sin x)]/x^3$ 的极限，说明这两个极限的联系。
\end{exercise}
\begin{exercise}{331}{无穷量相减的恒等化简}
求 $\lim_{x\to0}x^{-1}(1/\sin x-1/\tan x)$。先合并为一个分式，再用半角公式或洛必达法则求值；说明两个倒数项分别无界时为何不能单独取极限再作减法。
\end{exercise}
\begin{exercise}{332}{多项式与指数衰减的参数条件}
设 $n$ 为正整数、$\lambda>0$，求 $\lim_{x\to+\infty}x^n\mathrm e^{-\lambda x}$，给出能处理任意 $n$ 的论证。再讨论 $\lambda=0$ 与 $\lambda<0$ 时的极限，说明正参数假设的必要性。
\end{exercise}
\begin{exercise}{333}{零乘无穷与线性变量替换}
求 $\lim_{x\to1}(1-x)\tan(\pi x/2)$，分别说明左右两侧的乘积符号。通过令 $t=\pi(1-x)/2$ 化为基本三角极限，再用洛必达法则核验。
\end{exercise}
\begin{exercise}{334}{三角无穷量之差与可去间断}
求 $\lim_{x\to\pi/2}(\sec x-\tan x)$。先用共轭恒等式消去无穷量相减，给出左右极限，并判断原表达式在 $\pi/2$ 补值后是否可导；可导时求导数。
\end{exercise}
\begin{exercise}{335}{倒数差的补值与导数}
设 $h(x)=1/x-1/(\mathrm e^x-1)$，$x\ne0$。求 $\lim_{x\to0}h(x)$，补定义 $h(0)$ 使其连续，再判断补值后的函数在 $0$ 是否可导并求导数。每次合并分式后说明所用极限方法的条件。
\end{exercise}
\begin{exercise}{336}{二阶主部消去与偶性}
求 $\lim_{x\to0}[1/x^2-1/(x\tan x)]$。给出连续补值后函数在 $0$ 的导数，并说明偶性只有在可导性已经建立后才能用于确定导数值。
\end{exercise}
\begin{exercise}{337}{零的零次幂型极限}
求 $\lim_{x\to0^+}x^x$，通过取对数把它化为可应用洛必达法则的形式。进一步求 $\lim_{x\to0^+}(x^x-1)/(x\ln x)$，说明原底数必须取正的原因。
\end{exercise}
\begin{exercise}{338}{一的无穷次幂型与符号}
求 $\lim_{x\to1}x^{1/(1-x)}$，明确取对数后分母的符号。对实参数 $a$ 求 $\lim_{x\to1}x^{a/(1-x)}$，说明 $a=0$ 时与一般公式是否相容。
\end{exercise}
\begin{exercise}{339}{三角分母指数的幂极限}
求 $\lim_{x\to0}(1+3x)^{2/\sin x}$，给出底数为正的邻域。再求对实参数 $a,b$ 的 $\lim_{x\to0}(1+ax)^{b/\sin x}$，说明双侧取极限的合法性与指数为零的情形。
\end{exercise}
\begin{exercise}{340}{二阶对数之比的幂极限}
求 $\lim_{x\to0}(\cos x)^{1/\ln(1+x^2)}$。先确定取实数幂时底数为正的邻域，再处理对数之比。将分母改为 $\ln(1+ax^2)$、$a>0$ 后给出参数结果。
\end{exercise}
\begin{exercise}{341}{趋零底数与趋零指数}
求 $\lim_{x\to+\infty}(x^{1/x}-1)^{1/\ln x}$，先证明底数在所考虑范围为正。取对数后使用 $\mathrm e^u-1\sim u$，分开处理 $\ln x$ 与 $\ln\ln x$ 的增长，不能把指数趋于零直接解释为极限为 $1$。
\end{exercise}
\begin{exercise}{342}{双指数平均的幂极限}
设 $a,b>0$，求 $\lim_{x\to0}[(a^x+b^x)/2]^{3/x}$。进一步讨论一般实参数 $p$ 下指数改为 $p/x$ 的结果，并核验 $a=b$ 时可由原表达式直接得到同一结论。
\end{exercise}
\begin{exercise}{343}{复合指数差的无穷小阶数}
求实数 $n>0$，使 $x\to0^+$ 时 $\mathrm e^{\tan x}-\mathrm e^x$ 与 $x^n$ 为同阶无穷小，并求两者比值的非零极限。说明若把“同阶”加强为“等价”，是否仍能直接使用 $x^n$，应如何调整主项系数。
\end{exercise}
\begin{exercise}{344}{分段函数的连续补值与高阶信息}
设 $f(x)=[\sin2x+\mathrm e^{2ax}-1]/x$（$x\ne0$）、$f(0)=a$。求使 $f$ 在实数轴连续的全部实参数 $a$，再判断相应函数在 $0$ 是否可导并求导数。注意参数同时出现于非零分支和补定义值中，必须解出由连续性产生的参数方程。
\end{exercise}
\begin{exercise}{345}{幂函数的连续延拓}
在 $0<|x|<\pi/2$ 设 $f(x)=(\cos x)^{x^2}$，令 $f(0)=a$。求使其在 $0$ 连续的 $a$，再求 $\lim_{x\to0}[f(x)-a]/x^4$。比较将指数改为 $1/x^2$ 时的连续补值，说明指数中幂次对结果的影响。
\end{exercise}
\begin{exercise}{346}{端点奇异项的抵消}
设 $f(x)=1/(\pi x)+1/\sin\pi x-1/[\pi(1-x)]$，$1/2\le x<1$。求连续延拓值 $f(1)$，再求 $\lim_{x\to1^-}[f(x)-f(1)]/(x-1)$。令 $t=1-x$ 后明确哪些奇异项相消，不能分别代入取值。
\end{exercise}
\begin{exercise}{347}{迭代数列与相邻项的幂极限}
设 $0<x_1<\pi$，$x_{n+1}=\sin x_n$。证明数列从某项起严格递减且为正，求其极限；再分别求 $\lim_{n\to\infty}(x_{n+1}/x_n)^{1/x_n}$ 与 $\lim_{n\to\infty}(x_{n+1}/x_n)^{1/x_n^2}$，并说明即使 $x_1>\pi/2$ 也不影响论证。
\end{exercise}
\begin{exercise}{348}{洛必达法则失效时的原极限}
设 $f(x)=x+\sin x$、$g(x)=x$，$x>0$。求 $\lim_{x\to+\infty}f(x)/g(x)$，并证明 $f'(x)/g'(x)$ 没有极限。由此说明洛必达法则的导数比极限条件是应用定理的条件，能否把它当成原商有极限的必要条件。
\end{exercise}
\begin{exercise}{349}{参数主项消去与幂极限}
设 $a\in\mathbb R$，求 $\lim_{x\to0}(1+x+ax^2)^{1/x}$，再讨论 $\lim_{x\to0}(1+x+ax^2)^{1/x^2}\mathrm e^{-1/x}$。说明底数在零点附近为正，并用对数变换解释两问需要的不同展开阶数。
\end{exercise}
\subsection{泰勒公式}
\begin{exercise}{350}{两类泰勒余项与误差估计}
对任意正整数 $n$，写出 $\mathrm e^x$ 的 $n$ 阶麦克劳林公式，分别使用拉格朗日型余项与佩亚诺型余项。用三次麦克劳林多项式近似 $\mathrm e^{0.1}$，给出保证成立的绝对误差上界；说明仅写 $o(x^3)$ 能否直接提供这一数值误差界。
\end{exercise}
\begin{exercise}{351}{展开阶数与商的极限}
用泰勒公式求 $\lim_{x\to0}x\ln(1+x)/(\mathrm e^x-x-1)$，并求
\[\lim_{x\to0}\frac1x\left[\frac{x\ln(1+x)}{\mathrm e^x-x-1}-2\right].\]
说明求第二个极限时分子、分母分别至少需要展开到哪一阶。
\end{exercise}
\begin{exercise}{352}{三角函数的高阶消去}
用带佩亚诺余项的展开求 $\lim_{x\to0}(\sin x-x\cos x)/\tan^3x$。进一步求分子减去其三次主项后的五次主项，并说明仅将 $\sin x$ 换为 $x$ 为何会丢失所求信息。
\end{exercise}
\begin{exercise}{353}{对称展开与偶次误差}
求 $\lim_{x\to0}[\sqrt{1+x}+\sqrt{1-x}-2]/x^2$。再求实数 $c$，使 $\sqrt{1+x}+\sqrt{1-x}-2-cx^2=o(x^2)$，并计算该差与 $x^4$ 的比值极限。
\end{exercise}
\begin{exercise}{354}{有限极限确定多项式系数}
求使 $\lim_{x\to0}[\ln(1+x)-ax-bx^2]/x^2=2$ 的全部实数 $(a,b)$。对所得参数，继续求分子减去 $2x^2$ 后与 $x^3$ 的比值极限，写出保证运算有效的余项阶数。
\end{exercise}
\begin{exercise}{355}{四阶相消与复合展开}
用泰勒公式求
\[\lim_{x\to0}\frac{1+\tfrac12x^2-\sqrt{1+x^2}}{(\cos x-\mathrm e^{x^2})\sin(x^2)}.\]
再将分母中的 $\mathrm e^{x^2}$ 改为 $\mathrm e^{\alpha x^2}$，讨论极限随实参数 $\alpha$ 的变化，必要时区分左右极限。
\end{exercise}
\begin{exercise}{356}{指数与余弦的四阶比较}
讨论 $\lim_{x\to0}[\cos x-\mathrm e^{x^2/2}]/x^4$，允许结果为无穷大。进一步确定 $a,b$，使 $\cos x-\mathrm e^{ax^2+bx^4}=o(x^4)$，并解释二次项匹配后仍须比较四次项的原因。
\end{exercise}
\begin{exercise}{357}{隐含余项与函数极限恢复}
设 $f$ 在 $0$ 的某去心邻域有定义，且 $\lim_{x\to0}[\sin(6x)+xf(x)]/x^3=0$。不用假设 $f$ 可导，求 $\lim_{x\to0}[6+f(x)]/x^2$。若已知条件中的分母仅为 $x^2$，这一极限是否仍能唯一确定？给出两个具体函数说明。
\end{exercise}
\begin{exercise}{358}{高阶无穷小与参数消去}
求使 $\mathrm e^x-(ax^2+bx+1)=o(x^2)$ 的全部实数 $(a,b)$。对所得参数，求余项的无穷小阶数和非零主项系数；说明 $o(x^2)$ 与 $O(x^2)$ 两种条件对参数的约束有何不同。
\end{exercise}
\begin{exercise}{359}{根式差的阶数与连续消去}
当 $x\to0$ 时，求 $\sqrt{1-2x}-\sqrt[3]{1-3x}$ 相对于 $x$ 的无穷小阶数及主项。进一步求常数 $c$，使减去 $cx^2$ 后成为三阶无穷小，并求相应三阶主项系数。
\end{exercise}
\begin{exercise}{360}{泰勒系数与高阶导数}
设 $f(x)=\ln(1-3x)$。求任意正整数 $n$ 对应的 $f^{(n)}(0)$，并分别用归纳求导和麦克劳林系数验证。写出四阶麦克劳林多项式，说明其在 $0$ 附近的实函数定义范围。
\end{exercise}
\begin{exercise}{361}{泰勒余项的符号}
用带拉格朗日余项的泰勒公式证明 $\mathrm e^x>1+x$ 对 $x\ne0$ 成立，须同时处理 $x<0$。进一步证明对 $x\ge0$ 有 $\mathrm e^x\ge1+x+x^2/2$，并确定等号成立条件。
\end{exercise}
\begin{exercise}{362}{两端泰勒公式与三阶中值}
设 $f\in C^3[-1,1]$，$f(-1)=0$、$f(1)=1$、$f'(0)=0$。分别在 $0$ 处展开至二阶，证明存在 $\xi\in(-1,1)$ 使 $f'''(\xi)=3$。若将 $f'(0)$ 改为 $m$，推导相应三阶导数中值结论，并给出常数 $3$ 的来源。
\end{exercise}
\begin{exercise}{363}{最低点与二阶导数下界}
设 $f\in C^2[0,1]$，$f(0)=f(1)=0$，且在 $[0,1]$ 上的最小值为 $-1$。证明存在 $\xi\in(0,1)$ 使 $f''(\xi)\ge8$。构造满足条件且处处 $f''(x)=8$ 的函数，说明常数 $8$ 无法提高；证明中应明确最低点位于内部的依据。
\end{exercise}
\begin{exercise}{364}{双参数相消与极限分类}
对实数 $a,b$，讨论 $\lim_{x\to0}[\mathrm e^{ax}-1-ax-bx^2]/x^3$ 何时为有限实数，并求其值。再求使分子为 $o(x^3)$ 的全部 $(a,b)$，包括 $a=0$ 的情形。
\end{exercise}
\begin{exercise}{365}{中心差分的误差保证}
设 $f\in C^3[x_0-h,x_0+h]$，$h>0$，且 $|f'''(x)|\le M$。用泰勒公式证明中心差分 $[f(x_0+h)-f(x_0-h)]/(2h)$ 与 $f'(x_0)$ 的误差不超过 $Mh^2/6$。构造一个三次多项式使此界取等号，并讨论仅有 $f''(x_0)$ 存在是否足以保证该误差界。
\end{exercise}
\subsection{函数的单调性与曲线的凹凸性}
\begin{exercise}{366}{分数幂的定义域与单调性}
讨论 $f(x)=x-\tfrac32x^{2/3}$ 在实数范围内的单调区间，指出驻点与不可导点。若将定义域限制为 $[0,+\infty)$，说明单调区间的结论应如何表述，并比较两种定义域下 $x=0$ 附近的行为。
\end{exercise}
\begin{exercise}{367}{逐层求导证明严格不等式}
用导数证明 $\tan x>x+x^3/3$ 对 $0<x<\pi/2$ 成立。进一步证明 $\tan x-x-x^3/3=o(x^3)$，并解释严格正性与高阶小量这两个结论为何可以同时成立。
\end{exercise}
\begin{exercise}{368}{函数比较与商的单调性}
证明当 $x>0$ 时 $\ln(1+x)>\arctan x/(1+x)$。将不等式化为一个在 $0$ 处取零值的函数的正性，再证明该函数在 $(0,+\infty)$ 严格递增，给出导数符号的完整依据。
\end{exercise}
\begin{exercise}{369}{严格单调与导数符号的逻辑关系}
设 $f$ 在 $\mathbb R$ 可导且严格递增。证明 $f'(x)\ge0$，判断能否进一步断言处处 $f'(x)>0$；分别判断 $f(-x)$、$-f(-x)$ 的严格单调性，以及 $f'(-x)$ 的符号。对不成立的加强结论给出具体反例。
\end{exercise}
\begin{exercise}{370}{偶函数方程的根数}
讨论方程 $|x|^{1/4}+|x|^{1/2}-\cos x=0$ 的实根个数。先将根限制在一个有界区间，再利用偶性与正半轴上的单调性证明结论；不以数值近似代替唯一性证明。
\end{exercise}
\begin{exercise}{371}{三次函数与零点个数分岔}
对全部实参数 $a$，讨论方程 $x^3-3x-a=0$ 的不同实根个数，特别求出恰有两个不同实根的参数。给出临界参数下各根的重数，并用单调区间解释根数发生变化的原因。
\end{exercise}
\begin{exercise}{372}{参数约束与高次方程唯一根}
设实数 $a,b,c$ 满足 $3a^2-5b<0$。证明方程 $x^5-2ax^3+3bx+4c=0$ 有且仅有一个实根。若条件放宽为 $3a^2-5b\le0$，判断结论是否仍成立，并处理导数可能为零的情形。
\end{exercise}
\begin{exercise}{373}{严格凸函数的割线斜率}
设 $f\in C^2[a,b]$，$a<b$，且 $f''(x)>0$ 对 $a<x<b$ 成立。令 $F(x)=[f(x)-f(a)]/(x-a)$，$a<x\le b$。证明 $F$ 严格递增，并给出它在 $a$ 处的连续延拓值。若仅知 $f'(x)>0$，该单调性是否仍有保证？给出反例。
\end{exercise}
\begin{exercise}{374}{幂指数比较与参数边界}
设 $a>\mathrm e$，证明对每个 $x>0$ 有 $(a+x)^a<a^{a+x}$。进一步求使这一严格不等式对所有 $x>0$ 成立的全部 $a>0$，并单独处理临界值 $a=\mathrm e$。
\end{exercise}
\begin{exercise}{375}{导数下界与最优统一常数}
设 $\mathrm e<a<b<\mathrm e^2$。证明 $\ln^2b-\ln^2a>4(b-a)/\mathrm e^2$。再证明 $4/\mathrm e^2$ 是使该不等式对全部上述 $a,b$ 成立的最大常数；说明为何开区间中的严格不等号仍允许最优常数由端点极限确定。
\end{exercise}
\begin{exercise}{376}{二阶导数不存在处的凹凸变化}
设 $f(x)=x^{8/3}-x^{5/3}$，其中分数幂按实立方根定义。求定义域、所有凹凸区间和拐点坐标，分别处理二阶导数为零与二阶导数不存在的点。表述凹凸性时同时注明 $f''$ 的符号，避免名称约定造成歧义。
\end{exercise}
\begin{exercise}{377}{绝对值与三角曲线凹凸}
设 $f(x)=|1+\sin x|$。判断它在 $(\pi,3\pi/2)$ 内的严格单调性与凹凸性，写出导数符号依据。再讨论 $x=3\pi/2$ 处的可导性与二阶可导性，说明绝对值能否在全实数范围内直接去掉。
\end{exercise}
\begin{exercise}{378}{严格凸性与中点不等式}
证明对实数 $x\ne y$ 有 $(\mathrm e^x+\mathrm e^y)/2>\mathrm e^{(x+y)/2}$。进一步设 $0<t<1$，用导数或中值定理证明 $t\mathrm e^x+(1-t)\mathrm e^y>\mathrm e^{tx+(1-t)y}$，并完整说明等号出现的条件。
\end{exercise}
\begin{exercise}{379}{分段曲线的凹凸与不可导点}
设 $f(x)=|x^2-3x+2|$。写出分段表达式，求二阶导数存在的全部凹凸区间，并检查 $x=1,2$ 的连续性和可导性。本题将拐点定义为曲线上连续且两侧凹凸性发生改变的点，按此定义求全部拐点，并说明它们是否均满足 $f''(x)=0$。
\end{exercise}
\begin{exercise}{380}{拐点的参数充要条件}
求使 $(0,1)$ 成为曲线 $y=ax^3+bx^2+c$ 拐点的全部实数 $(a,b,c)$。证明条件的充分性与必要性，并解释为何仅有 $f''(0)=0$ 不能排除退化的直线或常函数。
\end{exercise}
\begin{exercise}{381}{偶性与各阶导数的符号}
设 $f\in C^2(\mathbb R)$ 为偶函数，在 $(-\infty,0)$ 有 $f'>0$、$f''<0$。推导 $f'$、$f''$ 的奇偶性，确定它们在 $(0,+\infty)$ 的符号。判断 $0$ 是否为极值点或拐点，并说明依据。
\end{exercise}
\begin{exercise}{382}{三阶判别与水平拐点}
设 $f$ 在 $x_0$ 的邻域三阶可导，$f'(x_0)=f''(x_0)=0$、$f'''(x_0)>0$。证明 $(x_0,f(x_0))$ 为拐点，判断 $x_0$ 是否为极值点。分别通过 $f''$ 的符号变化与泰勒展开说明结论，并指出此处二阶极值判别法不能直接作出判断。
\end{exercise}
\begin{exercise}{383}{微分关系与局部几何形态}
设 $f\in C^2(-\delta,\delta)$，$\delta>0$，满足 $f''(x)+[f'(x)]^2=x$，且 $f'(0)=0$。由该关系证明 $f'''(0)$ 存在并求其值，判断 $0$ 是否为极值点及 $(0,f(0))$ 是否为拐点。论证时须给出足够小邻域内的导数或函数增量符号。
\end{exercise}
\begin{exercise}{384}{函数图形信息与导函数图形}
设 $f(x)=x^3-3x$。先由单调区间和驻点说明 $f'$ 应有的符号与零点，再由凹凸性判断 $f'$ 的单调性，据此手绘 $f'$ 的图形。最后计算 $f'$ 核验，说明 $f$ 的极值点与 $f'$ 的零点之间是否总为一一对应。
\end{exercise}
\begin{exercise}{385}{参数函数的单调区间与严格单调}
设 $f_a(x)=x^3-3ax$，$a\in\mathbb R$。求全部单调区间，确定使 $f_a$ 在整个实数轴上严格递增的参数。对临界参数解释严格递增为何允许导数在孤立点为零，并求该函数的拐点。
\end{exercise}
\begin{exercise}{386}{凸性与零点定位}
设 $f\in C^2[0,1]$，$f(0)=f(1)=0$，且 $f''(x)>0$ 对 $0<x<1$ 成立。证明内部恒有 $f(x)<0$，并证明 $f'$ 在 $(0,1)$ 内恰有一个零点。若将条件改为 $f''\ge0$，判断两个严格结论分别是否仍成立，给出反例或证明。
\end{exercise}
\subsection{函数的极值和最值}
\begin{exercise}{387}{极大值的必要条件与二阶退化}
设 $f$ 在 $x_0$ 的邻域二阶可导，并在 $x_0$ 取得局部极大值。证明 $f'(x_0)=0$、$f''(x_0)\le0$。分别举出二阶导数严格小于零和等于零的例子，说明这两个必要条件能否反过来保证极大值。
\end{exercise}
\begin{exercise}{388}{驻点与不可导点的极值判别}
设 $f(x)=2x-\tfrac32x^{2/3}+1$，分数幂按实立方根定义。求所有驻点、不可导点和极值，用一阶导数符号进行判别。进一步判断该函数在 $\mathbb R$ 上是否有最大值、最小值，说明极值与最值的区别。
\end{exercise}
\begin{exercise}{389}{指数乘积的全局最小值}
求 $f(x)=x2^x$ 的极小值点与极小值，证明该极小值也是全实数范围的最小值。再判断是否存在最大值，并用两端极限支持结论。
\end{exercise}
\begin{exercise}{390}{导函数符号与极值数量}
设 $f$ 在 $\mathbb R$ 可导，且 $f'(x)=x(x+2)(x-1)^2(x-3)$。不用求出 $f$，列出导函数符号表，确定 $f$ 的单调区间以及极大值点、极小值点的数量与位置。说明 $x=1$ 处导数为零却不产生极值的原因，判断仅知导函数能否给出各极值的数值。
\end{exercise}
\begin{exercise}{391}{微分方程约束下的驻点分类}
设 $f$ 在实数轴二阶可导，满足 $xf''(x)+3x[f'(x)]^2=1-\mathrm e^{-x}$。若 $f'(x_0)=0$ 且 $x_0\ne0$，证明 $x_0$ 必为严格极小值点，并排除拐点的可能性。进一步利用关系式在 $0$ 附近的极限讨论 $f''(0)$ 与 $f'(0)$ 满足的等式。
\end{exercise}
\begin{exercise}{392}{导数比值极限与极值}
设 $f$ 在 $a$ 的邻域可导，$f'$ 在 $a$ 连续，且 $\lim_{x\to a}f'(x)/(x-a)=-1$。求 $f'(a)$、$f''(a)$，判断 $a$ 处的极值类型。若已知极限改为 $0$，说明极值类型能否确定，并给出例子。
\end{exercise}
\begin{exercise}{393}{二阶增量极限的信息边界}
设 $f$ 在 $a$ 的邻域有定义，且 $\lim_{x\to a}[f(x)-f(a)]/(x-a)^2=-1$。证明 $f$ 在 $a$ 可导，求 $f'(a)$ 并判断极值类型。判断该极限能否保证 $f''(a)$ 存在；若不能，构造在 $a$ 的邻域可导的反例。
\end{exercise}
\begin{exercise}{394}{分段指数函数与不可导极值点}
设 $f(x)=x^{2x}$（$x>0$）、$f(x)=x\mathrm e^x+1$（$x\le0$）。求各分支导数并检查 $0$ 处的连续性和可导性，确定所有局部极值。再比较各极小值与两端极限，判断全实数范围的最大值和最小值是否存在。
\end{exercise}
\begin{exercise}{395}{隐函数驻点与局部分支}
曲线由 $2y^3-2y^2+2xy-x^2=1$ 给出。在满足隐函数定理条件的点附近把 $y$ 看作 $x$ 的函数，求全部驻点及对应坐标，并判断极值类型。每个候选点均须检验关于 $y$ 的偏导数非零，不能只解 $y'=0$ 的分子。
\end{exercise}
\begin{exercise}{396}{无极值的参数范围}
求使 $f(x)=x^3+3ax^2-ax-1$ 在实数轴上无局部极值的全部实参数 $a$。处理导数有两个实根、重根、无实根三种情况，说明临界参数下驻点为何不构成极值。
\end{exercise}
\begin{exercise}{397}{闭区间最值与根式端点}
求 $f(x)=x+\sqrt{1-x}$ 在 $[-5,1]$ 上的最大值和最小值及全部取到位置。列全内部驻点与端点，并说明 $x=1$ 处导数不存在为何不影响闭区间最值的存在。
\end{exercise}
\begin{exercise}{398}{球外切圆锥的最小体积}
半径 $r>0$ 的球内切于正圆锥，求圆锥高 $h$ 使体积最小，并求最小体积。先由轴截面建立底面半径与 $h,r$ 的关系及 $h$ 的可行范围，再证明所得驻点给出全局最小值，同时检验可行区间两端的体积极限。
\end{exercise}
\begin{exercise}{399}{利润最值与整数产量}
某企业月产量为 $Q$，总成本 $C(Q)=6Q+200$，总收益 $R(Q)=12Q-0.01Q^2$，金额单位均为万元。先把 $Q\ge0$ 视为连续变量，求利润最大时的产量和最大利润，再在 $Q$ 必须为非负整数时核对结论。求企业盈利的产量范围，并区分利润与收益。
\end{exercise}
\begin{exercise}{400}{收入成本与边际量的最优条件}
某企业的总收入为 $R(x)=26x-2x^2-4x^3$，总成本为 $C(x)=8x+x^2$，产量 $x\ge0$ 视为连续变量。求利润函数、边际收入和边际成本，确定利润最大时的产量与最大利润。核验最优点处边际收入等于边际成本，并说明还需哪些符号或边界检查才能证明为全局最大值。
\end{exercise}
\begin{exercise}{401}{共轭指数与等号条件}
设 $p,q>1$ 且 $1/p+1/q=1$。利用函数最值证明对 $x>0$ 有 $x^p/p+1/q\ge x$，确定等号条件。再推出对 $u,v>0$ 有 $uv\le u^p/p+v^q/q$，说明代换过程中使用共轭指数关系的位置。
\end{exercise}
\begin{exercise}{402}{极限给定的切线与严格凸性}
设 $f\in C^2(\mathbb R)$，$\lim_{x\to0}f(x)/x=1$，且 $f''(x)>0$ 对所有实数 $x$ 成立。由极限确定 $f(0)$、$f'(0)$，利用最值或中值定理证明 $f(x)\ge x$，并确定等号条件。若仅有 $f''\ge0$，说明等号条件会怎样变化。
\end{exercise}
\begin{exercise}{403}{约束区间与边界最优}
设 $a\in\mathbb R$，$f_a(x)=x^2-2ax$，生产控制量满足 $0\le x\le1$。求最小值与全部最小值点随 $a$ 的分段表达式，再求最大值和全部最大值点。解释为何只求驻点会遗漏部分参数下的最优解。
\end{exercise}
\begin{exercise}{404}{箱体设计与可行域端点}
从边长为 $L>0$ 的正方形薄板四角各剪去边长为 $x$ 的小正方形，折成无盖长方体。写出体积与严格可行域，求最大体积及对应 $x$，证明全局最优性。若要求底面边长至少为 $L/2$，判断最优方案是否改变并说明依据。
\end{exercise}
\subsection{函数图形的描绘}
\begin{exercise}{405}{水平渐近线与交点}
求 $y=(x+1)/x$ 在 $x\to+\infty$ 与 $x\to-\infty$ 时的水平渐近线，并求铅直渐近线。判断曲线是否与水平渐近线相交，分别说明曲线在两支上从哪一侧逼近该直线。
\end{exercise}
\begin{exercise}{406}{单侧无穷极限与铅直渐近线}
求 $y=1/(1-x^2)$ 的全部铅直渐近线，分别写出每条直线左右两侧的无穷极限。再求水平渐近线，并说明有界的间断点为何不能仅因函数未定义而判为铅直渐近线。
\end{exercise}
\begin{exercise}{407}{约分、可去间断与渐近线}
求 $f(x)=(x^2-x-2)/(x^2-1)$ 的定义域、水平渐近线与铅直渐近线。分别判断 $x=1,-1$ 处的间断类型，并说明代数约分为何不能恢复原函数定义域中缺失的点。
\end{exercise}
\begin{exercise}{408}{斜渐近线与曲线位置}
求 $y=x^2/(x-1)$ 的斜渐近线和铅直渐近线，分别检验正负无穷方向的极限。利用函数与斜渐近线的差，判断两支曲线位于该直线的哪一侧，以及是否存在交点。
\end{exercise}
\begin{exercise}{409}{指数复合函数的单侧渐近线}
求 $y=(2x-1)\mathrm e^{1/x}$ 的全部水平、铅直和斜渐近线。重点分别计算 $x\to0^+$、$x\to0^-$ 的极限与 $x\to\pm\infty$ 时的斜截距，说明同一铅直渐近线的两侧是否都必须发散。
\end{exercise}
\begin{exercise}{410}{振荡因子与渐近线判定}
对定义在 $x>0$ 的函数 $f(x)=x\sin(1/x)$，求全部水平与铅直渐近线，并讨论 $x\to0^+$ 时能否连续补值。说明 $\sin(1/x)$ 的振荡为何不妨碍 $f$ 在右端点有有限极限。
\end{exercise}
\begin{exercise}{411}{两端斜渐近线与单侧发散}
求 $f(x)=x\mathrm e^{1/x}$ 的全部水平、铅直和斜渐近线。写出 $0$ 两侧的极限，并证明在正负无穷处得到的斜渐近线相同。
\end{exercise}
\begin{exercise}{412}{对数指数函数的渐近线条数}
求 $f(x)=\ln(1+\mathrm e^x)/x$ 的全部渐近线，分别处理 $x\to0^\pm$、$x\to\pm\infty$。对正无穷方向先抽出对数中的主导指数，再说明是否还可能存在斜率非零的斜渐近线。
\end{exercise}
\begin{exercise}{413}{反正切复合与有限间断辨析}
求 $f(x)=\mathrm e^{1/x}\arctan[(x^2-x+2)/((x-1)(x+2))]$ 的定义域及全部渐近线，给出渐近线总条数。分别检查 $x=0,1,-2$ 两侧的极限，区分无穷间断和有限跳跃，并计算正负无穷方向的极限。
\end{exercise}
\begin{exercise}{414}{有理函数的完整图形分析}
对 $f(x)=(x-1)^3/(x+1)^2$，求定义域、截距、全部渐近线、单调区间、极值、凹凸区间和拐点，注明关键点及两端趋势，据此手绘图形。用精确计算说明曲线是否能够穿过其斜渐近线；若能，求交点。
\end{exercise}
\begin{exercise}{415}{周期函数的图形描绘}
对 $f(x)=\cos x/\cos2x$，确定定义域、奇偶性与最小正周期。在 $[-\pi,\pi]$ 内求全部铅直渐近线、与坐标轴的交点、单调区间、极值、凹凸区间和拐点，据此手绘该区间图形并说明如何延拓至全定义域。
\end{exercise}
\begin{exercise}{416}{渐近线可被穿越}
设 $f(x)=x+\sin x/x$，$x\ne0$。求它的全部渐近线，并判断 $0$ 处的间断类型。求与斜渐近线的全部交点，证明在任意充分大的正半轴尾部仍存在交点，说明渐近线的定义实际要求的极限是什么。
\end{exercise}
\begin{exercise}{417}{参数抵消与渐近线变化}
设 $f_a(x)=(x^2+ax+1)/(x-1)$，定义域为 $\mathbb R\setminus\{1\}$。求全部渐近线随 $a$ 的变化，并确定何时 $x=1$ 为可去间断点。对非退化参数，判断曲线能否与其斜渐近线相交，证明结论。
\end{exercise}
\subsection{曲率}
\begin{exercise}{418}{显式与参数表示的曲率一致性}
求曲线 $xy=1$ 在 $(1,1)$ 处的曲率与曲率半径，分别用显式表示 $y=1/x$ 和参数表示 $(x,y)=(\mathrm e^t,\mathrm e^{-t})$ 验证。说明曲率为何与参数行进的快慢无关。
\end{exercise}
\begin{exercise}{419}{星形线的正则参数与曲率半径}
设 $a>0$，曲线为 $x=a\cos^3t$、$y=a\sin^3t$。求满足 $\sin t_0\cos t_0\ne0$ 的点处的曲率和曲率半径，确定第一象限弧上曲率的最小值。讨论 $t_0=k\pi/2$ 时通常的正则参数曲率公式为何不能直接使用，并求沿相邻弧趋近时的曲率极限。
\end{exercise}
\begin{exercise}{420}{法向方向与曲率圆}
求抛物线 $y=2x^2-3x+3$ 在 $M(1,2)$ 处的曲率、曲率中心与曲率圆方程。证明圆心位于该点的法线上，并核验该圆经过 $M$；说明仅知道曲率半径为何还不能确定圆心的位置。
\end{exercise}
\begin{exercise}{421}{弧长变量的链式求导}
设 $\rho(x)$ 是抛物线 $y=\sqrt x$ 在 $M(x,\sqrt x)$、$x\ge1$ 处的曲率半径，$s(x)$ 是从 $A(1,1)$ 到 $M$ 沿曲线的弧长。求 $\rho(x)$、$\dd s/\dd x$，再计算 $3\rho\,\dd^2\rho/\dd s^2-(\dd\rho/\dd s)^2$。明确二阶弧长导数的链式法则，不能把它直接替换为关于 $x$ 的二阶导数。
\end{exercise}
\begin{exercise}{422}{抛物线最大曲率与长度缩放}
设 $a>0$，曲线为 $y=ax^2$。求曲率函数、最大曲率及对应点，写出该点的曲率圆。若把曲线上所有坐标同时放大 $\lambda>0$ 倍，证明对应点的曲率变为原来的 $1/\lambda$，并解释量纲关系。
\end{exercise}
\begin{exercise}{423}{椭圆曲率的极值}
设 $a>b>0$，椭圆由 $x=a\cos t$、$y=b\sin t$ 表示。求曲率函数及其最大、最小值，给出对应点和曲率半径。再讨论 $a=b$ 时结论的变化，并将结果与圆的几何曲率核对。
\end{exercise}

\clearpage
\section{第4章 不定积分}
\subsection{不定积分的概念与性质}
\begin{exercise}{424}{原函数关系与二阶求导}
设 $f$ 在 $(1,+\infty)$ 上的一个原函数为 $F(x)=\ln(x+\sqrt{x^2-1})$。求 $f(x)$ 和 $f'(x)$，说明两次求导所用的定义域条件。若把原函数改为 $F(x)+C$，说明所求两个函数是否改变。
\end{exercise}
\begin{exercise}{425}{原函数族与参数退化}
已知 $\ln x$ 是 $f$ 在 $(0,+\infty)$ 上的一个原函数。对 $a>0$，判断 $\ln(ax)$、$a^{-1}\ln x$、$\ln(a+x)$、$\tfrac12(\ln x)^2$ 中哪些仍为 $f$ 的原函数，必要时给出对 $a$ 的限制。用求导与两原函数之差两种方法核对结论。
\end{exercise}
\begin{exercise}{426}{不同表达式的原函数与端点}
证明 $\arcsin(2x-1)$ 与 $2\arcsin\sqrt x$ 在 $(0,1)$ 上为同一函数的原函数，并求两者之差。讨论两个表达式能否连续延拓至 $[0,1]$，以及能否在端点仍按通常意义作为同一函数的可导原函数。
\end{exercise}
\begin{exercise}{427}{分段原函数的拼接条件}
设 $f(x)=3x^2+1$（$x>0$），$f(x)=2x$（$x\le0$）。判断
\[F(x)=\begin{cases}x^3+x+1,&x>0,\\x^2,&x\le0\end{cases}\]
是否为 $f$ 在 $\mathbb R$ 上的原函数，分别检查连续性与导数。将两个分支各自加上任意常数后，是否可能得到全实数范围的原函数？证明结论。
\end{exercise}
\begin{exercise}{428}{原函数存在性与连续性的区别}
设 $F'=f$ 在开区间 $I$ 上成立。判断 $f+F$ 是否必然可导、是否必然连续、是否必然存在原函数、是否必为初等函数，逐项说明依据。证明原函数存在性时，可使用连续函数必有原函数这一结论；不可把 $f$ 的连续性当作已知条件。
\end{exercise}
\begin{exercise}{429}{跳跃间断与原函数拼接}
设 $f(x)=x^2$（$x\ge0$）、$f(x)=\cos x$（$x<0$），另设 $g(x)=x\sin(1/x)$（$x\ne0$）、$g(0)=0$。分别判断两函数在 $\mathbb R$ 上是否存在原函数。对 $f$ 通过两侧原函数的拼接与 $0$ 处导数给出证明；对 $g$ 验证原函数存在定理所需的条件。
\end{exercise}
\begin{exercise}{430}{导函数与原函数的关系辨析}
已知 $\sin x$ 是 $f(x)$ 的一个原函数，求 $\int f(x)\,\dd x$ 与 $f(x)$。再将条件改为 $f'(x)=\sin x$，写出此时 $f$ 及其全部原函数，说明两种条件提供的信息有何区别。
\end{exercise}
\begin{exercise}{431}{切线斜率与初值确定曲线}
一条可导曲线经过 $(\mathrm e,-1)$，且在每个 $x>0$ 处的切线斜率为 $\ln x/x$。求该曲线方程，证明初值条件使积分常数唯一。再求曲线的单调区间与最低点，以检查所得方程与斜率信息相容。
\end{exercise}
\begin{exercise}{432}{不定积分的微分与量的类型}
设 $f$ 在 $\mathbb R$ 上可微。分别求 $\int f'(x)\,\dd x$、$\dd(\int f'(x)\,\dd x)$、$(\int f'(x)\,\dd x)'$，说明三者中哪些是函数族、哪些是微分或确定函数。用 $f(x)=x^2$ 检验符号及积分常数的处理。
\end{exercise}
\begin{exercise}{433}{积分、求导与微分的逆运算}
设 $f\in C^1(I)$，$I$ 为开区间。判断下列等式在通常不定积分记号下是否成立，错误者改正：$\dd(\int f\,\dd x)=f$；$(\int f\,\dd x)'=\int f'\,\dd x$；$\int f'\,\dd x=f$；$\dd(\int f\,\dd x)=f\,\dd x$。说明等式两侧的任意常数为何不能随意遗漏。
\end{exercise}
\begin{exercise}{434}{导数相等与原函数族相等}
设 $f,g$ 在开区间 $I$ 可导，且 $f'=g'$。证明 $f-g$ 为常数，判断是否必有 $f=g$、$f=g+1$、$(\int f\,\dd x)'=(\int g\,\dd x)'$，并判断 $\int\dd f=\int\dd g$ 作为函数族是否成立。若再给出一点 $f(x_0)=g(x_0)$，哪些结论随之加强？
\end{exercise}
\begin{exercise}{435}{嵌套积分微分与常数}
设 $f$ 在开区间可导，求 $\int\dd[\int\dd f(x)]$。逐层说明每次积分引入的常数和微分消去的常数，再判断交换两次任意常数的记号是否会改变最终函数族。
\end{exercise}
\begin{exercise}{436}{复合自变量恢复与定义范围}
设 $f$ 在 $(0,1)$ 可导，且对 $0<t<\pi/2$ 有 $f'(\sin^2t)=\cos2t+\tan^2t$。求 $f(x)$，再由 $f(1/2)=0$ 确定常数。解题时说明代换 $x=\sin^2t$ 在所选区间的值域，不能将 $t$ 与 $x$ 混作同一变量。
\end{exercise}
\begin{exercise}{437}{基本积分公式与定义域分支}
求
\[\int\frac{(x^2-1)\sqrt{1-x^2}-3x}{x\sqrt{1-x^2}}\,\dd x.\]
先化为基本函数之和，再说明最大实定义区间，以及跨过 $x=0$ 时两侧积分常数是否必须相同。
\end{exercise}
\begin{exercise}{438}{三角恒等变形与原函数核验}
求 $\int\cos2x/(\sin^2x\cos^2x)\,\dd x$，先把被积函数拆成两个基本三角函数的平方。说明答案在哪些开区间成立，并对结果求导检验拆分时的符号。
\end{exercise}
\begin{exercise}{439}{半角恒等式与奇点}
求 $\int(1-2\cos x)/(1-\cos2x)\,\dd x$。先写出被积函数的实定义域，再利用 $1-\cos2x=2\sin^2x$ 积分。检查哪些被排除的点为真正奇点，积分常数应如何按区间理解。
\end{exercise}
\begin{exercise}{440}{指数基数与代数约分}
求 $\int[1/(5^x\mathrm e^x)+(1-\mathrm e^{2x})/(1+\mathrm e^x)]\,\dd x$。先说明指数乘积与分式的化简，再确定满足 $F(0)=0$ 的原函数。若将第一项的 $5$ 改为实参数 $a>0$，说明何时需要单独处理常数被积项。
\end{exercise}
\begin{exercise}{441}{基本积分与降幂公式}
求 $\int[6\sqrt x+5\mathrm e^x+\cos^2(x/2)]\,\dd x$，说明余弦平方项的处理。进一步求满足 $F(0)=0$ 的连续原函数，并核验右端点处的右导数。
\end{exercise}
\begin{exercise}{442}{多项式除法与反三角积分}
求 $\int[x^4/(1+x^2)-2/\sqrt{1-x^2}]\,\dd x$。先对有理项作多项式除法，明确共同实定义区间；再确定满足 $F(0)=1$ 的原函数。
\end{exercise}
\begin{exercise}{443}{不连通定义域上的积分常数}
在 $\mathbb R\setminus\{0\}$ 上求 $1/x$ 的全部原函数。证明将答案统一写成 $\ln|x|+C$ 且只允许一个常数会遗漏原函数，并判断是否存在能连续延拓至 $0$ 的原函数。
\end{exercise}
\begin{exercise}{444}{拼接处可导性的必要条件}
设 $f_a(x)=x$（$x<0$）、$f_a(0)=a$、$f_a(x)=2x$（$x>0$）。求使 $f_a$ 在 $\mathbb R$ 上有原函数的全部实数 $a$，并求满足 $F(0)=1$ 的原函数。分别证明条件的必要性和充分性。
\end{exercise}
\subsection{换元积分法}
\begin{exercise}{445}{线性与二次内函数的凑微分}
分别求
\[\int\frac{\dd x}{\sqrt{2x+3}},\qquad \int\frac{x+1}{x^2+2x+5}\,\dd x,\qquad \int\frac{x}{(2x+3)^2}\,\dd x.\]
每式写出代换及实定义区间；比较第三式与第一式为何不能使用相同的幂函数积分步骤直接得到答案。
\end{exercise}
\begin{exercise}{446}{根式内函数与代换值域}
分别求
\[\int\frac{\mathrm e^{2\sqrt x}}{\sqrt x}\,\dd x,\qquad \int\frac{\sin\sqrt x}{\sqrt x}\,\dd x,\qquad \int\frac{\dd x}{(1+x)\sqrt x}.\]
使用同一个根式代换，写出新变量范围。讨论各被积函数在 $x\to0^+$ 时的有限性，以及所得原函数能否连续延拓至 $0$。
\end{exercise}
\begin{exercise}{447}{幂指数、参数与抽象复合换元}
设 $n\in\mathbb R$，$f\in C^1(\mathbb R)$。求
\[\int x^2\mathrm e^{x^3}\,\dd x,\quad \int\frac{x}{1+x^4}\,\dd x,\quad \int\frac{x^2}{(2x^3-1)^n}\,\dd x,\quad \int xf(x^2)f'(x^2)\,\dd x.\]
第三式先限制 $2x^3-1>0$，分别处理 $n=1$ 与 $n\ne1$；第四式说明原函数中 $f$ 应取何种幂次，并以求导验证。
\end{exercise}
\begin{exercise}{448}{三角积分的幂次与降幂选择}
求下列四个不定积分，并分别注明采用的恒等变形或代换：
\[\int\sin^2x\cos^2x\,\dd x,\quad \int\sin^5x\cos^2x\,\dd x,\quad \int\cot^3x\csc x\,\dd x,\quad \int\sin5x\sin3x\,\dd x.\]
对第三式说明为何需要分开理解跨越 $k\pi$ 的积分常数；其余各式给出可在全实数范围成立的原函数。
\end{exercise}
\begin{exercise}{449}{平方和与平方差的参数积分}
对实数 $a\ne0$，求 $\int\dd x/(a^2+x^2)$ 与 $\int\dd x/(a^2-x^2)$，答案须对 $a<0$ 仍成立。说明两式的实定义区间和绝对值的必要性；再单独处理 $a=0$，判断能否直接把它代入非零参数公式。
\end{exercise}
\begin{exercise}{450}{复合指数函数的链式识别}
求 $\int\mathrm e^{\tan x}\sec^2x\,\dd x$ 与 $\int(1-x^{-2})\mathrm e^{x+1/x}\,\dd x$。逐式写出内函数及其导数，注明答案成立的区间，并说明第二式中不能跨越 $x=0$ 统一使用一个积分常数的原因。
\end{exercise}
\begin{exercise}{451}{多层对数代换与实定义域}
分别求
\[\int\frac{\dd x}{x\sqrt{1-\ln^2x}},\quad \int\frac{\dd x}{x\ln x[1+\ln(\ln x)]^2},\quad \int\frac{1+\ln x}{1+x^2\ln^2x}\,\dd x.\]
写出每层代换，并逐式给出实定义域，特别检查第二式中的对数与分母条件。
\end{exercise}
\begin{exercise}{452}{反三角内函数与链式系数}
分别求
\[\int\frac{\dd x}{\sqrt{1-x^2}(\arcsin x)^2},\quad \int\frac{\arcsin\sqrt x}{\sqrt{x(1-x)}}\,\dd x,\quad \int\frac{\arcsin(1-x)}{\sqrt{2x-x^2}}\,\dd x.\]
对每式注明实定义区间；通过求导核验第二式中的系数与第三式中的负号。
\end{exercise}
\begin{exercise}{453}{配方与两个原函数的常数差}
求 $\int\dd x/\sqrt{x-x^2}$，分别使用配方三角代换和 $x=\sin^2t$、$0<t<\pi/2$。说明两种所得表达式在哪一开区间为同一被积函数的原函数，并求两表达式之间的常数差。
\end{exercise}
\begin{exercise}{454}{有界内函数与反正切原函数}
求 $\int\cos x/(9+\sin^2x)\,\dd x$，并确定满足 $F(0)=0$ 的原函数。说明代换 $u=\sin x$ 在全实数轴不单调，为何仍可通过求导确认所得表达式在全实数轴成立。
\end{exercise}
\begin{exercise}{455}{根式分母的拆分代换}
求 $\int(1-x)/\sqrt{9-4x^2}\,\dd x$，分别处理常数分子项与一次分子项。明确实定义区间，求满足 $F(0)=0$ 的原函数，并核验两个部分的系数。
\end{exercise}
\begin{exercise}{456}{三角根式内函数与连通区间}
求 $\int(\sin x+\cos x)/\sqrt[3]{\sin x-\cos x}\,\dd x$，立方根按实数定义。给出被积函数全部实定义区间，并证明原函数在各区间可连续延拓至端点。比较分母改为平方根时实定义域的变化，说明不能给立方根附加非负真数条件。
\end{exercise}
\begin{exercise}{457}{根式三角复合与负号}
求 $\int x\tan(\sqrt{1-x^2})/\sqrt{1-x^2}\,\dd x$。写出代换变量的范围并检验正切是否在此范围产生额外奇点，再讨论被积函数在 $x=\pm1$ 的极限及可补值性。
\end{exercise}
\begin{exercise}{458}{对数根式的复合积分}
求
\[\int\sqrt{\frac{\ln(x+\sqrt{1+x^2})+1}{1+x^2}}\,\dd x.\]
说明对数真数在全实数范围恒为正，再由外层根式确定实定义范围。选择对数整体作新变量，完成积分并确定满足 $F(0)=0$ 的原函数；说明为何不能忽略根式对定义域的额外限制。
\end{exercise}
\begin{exercise}{459}{指数根式代换与补值}
求 $\int\mathrm e^{\sqrt{1+\sin x}}\cos x/\sqrt{1+\sin x}\,\dd x$。在 $1+\sin x>0$ 的各开区间给出原函数，再讨论被排除点处原函数能否连续、可导延拓，说明根式中零点对链式求导的影响。
\end{exercise}
\begin{exercise}{460}{拆项与反正切复合}
求 $\int[x+(\arctan x)^2]/(1+x^2)\,\dd x$，说明为何应先拆分再分别代换。确定满足 $F(0)=1$ 的原函数，并用求导检验平方项积分后应出现的幂次与系数。
\end{exercise}
\begin{exercise}{461}{三角幂与有理化}
求 $\int\sin x\cos^3x/(1+\cos^2x)\,\dd x$，选取代换把它化为有理函数积分。再给出所得原函数的最小正周期，并说明原函数的周期性为何需要实际检查，不能仅由被积函数周期性推出。
\end{exercise}
\begin{exercise}{462}{正切代换的分支与全局原函数}
求 $\int\dd x/(\sin^2x+2\cos^2x)$，先在 $(-\pi/2,\pi/2)$ 使用正切代换。讨论直接使用反正切主值表达式为何在 $\pi/2$ 附近需要调整常数，并说明如何拼接得到全实数范围的一个原函数。
\end{exercise}
\begin{exercise}{463}{指数代换与两种等价表达式}
求 $\int\dd x/(1+\mathrm e^x)$，分别用 $u=\mathrm e^x$ 和 $v=\mathrm e^{-x}$ 代换。证明两种所得原函数只相差常数，并求满足 $F(0)=0$ 的表达式。
\end{exercise}
\begin{exercise}{464}{对数差的整体代换}
求 $\int[\ln(x+1)-\ln x]/[x(x+1)]\,\dd x$。先确定实定义域，计算对数差的导数后整体代换。讨论若改成 $\ln|x+1|-\ln|x|$，原函数公式可延伸到哪些额外区间。
\end{exercise}
\begin{exercise}{465}{对数代换后的多项式拆分}
求 $\int\ln^2x/[x(1+\ln^2x)]\,\dd x$，把代换后的有理函数先拆分再积分。求满足 $F(1)=0$ 的原函数，并证明它在 $(0,+\infty)$ 单调递增，须处理导数为零的点。
\end{exercise}
\begin{exercise}{466}{共轭有理化与根式分支}
求 $\int x/(x-\sqrt{x^2-1})\,\dd x$。先确定实定义域，使用共轭因子化简，分别说明在 $x>1$ 与 $x<-1$ 上计算所得原函数时哪些符号与对数绝对值必须保留。
\end{exercise}
\begin{exercise}{467}{指数复合与二倍角恒等式}
求 $\int\mathrm e^{\sin2x}\sin^2x/\mathrm e^{2x}\,\dd x$。先合并指数，再利用二倍角恒等式将其导数与剩余因子联系，给出在全实数范围成立的原函数，并说明关键系数的来源。
\end{exercise}
\begin{exercise}{468}{对数正切的整体代换}
在 $0<x<\pi/2$ 上求 $\int\ln(\tan x)/\sin2x\,\dd x$。证明 $\ln\tan x$ 在该区间严格递增且值域为 $\mathbb R$，用它作新变量积分，并确定满足 $F(\pi/4)=0$ 的原函数，须核对二倍角产生的系数。
\end{exercise}
\begin{exercise}{469}{商函数求导的逆向识别}
求 $\int[x+\sin x\cos x]/(\cos x-x\sin x)^2\,\dd x$。通过对 $\sin x/(\cos x-x\sin x)$ 求导寻找原函数，并说明答案应限制在哪些开区间；不要求显式求出所有分母零点。
\end{exercise}
\begin{exercise}{470}{三角指数乘积的有理结构}
求 $\int(\cos^2x-\sin x)/\{\cos x[1+\mathrm e^{\sin x}\cos x]\}\,\dd x$。以 $u=\mathrm e^{\sin x}\cos x$ 作整体代换，把积分化为有理函数形式；说明各分母非零条件与对数绝对值的作用，不要求显式求解全部奇点。
\end{exercise}
\begin{exercise}{471}{半角组合与指数乘积导数}
求 $\int(1+\sin x)\mathrm e^x/(1+\cos x)\,\dd x$。用半角恒等式将被积函数写成 $\mathrm e^x$ 与正切半角函数的乘积导数，给出最大实定义区间，并核验原函数在每个区间成立。
\end{exercise}
\begin{exercise}{472}{已知原函数后的复合换元}
在 $0<u<\pi$ 已知 $\int f(u)\,\dd u=\ln(\sin u)+C$。求 $\int xf(1-x^2)\,\dd x$，给出由 $0<1-x^2<\pi$ 确定的实定义区间。解释为何先由已知条件恢复 $f$ 或直接使用原函数换元可得到一致结果。
\end{exercise}
\begin{exercise}{473}{含抽象函数的商法则逆运算}
设 $f\in C^1(I)$，$0\notin I$。证明
\[\int\frac{xf'(x)-(1+x)f(x)}{x^2\mathrm e^x}\,\dd x=\frac{f(x)}x\mathrm e^{-x}+C.\]
明确使用的乘积与商法则，并用 $f(x)=x\mathrm e^x$ 检查两侧。若 $I$ 跨过 $0$，说明公式的适用区间必须如何调整。
\end{exercise}
\begin{exercise}{474}{偶次余项与根式代换}
求 $\int x^3\sqrt{4-x^2}\,\dd x$，使用 $u=4-x^2$ 后先将剩余 $x^2$ 消去。给出实积分区间，并确定满足 $F(0)=0$ 的原函数；说明代换在 $0$ 两侧不单调为何不妨碍最终求导核验。
\end{exercise}
\begin{exercise}{475}{平方和根式与双曲型对数}
求 $\int x^2/(1+x^2)^{3/2}\,\dd x$。分别用三角代换与拆成 $1/\sqrt{1+x^2}-1/(1+x^2)^{3/2}$ 的方法，证明所得答案相差常数；确定满足 $F(0)=0$ 的原函数。
\end{exercise}
\begin{exercise}{476}{平方差根式与参数绝对值}
设 $a>0$，求 $\int x^2/(x^2-a^2)^{3/2}\,\dd x$。在 $x>a$ 上使用三角代换，并给出在 $x<-a$ 上也成立的带绝对值表达式，明确两侧常数可独立选择。
\end{exercise}
\begin{exercise}{477}{根式有理化与变量范围}
求 $\int\dd x/[1+\sqrt[3]{x+1}]$，立方根按实数定义，采用 $t=\sqrt[3]{x+1}$ 并写出 $\dd x$。给出全部最大实定义区间，再检查所得原函数在 $x=-1$ 处是否可导，说明代换的逆函数导数在该点不存在是否妨碍最终答案成立。
\end{exercise}
\begin{exercise}{478}{不同根指数的统一有理化}
求 $\int\dd x/(\sqrt x+\sqrt[3]{x^2})$。先确定实定义区间，再根据两个根指数的最小公倍数选取代换，写出有理化后的完整积分。比较分别只代换平方根或立方根时仍会留下的根式，并检查最终结果的导数。
\end{exercise}
\begin{exercise}{479}{根式分式的有理代换}
求 $\int(3-x)^{-1}\sqrt{(5-x)/(x-1)}\,\dd x$。用根式本身作新变量，写出反解关系及新变量范围；给出被积函数的各个实定义区间，说明 $x=3$ 对积分常数的影响。
\end{exercise}
\begin{exercise}{480}{指数代换后的部分分式}
求 $\int\dd x/(\mathrm e^x+\mathrm e^{2x})$，将它化为新变量的有理函数积分。确定满足 $F(0)=0$ 的原函数，并由被积函数符号判断其单调性。
\end{exercise}
\begin{exercise}{481}{指数加常数的根式代换}
求 $\int\dd x/\sqrt{\mathrm e^x+1}$，令 $t=\sqrt{\mathrm e^x+1}$ 后写出完整有理积分。给出新变量范围和实对数表达式，并说明为何不能把 $\dd x$ 误写为 $2t\dd t$。
\end{exercise}
\begin{exercise}{482}{幂代换与不连通区间}
求 $\int\dd x/[x(x^6+1)]$。分别使用 $t=x^6$ 与先拆成 $1/x-x^5/(x^6+1)$ 的方法，比较所得表达式。解释 $x>0$、$x<0$ 两个区间中 $\ln|x|$ 和积分常数的处理。
\end{exercise}
\begin{exercise}{483}{倒代换与根式绝对值}
求 $\int\dd x/[x^4\sqrt{x^2+1}]$。尝试 $t=1/x$，明确 $\sqrt{x^2+1}$ 变换后出现的绝对值，分别在正负半轴积分；最终给出可在两个半轴分别使用的表达式并求导核验。
\end{exercise}
\begin{exercise}{484}{对数与代数的混合内函数}
求 $\int(1+x)/(x^2+x\ln x)\,\dd x$。选择 $u=x+\ln x$ 化简，证明分母在正半轴只有一个零点，从而说明答案成立的最大开区间，不要求写出该零点的初等表达式。
\end{exercise}
\begin{exercise}{485}{配方、三角代换与半圆结构}
求 $\int\sqrt{2x-x^2}\,\dd x$，先配方并给出实定义范围。确定满足 $F(1)=0$ 的原函数，并利用被积函数关于 $x=1$ 的对称性检验所得原函数的奇对称关系。
\end{exercise}
\begin{exercise}{486}{指数根式与两步代换}
求 $\int\dd x/\sqrt{1+\mathrm e^{2x}}$，先作指数代换，再选用根式或三角代换。写出对数形式的实数原函数，说明所有对数真数在全实数范围的符号并求导核验。
\end{exercise}
\begin{exercise}{487}{三角代换后的有理积分}
求 $\int\dd x/[(1+x^2)\sqrt{1-x^2}]$。在 $(-1,1)$ 作 $x=\sin t$，再使用正切代换完成积分。明确两次代换的变量范围，避免反函数分支导致常数跳跃。
\end{exercise}
\begin{exercise}{488}{参数退化与复合换元}
对 $a\in\mathbb R$ 求 $\int x/(1+ax^2)\,\dd x$，完整处理 $a=0$。在包含 $0$ 的最大实定义区间上选取 $F_a(0)=0$ 的原函数，并求 $a\to0$ 时 $F_a(x)$ 的点态极限，检验参数退化的一致性。
\end{exercise}
\begin{exercise}{489}{反三角主值与分支常数}
在 $(-1,1)$ 上求 $\int\dd x/\sqrt{1-x^2}$。若代换 $x=\sin t$ 分别取 $t\in(-\pi/2,\pi/2)$ 与 $t\in(\pi/2,3\pi/2)$，写出两种计算过程，并说明 $\sqrt{\cos^2t}=|\cos t|$ 与反函数分支如何共同保证所得函数只相差常数。
\end{exercise}
\subsection{分部积分法}
\begin{exercise}{490}{多项式与指数三角函数的分部积分}
求 $\int x\mathrm e^{ax}\,\dd x$、$\int x^2\cos x\,\dd x$、$\int(x^2-x)\sin x\,\dd x$，其中 $a\in\mathbb R$。第一式须单独处理 $a=0$，其余两式明确每次分部积分中次数降低的部分，并通过求导核验符号。
\end{exercise}
\begin{exercise}{491}{反函数与对数函数的分部选择}
求 $\int x/\sin^2x\,\dd x$、$\int\ln(x+\sqrt{x^2+1})\,\dd x$、$\int\arctan x\,\dd x$。写出每式分部积分中 $u,\dd v$ 的选择，说明第一式的定义区间与第二式对数真数的符号。
\end{exercise}
\begin{exercise}{492}{循环积分、恒等式与抽象函数}
求 $\int\mathrm e^{-x}\cos x\,\dd x$ 与 $\int x\cot^2x\,\dd x$。另设 $\int f(x)\,\dd x=(1+\sin x)\ln x+C$（$x>0$），求 $\int xf'(x)\,\dd x$；说明第三式为何可利用已知原函数而不必完整展开 $f'$。
\end{exercise}
\begin{exercise}{493}{对数移位与多项式分部}
求 $\int x\ln(x+1)\,\dd x$，给出实定义区间。分别采用直接分部积分与先令 $u=x+1$ 的方法，验证两种答案之差为常数，再确定满足 $F(0)=0$ 的原函数。
\end{exercise}
\begin{exercise}{494}{平移反正弦与根式符号}
求 $\int\arcsin(x+1)\,\dd x$，明确可导原函数的开区间。将答案连续延拓至实定义域端点，求两个端点的单侧导数，说明被积函数在端点有限与自身导数无界的区别。
\end{exercise}
\begin{exercise}{495}{凑微分与分部积分的配合}
求 $\int x\cot x/\sin^2x\,\dd x$，先把 $\cot x/\sin^2x$ 识别为一个基本函数的导数，再进行分部积分。注明积分成立的区间，并解释只作 $u=\cot x$ 代换为何仍会留下需要处理的 $x$。
\end{exercise}
\begin{exercise}{496}{半角恒等式后的分部积分}
求 $\int x\cos^4(x/2)/\sin^3x\,\dd x$。先用半角公式化简三角因子，再选择分部积分；明确原被积函数的定义域，并检查化简是否在原来无定义的点造成表面约分。
\end{exercise}
\begin{exercise}{497}{对数三角函数与分部积分}
在 $0<x<\pi$ 上求 $\int\ln(\sin x)/\sin^2x\,\dd x$。通过分部积分把剩余项化为基本三角积分，并用求导检验。说明若要延伸到 $\sin x<0$ 的区间，题面中的对数必须作何调整。
\end{exercise}
\begin{exercise}{498}{根式积分的两种路线}
求 $\int x^3/\sqrt{1+x^2}\,\dd x$，分别用整体根式代换与取 $u=x^2$、$\dd v=x\dd x/\sqrt{1+x^2}$ 的分部积分法。证明两种结果等价，并说明本题中哪一步使剩余积分的次数降低。
\end{exercise}
\begin{exercise}{499}{指数三角循环积分}
用两次分部积分求 $\int\mathrm e^{2x}\sin3x\,\dd x$。进一步推导实参数 $a,b$ 满足 $a^2+b^2>0$ 时 $\int\mathrm e^{ax}\sin bx\,\dd x$ 的公式，并检查 $b=0$ 与 $a=0$ 的退化情形。
\end{exercise}
\begin{exercise}{500}{平方和根式与参数退化}
设 $a>0$，用分部积分求 $\int\sqrt{x^2+a^2}\,\dd x$，将出现的同类积分移项求解，说明对数真数的正性。再单独求 $a=0$ 时的 $\int|x|\,\dd x$，并讨论两个半轴上的表达式如何拼接。
\end{exercise}
\begin{exercise}{501}{平方差根式与区间约束}
设 $a>0$，求 $\int\sqrt{a^2-x^2}\,\dd x$，分别用三角代换与分部积分核验。明确实定义域和反正弦的主值范围，再确定满足 $F(0)=0$ 的原函数及其奇偶性。
\end{exercise}
\begin{exercise}{502}{反正弦与剩余根式积分}
求 $\int x\arcsin x\,\dd x$，明确分部积分后的根式积分如何继续化简。选取满足 $F(0)=0$ 的原函数，检验它的奇偶性，并求两端连续延拓值。
\end{exercise}
\begin{exercise}{503}{反正切乘负幂与对数绝对值}
求 $\int x^{-2}\arctan x\,\dd x$，说明 $\int\dd x/[x(1+x^2)]$ 的分解。分别给出正负半轴上的原函数族，并讨论被积函数在 $x\to0$ 的主项是否允许原函数有限连续延拓。
\end{exercise}
\begin{exercise}{504}{对数分部与内部极点}
求 $\int\ln(1+x)/(2-x)^2\,\dd x$。先给出最大实定义区间，再作分部积分并处理所得有理积分。明确 $\ln(1+x)$ 与 $\ln|2-x|$ 的适用条件，说明 $x=2$ 两侧的常数为何独立。
\end{exercise}
\begin{exercise}{505}{指数根式的分部与有理化}
在 $x>0$ 上求 $\int x\mathrm e^x/\sqrt{\mathrm e^x-1}\,\dd x$。先把指数根式凑为一个微分，再对剩余根式积分换元。确定能连续延拓至 $0$ 且取值为 $0$ 的原函数，并检查右导数。
\end{exercise}
\begin{exercise}{506}{对数与根式的双重处理}
求 $\int\ln x/\sqrt{1+x}\,\dd x$（$x>0$），用分部积分后作 $t=\sqrt{1+x}$。说明对数中 $t-1$ 的符号，并讨论所求原函数在 $x\to0^+$ 时是否具有有限极限。
\end{exercise}
\begin{exercise}{507}{反正切根式与端点补值}
求 $\int\arctan\sqrt x\,\dd x$，给出在 $x>0$ 上的计算过程。选择满足 $F(0)=0$ 的连续延拓，核验其在 $0$ 的右导数与被积函数在该点的值相等。
\end{exercise}
\begin{exercise}{508}{有理指数乘积的导数识别}
求 $\int x^2\mathrm e^x/(x+2)^2\,\dd x$。把被积函数写成 $\mathrm e^x$ 与一个有理函数的导数结构，可设原函数为 $\mathrm e^x(A+B/(x+2))$ 求系数。注明定义区间，并用求导检查该设定确实覆盖本题。
\end{exercise}
\begin{exercise}{509}{三角有理组合与逆向求导}
求 $\int[(1+x)\sin x+\cos x]/(1+x)^2\,\dd x$。分别用分部积分的消项与识别商函数导数的方法，说明两条路线的等价性，并明确 $x=-1$ 两侧的积分常数。
\end{exercise}
\begin{exercise}{510}{复合导函数恢复与分部积分}
设 $f'(\mathrm e^t)=\cos t$ 对所有实数 $t$ 成立。求 $f(x)$ 的全部表达式，并由 $f(1)=0$ 确定常数。解题时先写出 $f'$ 的定义区间，再处理含 $\cos\ln x$ 的积分。
\end{exercise}
\begin{exercise}{511}{正割幂的递推与初值}
在不含 $\cos x=0$ 的开区间上，令 $I_n=\int\sec^nx\,\dd x$。对整数 $n\ge2$ 证明
\[I_n=\frac{\sin x}{(n-1)\cos^{n-1}x}+\frac{n-2}{n-1}I_{n-2},\]
其中等式按相差任意常数理解。求 $I_0,I_1$ 并用递推式求 $I_3,I_4$，说明递推为何需要两个初始积分。
\end{exercise}
\begin{exercise}{512}{参数分部积分与零参数极限}
对实参数 $a$ 求 $\int x\mathrm e^{ax}\,\dd x$，选取满足 $F_a(0)=0$ 的原函数。单独处理 $a=0$，并用泰勒公式证明非零参数表达式在 $a\to0$ 时趋于该结果，解释积分常数选择对极限的影响。
\end{exercise}
\begin{exercise}{513}{对数幂积分的递推}
设 $n$ 为非负整数，$J_n=\int(\ln x)^n\,\dd x$（$x>0$）。建立 $J_n$ 与 $J_{n-1}$ 的递推关系，并求 $J_3$。证明所得表达式可取有限的 $x\to0^+$ 极限，不能仅以每个对数项都无界判断原函数发散。
\end{exercise}
\subsection{有理函数的积分}
\begin{exercise}{514}{假分式除法与余式积分}
求 $\int x^3/(x^2+1)\,\dd x$，先写出多项式除法的商与余式。确定满足 $F(0)=0$ 的原函数，再用奇偶性检验该原函数的对称性。
\end{exercise}
\begin{exercise}{515}{不同一次因子的部分分式}
求 $\int(3x-2)/(x^2+3x-4)\,\dd x$，写出分母分解和部分分式系数。给出所有最大实定义区间，说明所得对数的绝对值及各区间常数的含义。
\end{exercise}
\begin{exercise}{516}{三不同实根分母的积分}
求 $\int(2x^2-x+1)/[(x+1)(x-3)(x+2)]\,\dd x$。用赋值法或待定系数法完整确定三个一次分式系数，再以通分核验分子恒等式，明确原函数的各个定义区间。
\end{exercise}
\begin{exercise}{517}{简单极点与原函数族}
求 $\int\dd x/[x(x-1)]$。在三个最大定义区间上分别描述全部原函数，判断是否能找到可连续延拓至 $0$ 或 $1$ 的原函数，并由对数主项说明理由。
\end{exercise}
\begin{exercise}{518}{重一次因子的完整分解}
求 $\int(x^2+1)/[(x+1)(x^2-1)]\,\dd x$。先按重根结构列全部分分式，说明若只保留 $1/(x+1)$ 与 $1/(x-1)$ 两项为何不足，再求系数并积分。
\end{exercise}
\begin{exercise}{519}{不可约二次因子的分子形式}
求 $\int\dd x/[(x^2+1)(1+2x)]$。写出包含二次因子对应一次分子的完整分解，积分后同时保留对数与反正切项。说明若把该一次分子擅自设为常数，会遗漏什么约束。
\end{exercise}
\begin{exercise}{520}{配方与分子拆成导数}
求 $\int(3x+1)/(x^2+2x+5)\,\dd x$，把分子拆成分母导数的倍数与常数。证明分母对所有实数非零，并确定满足 $F(-1)=0$ 的原函数。
\end{exercise}
\begin{exercise}{521}{重二次分母与标准积分}
求 $\int\dd x/(x^2+2x+5)^2$，先配方再使用三角代换或分部积分。写出有理项和反正切项，最后对整个表达式求导核验平方分母的系数。
\end{exercise}
\begin{exercise}{522}{四次分母分解与对称分子}
求 $\int(2x^3+2x^2+5x+5)/(x^4+5x^2+4)\,\dd x$。先分解分母，再利用分子的结构减少待定系数数量；说明为何答案可在全实数轴使用同一个积分常数。
\end{exercise}
\begin{exercise}{523}{配方后重二次因子的积分}
求 $\int x^2/(x^2+2x+2)^2\,\dd x$。先作平移代换，把分子拆成适合积分的三项，并给出完整原函数。通过检查大 $|x|$ 时被积函数的衰减阶，判断所求原函数在两端是否具有有限极限。
\end{exercise}
\begin{exercise}{524}{高次有理函数的实因式分解}
求 $\int(x^4+1)/(1+x^6)\,\dd x$。利用 $1+x^6=(x^2+1)(x^4-x^2+1)$，将四次因子进一步分解为两个实不可约二次因子，再作部分分式积分。注明实定义区间，并利用被积函数的偶性检查可取为奇函数的原函数。
\end{exercise}
\begin{exercise}{525}{高重零点与偶函数分解}
求 $\int\dd x/(x^4-x^6)$。先利用偶性寻找仅含 $1/x^2$、$1/x^4$ 与 $1/(1-x^2)$ 的分解，再积分。列出全部被排除点，说明不能将 $x=0$ 当作普通的简单极点处理。
\end{exercise}
\begin{exercise}{526}{倒幂代换与重因子结构}
求 $\int(1+2x^4)/[x^3(1+x^4)^2]\,\dd x$。先作 $t=1/x^2$ 以降低次数，再积分得到实原函数，分别说明正负半轴上的适用性，并通过求导核验重因子对应的有理项。
\end{exercise}
\begin{exercise}{527}{四次平方和的实分解}
求 $\int\dd x/(x^4+1)$，将分母分解为两个实不可约二次因子后积分。给出在 $0$ 附近连续且取值为 $0$ 的表达式，并讨论合并两个反正切为一个主值反正切时可能需要补加的常数。
\end{exercise}
\begin{exercise}{528}{原函数存在性与奇偶性}
逐项判断并证明或举例：初等函数在其连续的定义区间内存在原函数；不连续函数也可能存在原函数；奇函数在关于 $0$ 对称的开区间上的原函数都是偶函数；偶函数在这样的区间上的原函数都是奇函数。对最后一项给出使之成立的积分常数条件，并说明“存在原函数”与“有初等原函数”有何区别。
\end{exercise}
\begin{exercise}{529}{闭区间连续性保证的结论}
设 $f\in C[a,b]$，$a<b$。判断 $f$ 在 $(a,b)$ 内是否必有导函数、原函数、最大值或最小值、局部极值。对必然成立者说明定理依据；对其余结论给出具体反例，注意闭区间最值定理的取值范围。
\end{exercise}
\begin{exercise}{530}{分段函数原函数的参数拼接}
设 $f(x)=2(x-1)$（$x<1$）、$f(x)=\ln x$（$x\ge1$）。求实数 $a,b$ 使
\[F_{a,b}(x)=\begin{cases}(x-1)^2,&x<1,\\x(\ln x+a)+b,&x\ge1\end{cases}\]
为 $f$ 在 $\mathbb R$ 上的原函数。必须分别检查拼接处连续、左右导数及非拼接处导数，不能仅由两侧分别可积确定结论。
\end{exercise}
\begin{exercise}{531}{反正切指数函数的分部积分}
求 $\int\mathrm e^{-2x}\arctan(\mathrm e^x)\,\dd x$。先作指数代换，再通过分部积分处理反正切项，明确新变量恒正对对数绝对值的影响。
\end{exercise}
\begin{exercise}{532}{反正切与有理因子的配合}
求 $\int\arctan x/[x^2(1+x^2)]\,\dd x$。先分解 $1/[x^2(1+x^2)]$，再分别采用分部与整体代换；写出在两个半轴上的原函数族，并讨论 $0$ 处是否可连续延拓。
\end{exercise}
\begin{exercise}{533}{指数有理混合的对数导数}
求 $\int(1+x)/[x(1+x\mathrm e^x)]\,\dd x$，尝试将被积函数写成两个对数导数之差。证明 $1+x\mathrm e^x$ 在全实数范围恒正，从而给出答案的全部最大实定义区间。
\end{exercise}
\begin{exercise}{534}{对数商函数的逆向求导}
求 $\int(1-\ln x)/(x-\ln x)^2\,\dd x$。尝试对 $x/(x-\ln x)$ 求导识别原函数，并用最值证明分母在正半轴不为零，确定全实定义区间。
\end{exercise}
\begin{exercise}{535}{指数根式与反正切分部积分}
在 $x>0$ 上求 $\int\mathrm e^{2x}\arctan\sqrt{\mathrm e^x-1}\,\dd x$。先令 $t=\sqrt{\mathrm e^x-1}$ 化为多项式乘反正切的积分，再作分部积分。明确代换后 $\mathrm e^{2x}\dd x$ 的表达，避免遗漏雅可比因子。
\end{exercise}
\begin{exercise}{536}{三角商函数的导数结构}
求 $\int[x+\sin x\cos x]/(x\sin x+\cos x)^2\,\dd x$。通过求导 $\cos x/(x\sin x+\cos x)$ 寻找原函数，说明答案的适用开区间，不要求用初等函数写出分母的所有零点。
\end{exercise}
\begin{exercise}{537}{三角有理函数与万能代换}
求 $\int\dd x/(\sin2x+2\sin x)$。先因式分解并确定全部奇点，再作 $t=\tan(x/2)$ 积分。说明万能代换遗漏的点是否都是原函数的奇点，以及积分常数如何在定义区间上选择。
\end{exercise}
\begin{exercise}{538}{万能代换与可去的参数断点}
求 $\int\dd x/(1+\sin x+\cos x)$，用 $t=\tan(x/2)$ 化为有理函数。列出原被积函数的全部奇点，区别原函数的奇点与代换坐标在奇数倍 $\pi$ 处的失效，说明如何检验答案。
\end{exercise}
\begin{exercise}{539}{反正切原函数的跨区间拼接}
求 $\int\dd x/(2+\cos x)$，先在 $(-\pi,\pi)$ 采用万能代换。原被积函数在全实数连续，说明为何得到的反正切主值公式仍需按周期补常数，并给出拼接为全实数原函数的方法。
\end{exercise}
\begin{exercise}{540}{常数项与对数导数拆分}
求 $\int(3\cos x-\sin x)/(\cos x+\sin x)\,\dd x$，把分子拆成分母及其导数的线性组合。给出每个最大定义区间的原函数，并明确对数中绝对值的必要性。
\end{exercise}
\begin{exercise}{541}{对称拆分的三角积分}
求 $\int\sin x/(\sin x+\cos x)\,\dd x$，利用分母与其导数拆分。再求对应的 $\int\cos x/(\sin x+\cos x)\,\dd x$，检验两者相加与被积函数恒等式一致，说明常数的合并方式。
\end{exercise}
\begin{exercise}{542}{根式与线性项的三角代换}
求 $\int\dd x/[2x+\sqrt{1-x^2}]$，在 $(-1,1)$ 使用 $x=\sin t$。求出分母在该区间的唯一零点，分别给出两侧原函数，并说明代换中根式为何可写为 $\cos t$。
\end{exercise}
\begin{exercise}{543}{奇偶拆分与重二次分母}
求 $\int(x^3+1)/(1+x^2)^2\,\dd x$，先把奇函数部分与偶函数部分分开处理。确定满足 $F(0)=0$ 的原函数，并说明最终答案中哪些项分别来自这两部分。
\end{exercise}
\begin{exercise}{544}{反正切指数与根式分母}
求 $\int x\mathrm e^{\arctan x}/(1+x^2)^{3/2}\,\dd x$，先令 $t=\arctan x$，注意 $\cos t$ 的符号，再完成指数三角积分。说明该原函数可在全实数范围成立的依据。
\end{exercise}
\begin{exercise}{545}{指数几何和的因式分解}
求 $\int\dd x/(1+\mathrm e^{x/2}+\mathrm e^{x/3}+\mathrm e^{x/6})$。根据指数分母的最小公倍数选取 $t=\mathrm e^{x/6}$ 化为有理积分，利用 $1+t+t^2+t^3=(1+t)(1+t^2)$ 作部分分式分解，再积分并给出实定义范围，须明确 $\dd x$ 中的系数。
\end{exercise}
\begin{exercise}{546}{三角因子与有理代换}
求 $\int\dd x/[(2+\cos x)\sin x]$，先确定全部实定义区间，再用 $u=\cos x$ 化为有理函数积分。解释代换时 $\sin^2x=1-u^2$ 的作用，并核验所得对数系数。
\end{exercise}
\begin{exercise}{547}{反正弦分部与根式多项式}
求 $\int x^3\arcsin x/\sqrt{1-x^2}\,\dd x$。先求不含反正弦因子的原函数，再作分部积分，明确实定义区间。说明若先令 $t=\arcsin x$，为何会转化为多项式三角幂的分部积分。
\end{exercise}
\begin{exercise}{548}{隐函数约束下的原函数}
设 $y=y(x)$ 是关系 $y(x-y)^2=x$ 在某开区间上的可导分支，且 $x-3y\ne0$。求 $\int\dd x/(x-3y(x))$，可引入 $u=x-y$ 作参数。推导 $x,y$ 关于 $u$ 的表达式与 $\dd x$，并明确参数化需要排除的值及积分公式的分支适用条件。
\end{exercise}
\begin{exercise}{549}{函数关系的单射区间与换元}
设 $f(\sin^2t)=t/\sin t$ 对 $0<t<\pi/2$ 成立。求 $\int\sqrt x\,f(x)/\sqrt{1-x}\,\dd x$，明确 $f$ 在 $(0,1)$ 上由条件唯一确定的依据。先作 $x=\sin^2t$，再完成分部积分，并用直接恢复 $f(x)$ 的方法核验。
\end{exercise}
\begin{exercise}{550}{已知原函数与高阶导数积分}
设 $\int f(x)\,\dd x=\mathrm e^{x^2}+C$。求 $f(x)$、$f'(x)$ 及 $\int xf''(x)\,\dd x$。最后一个积分先用分部积分写为 $xf'(x)-f(x)$ 再代入，说明为什么不必展开 $f''$。
\end{exercise}
\begin{exercise}{551}{加权原函数恢复与奇点}
已知 $\int xf(x)\,\dd x=\arcsin x+C$ 在 $(-1,0)$ 或 $(0,1)$ 上成立。求 $f(x)$ 及 $\int f(x)\,\dd x$，给出带绝对值的实表达式。说明在包含 $0$ 的开区间上是否可能存在连续 $f$ 满足原条件。
\end{exercise}
\begin{exercise}{552}{指数复合导函数的恢复}
设 $\int f'(\mathrm e^t)\,\dd t=-(1+t)\mathrm e^{-t}+C$，且 $f(1)=0$。求 $f(x)$（$x>0$），分别说明对已知积分式求导和将 $x=\mathrm e^t$ 代回的步骤，再用初值确定常数。
\end{exercise}
\begin{exercise}{553}{抽象函数商的平方导数}
设 $f\in C^2(I)$，且 $f'(x)\ne0$。证明
\[\int\left[\frac{f(x)}{f'(x)}-\frac{f^2(x)f''(x)}{[f'(x)]^3}\right]\,\dd x=\frac12\left[\frac{f(x)}{f'(x)}\right]^2+C.\]
逐项检查商法则与链式法则产生的分母幂次，并解释条件 $f'\ne0$ 的必要性；再用 $f(x)=\mathrm e^x$ 验证。
\end{exercise}
\begin{exercise}{554}{线性三角组合的平方分母}
设实数 $a,b$ 满足 $a^2+b^2>0$。证明
\[\int\frac{\dd x}{(a\sin x+b\cos x)^2}=-\frac1{a^2+b^2}\frac{a\cos x-b\sin x}{a\sin x+b\cos x}+C.\]
给出结果成立的区间条件，并分别检查 $a=0$、$b=0$。再说明若 $a=b=0$，为何问题本身失去定义。
\end{exercise}
\begin{exercise}{555}{参数因子与部分分式退化}
对 $a\in\mathbb R$ 求 $\int\dd x/[(x-1)(x-a)]$，完整区分 $a=1$ 与 $a\ne1$。再取 $0<a<2$，在公共区间 $x<0$ 选取满足 $F_a(-1)=0$ 的原函数，求 $a\to1$ 时 $F_a(x)$ 的极限，说明积分常数归一化对该极限的作用。
\end{exercise}
\begin{exercise}{556}{一次因子与重二次因子的联合分解}
求 $\int\dd x/[x(x^2+1)^2]$，先列全部分分式并利用奇偶性减少系数。给出两个半轴上的全部原函数，最后通过求导及 $x\to0$ 的主项检查对数系数是否正确。
\end{exercise}

\clearpage
\section{第5章 定积分}
\subsection{定积分的概念与性质}
\begin{exercise}{557}{积分变量与参数求导}
设常数 $-1<a<b<1$，记 $I=\int_a^b\arcsin t\dd t$。分别求 $I$ 对无关变量 $x$、上限 $b$、下限 $a$ 的导数。若 $a=a(x),b=b(x)$ 为可导函数，写出 $\frac\dd{\dd x}\int_{a(x)}^{b(x)}\arcsin t\dd t$，说明哑变量命名不影响结果。
\end{exercise}
\begin{exercise}{558}{几何意义与尺度变换}
利用定积分的几何意义求 $\int_{-2}^{2}\sqrt{4-x^2}\dd x$，并推广计算 $\int_{-r}^{r}\sqrt{r^2-x^2}\dd x$（$r>0$）。若仅将积分区间改为 $[0,r]$，解释结果与原积分的关系。
\end{exercise}
\begin{exercise}{559}{黎曼和与端点影响}
利用定积分定义求 $\lim_{n\to\infty}\sum_{k=1}^{n-1}\frac{k}{n^2}$。对 $a>0$，再求 $\lim_{n\to\infty}\sum_{k=1}^{\lfloor an\rfloor}\frac{k}{n^2}$，说明取整及有限个端点项为何不影响极限。
\end{exercise}
\begin{exercise}{560}{黎曼和与积分夹逼}
将 $\sum_{k=1}^n\frac1{n+k}$ 写成黎曼和并求其极限。再利用 $f(x)=1/(1+x)$ 的单调性，证明该和与其极限的差的绝对值不超过 $1/(2n)$。
\end{exercise}
\begin{exercise}{561}{黎曼和与误差控制}
利用定积分定义求 $S_n=\sum_{k=1}^n\frac{n}{n^2+k^2}$ 的极限。将 $S_n$ 与相应积分比较，确定误差符号，并证明误差的绝对值不超过 $1/(2n)$。
\end{exercise}
\begin{exercise}{562}{积分比较与严格大小关系}
不计算原函数，分别比较以下各组积分：$\int_0^1x^2\dd x$ 与 $\int_0^1x^3\dd x$；$\int_0^1x\dd x$ 与 $\int_0^1\tan x\dd x$；$\int_1^2\ln x\dd x$ 与 $\int_1^2\ln^2x\dd x$；$\int_0^1e^x\dd x$ 与 $\int_0^1(1+x)\dd x$。给出每组严格不等式的依据，并指出端点相等是否影响结论。
\end{exercise}
\begin{exercise}{563}{根式积分的比较与区间影响}
不求原函数，比较 $I_1=\int_0^1\sqrt{1+x^2}\dd x$ 与 $I_2=\int_0^1\sqrt{1+x^3}\dd x$，证明严格大小关系。若积分区间同时改为 $[1,2]$，关系如何变化？说明理由。
\end{exercise}
\begin{exercise}{564}{积分估值与严格不等式}
不求原函数，先给出 $\int_{\pi/4}^{\pi}(1+\cos^2x)\dd x$ 的严格上、下界，所用界仅由被积函数最值与区间长度构成。再计算准确值，说明被积函数在有限个点达到最值为何不能使估值取等号。
\end{exercise}
\begin{exercise}{565}{带权积分估值}
不求原函数，证明 $\frac1{10\sqrt2}<\int_0^1\frac{x^9}{\sqrt{1+x}}\dd x<\frac1{10}$。把 $x^9$ 改为 $x^n$（$n$ 为非负整数），给出相应严格界，并求 $\lim_{n\to\infty}(n+1)\int_0^1\frac{x^n}{\sqrt{1+x}}\dd x$。
\end{exercise}
\begin{exercise}{566}{三角积分估值与单调性}
设 $I=\int_{\pi/4}^{\pi/2}\frac{\sin x}{x}\dd x$。证明 $\frac{\sqrt2}{4}\le I\le\frac{\sqrt2}{2}$，并判断两端能否取等号。再利用 $\sin x/x$ 的单调性改进下界。
\end{exercise}
\begin{exercise}{567}{积分中值定理与罗尔定理}
设 $a<b$，$f$ 在 $[a,b]$ 上连续、在 $(a,b)$ 可导，且 $\frac2{b-a}\int_a^{(a+b)/2}f(x)\dd x=f(b)$。证明存在 $\xi\in(a,b)$ 使 $f'(\xi)=0$，并说明可先在左半区间选出哪个中值点，再在哪个区间应用罗尔定理。
\end{exercise}
\begin{exercise}{568}{积分平均值的单调性}
设 $f$ 在 $[a,b]$ 上连续、在 $(a,b)$ 内可导，且 $f'\le0$。证明
\[
F(x)=\frac1{x-a}\int_a^xf(t)\dd t
\]
在 $(a,b)$ 内满足 $F'(x)\le0$。对固定的 $x_0\in(a,b)$，证明 $F'(x_0)=0$ 当且仅当 $f$ 在 $[a,x_0]$ 上为常数。
\end{exercise}
\begin{exercise}{569}{函数平均值与中值点}
求 $f(x)=\sqrt{4-x^2}$ 在 $[-2,2]$ 上的平均值，并求使函数值等于平均值的全部点。将半径 $2$ 改为 $r>0$ 后，求相应中值点与 $r$ 的比例关系。
\end{exercise}
\begin{exercise}{570}{积分定义与黎曼和}
设 $f$ 在 $[0,1]$ 上连续。用定积分表示极限
\[
\lim_{n\to\infty}\sum_{k=1}^n\frac{n}{n^2+k^2}f\!\left(\frac{k}{n}\right).
\]
当 $f(x)=x$ 时求极限值；当 $f(x)\ge0$ 且上述极限为零时，证明 $f\equiv0$。
\end{exercise}
\begin{exercise}{571}{积分不等式与等号条件}
设 $f$ 在 $[0,1]$ 上连续，且 $0\le f(x)\le1$、$\int_0^1f(x)\dd x=\frac12$。证明
\[
\frac14\le\int_0^1f^2(x)\dd x<\frac12,
\]
并完整讨论左端等号条件及右端不能取等号的原因。
\end{exercise}
\subsection{微积分基本公式}
\begin{exercise}{572}{变积分限与哑变量}
设 $F(a,b)=\int_a^b\sin t^2\dd t$。分别求其对无关变量 $x$、上限 $b$、下限 $a$ 的导数。再令 $a=x,b=x^2$，求复合函数的导数，说明不能把被积函数中的哑变量误作外部变量。
\end{exercise}
\begin{exercise}{573}{积分上限函数与凹凸性}
设 $y=\int_0^{x^2}\sqrt{1+t^2}\dd t$。求 $y'$、$y''$，判断 $y$ 的奇偶性及单调性，并判断是否存在拐点。
\end{exercise}
\begin{exercise}{574}{双变限积分的求导}
设 $y=\int_x^{x^2}\frac{\dd t}{\sqrt{1+t^2}}$。求 $y'$，并在 $x\ge0$ 上确定 $y$ 的零点与符号。计算 $y'(0)$、$y'(1)$，解释积分方向对符号判断的影响。
\end{exercise}
\begin{exercise}{575}{变限积分极限与高阶小量}
求 $\lim_{x\to0}\frac{\int_0^x\cos t^2\dd t}{x}$。进一步求 $\lim_{x\to0}\frac{x-\int_0^x\cos t^2\dd t}{x^5}$，并说明使用的局部展开及余项阶数。
\end{exercise}
\begin{exercise}{576}{变限积分极限与泰勒展开}
求 $\lim_{x\to0}\frac{\int_0^x\arctan t\dd t}{x^2}$，以及 $\lim_{x\to0}\frac{\int_0^x\arctan t\dd t-x^2/2}{x^4}$。论证从两侧趋于零时所得结果一致。
\end{exercise}
\begin{exercise}{577}{变限积分商的极限}
求
\[
\lim_{x\to0}\frac{\left(\int_0^xe^{t^2}\dd t\right)^2}{\int_0^xte^{2t^2}\dd t}.
\]
先证明分母在充分小的非零 $x$ 处不为零；再将分子中的指数 $t^2$ 改为 $at^2$、分母中的指数改为 $bt^2$，讨论实参数 $a,b$ 是否影响极限。
\end{exercise}
\begin{exercise}{578}{含参积分与连续性条件}
设 $f$ 在实轴上连续，求 $\lim_{x\to0}\frac{\int_0^xf(t)(x-t)\dd t}{x^2}$。分别用 $t=xu$ 与两次求导说明结果；再讨论将 $f$ 在原点连续弱化为在原点两侧有有限极限时，两侧极限相等的条件。
\end{exercise}
\begin{exercise}{579}{无穷远变限积分极限}
求 $\lim_{x\to+\infty}\frac{\int_0^x(\arctan t)^2\dd t}{\sqrt{1+x^2}}$。证明一般结论：若连续函数 $g(x)\to L$（$x\to+\infty$），则 $x^{-1}\int_0^xg(t)\dd t\to L$，并用该结论重新求本题极限。
\end{exercise}
\begin{exercise}{580}{隐函数与变限积分求导}
设 $\int_0^ye^t\dd t+\int_0^x\cos t\dd t=0$ 在原点附近确定隐函数 $y=y(x)$。求 $y'$、$y''$ 及其在 $x=0$ 的值，并确定包含 $0$ 的最大实定义区间。
\end{exercise}
\begin{exercise}{581}{高斯型变限积分的性质}
设 $f(x)=\int_0^xe^{-t^2/2}\dd t$，$x\in\mathbb R$。求其单调区间、奇偶性与凹凸区间；不用求出极限准确值，证明 $f$ 在实轴上有界。
\end{exercise}
\begin{exercise}{582}{含参变限积分与数据恢复}
设 $f$ 连续，$\varphi(x)=\int_0^{x^2}xf(t)\dd t$，且 $\varphi(1)=1$、$\varphi'(1)=5$。求 $f(1)$。若再知 $f$ 在 $1$ 处可导、$\varphi''(1)=9$，求 $f'(1)$，写清各次求导中的链式法则。
\end{exercise}
\begin{exercise}{583}{变限积分方程的唯一零点}
设 $f$ 在 $[a,b]$ 上连续且恒正，$a<b$。证明 $\int_a^xf(t)\dd t+\int_b^x\frac{\dd t}{f(t)}=0$ 在 $(a,b)$ 恰有一根。若 $f\equiv c>0$，求该根，并判断其关于 $c$ 的单调性。
\end{exercise}
\begin{exercise}{584}{基本公式与参数积分}
计算 $\int_{-1}^{1/2}e^{2x}\dd x$。对实参数 $a$，计算 $J(a)=\int_{-1}^{1/2}e^{ax}\dd x$，单独处理 $a=0$，并证明所得表达式在 $a=0$ 连续。
\end{exercise}
\begin{exercise}{585}{原函数与参数上限}
计算 $\int_0^2\frac{x}{\sqrt{1+x^2}}\dd x$。令 $G(a)=\int_0^a\frac{x}{\sqrt{1+x^2}}\dd x$，$a\in\mathbb R$，求 $G$ 并判断其奇偶性及最小值。
\end{exercise}
\begin{exercise}{586}{部分分式与积分取值}
计算 $\int_0^1\frac{\dd x}{x^2+5x+6}$。再令上限为 $a\ge0$，求积分关于 $a$ 的表达式，并确定其值域。
\end{exercise}
\begin{exercise}{587}{绝对值积分与参数最值}
计算 $\int_{-1}^3|2-x|\dd x$。对实数 $a$ 求 $H(a)=\int_{-1}^3|a-x|\dd x$ 的分段表达式，并求 $H$ 的最小值及取到最小值的全部 $a$。
\end{exercise}
\begin{exercise}{588}{绝对值积分与周期性}
计算 $\int_0^{2\pi}|\sin x|\dd x$。证明对任意实数 $a$ 与正整数 $n$，$\int_a^{a+n\pi}|\sin x|\dd x$ 与 $a$ 无关，并求其值。
\end{exercise}
\begin{exercise}{589}{分段函数的变限积分与可导性}
设 $f(x)=2x$（$0\le x\le1$），$f(x)=2+x$（$1<x\le2$）。求 $P(x)=\int_0^xf(t)\dd t$ 在 $[0,2]$ 上的表达式，判断 $P$ 在 $x=1$ 处的连续性和可导性，并求该处左右导数。
\end{exercise}
\begin{exercise}{590}{积分方程与自洽条件}
连续函数 $f$ 在 $[0,1]$ 上满足 $f(x)=\frac1{1+x^2}+\sqrt{1-x^2}\int_0^1f(t)\dd t$。求 $\int_0^1f(t)\dd t$ 及 $f$。若根式项前再乘实参数 $\lambda$，讨论连续解的存在性与唯一性。
\end{exercise}
\begin{exercise}{591}{积分方程与参数可解性}
设 $f$ 在 $[0,1]$ 上连续，满足 $f(x)=x\int_0^1f(t)\dd t-1$。求 $\int_0^1f(t)\dd t$ 及 $f$。若 $x$ 前的系数改为实数 $\lambda$，讨论方程何时存在唯一连续解。
\end{exercise}
\begin{exercise}{592}{复合函数恢复与换元}
设 $f$ 在 $[2,+\infty)$ 上满足 $f(x+1/x)=\frac{x+x^3}{1+x^4}$（$x>0$）。证明此条件唯一确定 $f$，求 $f$ 的表达式及 $\int_2^{2\sqrt2}f(u)\dd u$。说明 $x$ 与 $1/x$ 给出的函数值为何一致。
\end{exercise}
\begin{exercise}{593}{变限积分与极值}
设 $F(x)=\int_0^x(t^2-1)e^{-t}\dd t$，$x\ge0$。求 $F$ 的单调区间与极值，并计算 $\lim_{x\to\infty}F(x)$；由此确定 $F$ 在 $[0,+\infty)$ 上的零点个数。
\end{exercise}
\begin{exercise}{594}{积分方程与条件恢复}
设连续函数 $f$ 满足
\[
f(x)=1+\int_0^x(x-t)f(t)\dd t\qquad(x\in\mathbb R).
\]
证明 $f$ 二阶可导，求 $f$；再直接代回积分等式核验。
\end{exercise}
\subsection{定积分的换元积分法和分部积分法}
\begin{exercise}{595}{根式换元与定义域}
计算 $\int_4^9\frac{\sqrt x}{\sqrt x-1}\dd x$。对 $1<a<b$，求同一被积函数在 $[a,b]$ 上的积分，并说明当下限趋于 $1^+$ 时积分的变化。
\end{exercise}
\begin{exercise}{596}{根式换元与积分函数}
计算 $\int_0^1x\sqrt{3-2x}\dd x$。对 $0\le a\le3/2$，求 $J(a)=\int_0^ax\sqrt{3-2x}\dd x$，并确定 $J$ 在该区间上的最大值。
\end{exercise}
\begin{exercise}{597}{对数换元与定义域}
计算 $\int_1^e\frac{\dd x}{x\sqrt{1+\ln x}}$。将上限改为 $a>e^{-1}$，求积分表达式；再判断上限固定为 $1$、下限趋于 $e^{-1}$ 时的积分是否有限。
\end{exercise}
\begin{exercise}{598}{复合换元与单调条件}
计算 $\int_{-1}^0\frac{(1+x)e^x}{\sqrt{1+xe^x}}\dd x$。说明换元 $u=1+xe^x$ 在积分区间上是否合法；再求将下限改为 $a\in[-1,0]$ 后的积分及其关于 $a$ 的单调性。
\end{exercise}
\begin{exercise}{599}{三角换元与几何意义}
计算 $\int_0^{\sqrt2/3}\sqrt{2-9x^2}\dd x$，并用椭圆面积核验。进一步计算 $\int_0^{\sqrt2/3}x^2\sqrt{2-9x^2}\dd x$。
\end{exercise}
\begin{exercise}{600}{三角换元与无穷上限}
计算 $\int_1^{\sqrt3}\frac{\dd x}{x\sqrt{1+x^2}}$。求 $J(a)=\int_1^a\frac{\dd x}{x\sqrt{1+x^2}}$（$a\ge1$），并求 $\lim_{a\to+\infty}J(a)$。
\end{exercise}
\begin{exercise}{601}{根式积分与反三角函数}
计算 $\int_1^2\frac{\sqrt{x^2-1}}x\dd x$。将上限改为 $a\ge1$，求积分表达式；由此求 $\lim_{a\to+\infty}\left(\int_1^a\frac{\sqrt{x^2-1}}x\dd x-a\right)$。
\end{exercise}
\begin{exercise}{602}{根式积分与对称换元}
计算 $\int_{-1}^1\frac{\dd x}{1+\sqrt{1-x^2}}$。对 $-1\le a\le1$ 求 $J(a)=\int_{-a}^a\frac{\dd x}{1+\sqrt{1-x^2}}$，正确处理 $a<0$ 时的积分方向，并判断 $J$ 的奇偶性。
\end{exercise}
\begin{exercise}{603}{配方换元与圆面积}
计算 $\int_0^1\sqrt{2x-x^2}\dd x$。利用被积函数关于 $x=1$ 的对称性求 $\int_0^2x\sqrt{2x-x^2}\dd x$，并说明不必另求原函数的理由。
\end{exercise}
\begin{exercise}{604}{指数根式换元与积分方向}
计算 $\int_0^{-\ln2}\sqrt{1-e^{2x}}\dd x$。设 $J(a)=\int_{-a}^0\sqrt{1-e^{2x}}\dd x$（$a>0$），求 $J(a)$ 及 $\lim_{a\to+\infty}(J(a)-a)$，先由积分方向判断原积分的符号。
\end{exercise}
\begin{exercise}{605}{平移换元与分段积分}
设 $f(x)=1+x^2$（$x<0$），$f(x)=e^{-x}$（$x\ge0$）。求 $\int_1^3f(x-2)\dd x$；令 $J(c)=\int_1^3f(x-c)\dd x$，证明 $J$ 在 $c=2$ 可导并求 $J'(2)$。
\end{exercise}
\begin{exercise}{606}{分段函数的平移与奇偶消去}
设 $f(x)=xe^{x^2}$（$-1/2\le x\le1/2$），$f(x)=-1$（$x>1/2$）。求 $\int_{1/2}^2f(x-1)\dd x$。说明换元后如何利用奇偶性简化；若仅改变 $f(1/2)$ 的值，积分是否改变？给出理由。
\end{exercise}
\begin{exercise}{607}{卷积换元与函数恢复}
设连续函数 $f$ 满足 $\int_0^xf(x-t)\dd t=e^{-2x}-1$。求 $\int_0^1f(x)\dd x$ 及 $f(x)$，并证明此连续解唯一。
\end{exercise}
\begin{exercise}{608}{含参积分与二次求导}
设 $f$ 在实轴上连续，且 $\int_0^xf(x-t)t\dd t=e^x-x-1$。通过换元和微积分基本定理求 $f$，并直接核验原等式。说明推导时为何不需要预先假设 $f$ 可导。
\end{exercise}
\begin{exercise}{609}{奇函数积分与参数平移}
计算 $\int_{-\pi/2}^{\pi/2}\frac{x}{1+\cos x}\dd x$。将分子改为 $x+c$，求积分关于实参数 $c$ 的表达式，并确定积分为零的充要条件。
\end{exercise}
\begin{exercise}{610}{对数函数的奇偶性}
计算 $\int_{-\pi/2}^{\pi/2}\sin^2x\ln(x+\sqrt{1+x^2})\dd x$。先证明对数因子是奇函数；若将 $\sin^2x$ 改为任意连续函数 $g(x)$，说明积分只与 $g$ 的奇部还是偶部有关并证明。
\end{exercise}
\begin{exercise}{611}{对称换元与参数条件}
计算 $\int_{-1}^1\frac1{\sqrt{4-x^2}}\left(\frac1{1+e^x}-\frac12\right)\dd x$。若 $\frac12$ 改为实参数 $c$，求积分关于 $c$ 的表达式及其为零的充要条件。
\end{exercise}
\begin{exercise}{612}{反三角函数的对称关系}
计算 $\int_{-1}^1\cos x\arccos x\dd x$。证明对任意连续偶函数 $g$，$\int_{-1}^1g(x)\arccos x\dd x=\frac\pi2\int_{-1}^1g(x)\dd x$，并说明本题如何应用此式。
\end{exercise}
\begin{exercise}{613}{奇偶消去与积分排序}
设三积分的区间均为 $[-\pi/4,\pi/4]$，被积函数分别为 $\frac{\tan x}{1+x^4}+x^8$、$\sin^8x+\ln(x+\sqrt{1+x^2})$、$\tan^8x+e^x\cos x-e^{-x}\cos x$，对应积分为 $M,N,P$。不用求原函数，确定三者严格大小次序；证明奇函数部分消去后所用的点态不等式，并说明严格性的来源。
\end{exercise}
\begin{exercise}{614}{周期函数与绝对值}
对正整数 $n$，计算 $\int_0^{n\pi}\sqrt{1-\sin^2x}\dd x$。再求 $\int_a^{a+n\pi}\sqrt{1-\sin^2x}\dd x$，证明结果与实数 $a$ 无关，并指出去掉根号而直接写成 $\cos x$ 的错误。
\end{exercise}
\begin{exercise}{615}{周期表达式的瑕点与截断}
考察 $I=\int_0^\pi\sin^2x(\tan x+1)\dd x$。指出区间内部的瑕点，按反常积分定义判断 $I$ 是否存在；再求对称截断和 $\lim_{\varepsilon\to0^+}(\int_0^{\pi/2-\varepsilon}+\int_{\pi/2+\varepsilon}^{\pi})\sin^2x(\tan x+1)\dd x$，说明两种概念的区别。
\end{exercise}
\begin{exercise}{616}{对称换元恒等式与指数积分}
设 $f$ 在 $[0,1]$ 上连续。证明 $\int_0^{\pi/2}f(\sin x)\dd x=\int_0^{\pi/2}f(\cos x)\dd x$ 及 $\int_0^\pi xf(\sin x)\dd x=\pi\int_0^{\pi/2}f(\sin x)\dd x$。利用这两式计算 $I_n=\int_0^\pi\frac{x\sin^{2n}x}{\sin^{2n}x+\cos^{2n}x}\dd x$（$n$ 为正整数），证明结果与 $n$ 无关。
\end{exercise}
\begin{exercise}{617}{分部积分与参数恒等式}
计算 $\int_0^1\ln(1+x^2)\dd x$。对 $a>0$ 计算 $\int_0^a\ln(1+x^2)\dd x$，再对所得结果求导核验。
\end{exercise}
\begin{exercise}{618}{反三角函数的分部积分}
计算 $\int_0^{1/2}\arcsin x\dd x$。将上限改为 $a\in[-1,1]$，求积分表达式并判断其奇偶性；由所得结果求该积分在允许参数范围内的最大值。
\end{exercise}
\begin{exercise}{619}{分部积分与奇偶分解}
计算 $\int_{-1}^1x\arccos x\dd x$。分别用分部积分及 $\arccos(-x)=\pi-\arccos x$ 将积分化到 $[0,1]$ 上，说明两种计算的联系。
\end{exercise}
\begin{exercise}{620}{分部积分与有理积分}
计算 $\int_0^1\frac{\ln(1+x)}{(2-x)^2}\dd x$。先选择能使剩余积分化为有理函数的分部方式；再利用被积函数的符号核验所得结果为正。
\end{exercise}
\begin{exercise}{621}{循环分部积分}
计算 $\int_0^{\pi/2}e^{-x}\sin2x\dd x$。求 $F(a)=\int_0^ae^{-x}\sin2x\dd x$（$a\ge0$），并由 $F'(a)$ 判断 $a=\pi/2$ 是极大点还是极小点。
\end{exercise}
\begin{exercise}{622}{换元与分部积分的结合}
计算 $\int_0^1e^{-\sqrt x}\dd x$。对 $a>0$ 求 $\int_0^{a^2}e^{-\sqrt x}\dd x$，并求其在 $a\to+\infty$ 时的极限。
\end{exercise}
\begin{exercise}{623}{三角换元与端点处理}
计算 $\int_0^1\frac{x^2\arcsin x}{\sqrt{1-x^2}}\dd x$，在计算中说明 $x=1$ 处无界因子为何不导致发散；再求 $\int_0^1\frac{\arcsin x}{\sqrt{1-x^2}}\dd x$ 作为对照。
\end{exercise}
\begin{exercise}{624}{嵌套根式换元与分部积分}
计算 $\int_1^{16}\arctan\sqrt{\sqrt x-1}\dd x$。先说明 $u=\sqrt{\sqrt x-1}$ 将区间变为何区间，再用分部积分化为多项式积分；用被积函数取值范围给出结果应满足的粗略界。
\end{exercise}
\begin{exercise}{625}{倒数换元与指数分部积分}
计算 $\int_1^2x^{-3}e^{1/x}\dd x$。将上限改为 $a\ge1$ 求积分表达式，并判断该积分在 $a\to+\infty$ 时是否有有限极限，给出极限值。
\end{exercise}
\begin{exercise}{626}{根式换元与振荡符号}
计算 $\int_0^{\pi^2}\sqrt x\cos\sqrt x\dd x$。先换元化为多项式乘三角函数，再解释原被积函数虽在前半段为正、总积分仍可能为负的原因。
\end{exercise}
\begin{exercise}{627}{分部积分与端点数据}
设 $f\in C^2[0,2]$，$f(0)=1,f(2)=3,f'(2)=5$。求 $\int_0^2xf''(x)\dd x$，说明为什么无需知道 $f'(0)$。若再知 $f'(0)=a$，求 $\int_0^2(2-x)f''(x)\dd x$。
\end{exercise}
\begin{exercise}{628}{周期积分的常值与符号}
设 $F(x)=\int_x^{x+2\pi}e^{\sin t}\sin t\dd t$。证明 $F$ 是常函数，并确定其符号，不能只根据因子 $\sin t$ 在一周期积分为零判断。用正、负半周配对给出严格证明。
\end{exercise}
\begin{exercise}{629}{已知原函数与数据化积分}
已知 $e^{x^2}$ 是 $f(x)$ 的一个原函数。求 $\int_1^e\frac{\ln x\,f'(\ln x)}x\dd x$；计算中先换元，再用分部积分把结果化为原函数和 $f$ 的端点值，说明所需可导性由何保证。
\end{exercise}
\begin{exercise}{630}{有限秩积分方程}
设连续函数 $f$ 满足 $f(x)=x+2\int_0^{\pi/2}f(t)\cos t\dd t$。求 $f$ 及 $\int_0^1f(x)\dd x$。若积分前的系数 $2$ 改为实数 $\lambda$，讨论连续解的存在性和唯一性。
\end{exercise}
\begin{exercise}{631}{可去端点与积分次序转换}
设 $f(x)=\int_0^x\frac{\sin t}{\pi-t}\dd t$，$0\le x\le\pi$，在 $t=\pi$ 处按连续延拓取值。求 $\int_0^\pi f(x)\dd x$；先证明延拓存在，再用分部积分或交换累次积分次序计算，写明边界项。
\end{exercise}
\begin{exercise}{632}{变限积分的分部计算}
设 $f(x)=\int_1^xe^{-y^2}\dd y$。求 $\int_0^1x^2f(x)\dd x$，说明虽然 $e^{-x^2}$ 的原函数不能用初等函数表示，本题仍可通过分部积分化为初等积分；由 $f$ 的符号核验结果。
\end{exercise}
\begin{exercise}{633}{含导数积分方程与积分计算}
设 $f\in C^1[0,\pi/2]$，满足 $f(x)=x-\int_0^xf'(t)\cos t\dd t$。求 $f$ 及 $\int_0^{\pi/2}f(x)\dd x$，并代回原等式验证端点条件与函数表达式。
\end{exercise}
\begin{exercise}{634}{对称换元与参数积分}
对 $a>0$，计算
\[
I(a)=\int_0^{\pi/2}\frac{\sin x}{\sin x+a\cos x}\dd x.
\]
利用 $x\mapsto\frac\pi2-x$ 建立 $I(a)$ 与 $I(1/a)$ 的关系，并单独处理 $a=1$ 的情形。
\end{exercise}
\begin{exercise}{635}{分部积分与递推估计}
记 $I_n=\int_0^1x^ne^x\dd x$，$n=0,1,2,\ldots$。建立 $I_n$ 的递推关系，求 $I_3$，并证明 $\lim_{n\to\infty}(n+1)I_n=e$。
\end{exercise}
\subsection{反常积分}
\begin{exercise}{636}{无穷积分分类与判敛}
取 $a>1$，分别判断 $\int_a^{+\infty}x^{-1/2}\dd x$、$\int_a^{+\infty}\frac{\ln x}{x}\dd x$、$\int_a^{+\infty}\frac{\dd x}{x\ln x}$、$\int_a^{+\infty}xe^{-x}\dd x$ 的敛散性，给出理由；对收敛者求值。
\end{exercise}
\begin{exercise}{637}{对数换元与参数临界值}
计算 $\int_2^{+\infty}\frac{\dd x}{x\ln^2x}$。对 $a>1$、实参数 $p$，讨论 $\int_a^{+\infty}\frac{\dd x}{x\ln^px}$ 的敛散性，收敛时求值。
\end{exercise}
\begin{exercise}{638}{全轴反常积分与平移配方}
计算 $\int_{-\infty}^{+\infty}\frac{\dd x}{x^2+4x+5}$，写明全轴积分的定义。再求使 $\int_{-\infty}^{+\infty}\frac{\dd x}{x^2+4x+c}$ 收敛的实参数 $c$ 的范围，并在收敛时求值。
\end{exercise}
\begin{exercise}{639}{拆项积分与无穷端点抵消}
计算 $\int_1^{+\infty}\frac{\dd x}{x^2+x}$。说明为何部分分式分解后不能将积分写为两个发散积分之差；再求下限改为 $a>0$ 时的值及 $a\to0^+$ 时的极限。
\end{exercise}
\begin{exercise}{640}{反常根式积分与三角代换}
计算 $\int_3^{+\infty}\frac{\dd x}{(x-1)^4\sqrt{x^2-2x}}$。先用 $u=x-1$ 化简，再选择三角代换；求下限改为 $a>2$ 时的值，并判断 $a\to2^+$ 时的极限是否有限。
\end{exercise}
\begin{exercise}{641}{全轴换元与收敛证明}
计算 $\int_{-\infty}^{+\infty}\frac{e^x}{1+e^{2x}}\dd x$，分别说明两个无穷端点的收敛性。将指数中的 $x$ 换成 $ax$（$a\ne0$），求积分并讨论 $a$ 的符号。
\end{exercise}
\begin{exercise}{642}{分段积分与参数阈值}
计算 $\int_0^{+\infty}\min\{e^{-x},1/2\}\dd x$。将 $1/2$ 改为 $c>0$，分情况求积分关于 $c$ 的表达式，并核验其在 $c=1$ 连续。
\end{exercise}
\begin{exercise}{643}{反常分部积分与边界项}
计算 $\int_1^{+\infty}\frac{\ln x}{(1+x)^2}\dd x$。先证明积分收敛并处理分部积分中的无穷端边界项，再求下限改为 $a\ge1$ 时的积分。
\end{exercise}
\begin{exercise}{644}{三角换元与全轴推广}
计算 $\int_0^{+\infty}\frac{\dd x}{(1+x^2)^2}$。进一步求 $\int_{-\infty}^{+\infty}\frac{\dd x}{(a^2+x^2)^2}$（$a>0$），说明尺度代换中的系数。
\end{exercise}
\begin{exercise}{645}{反常三角换元}
计算 $\int_1^{+\infty}\frac{\dd x}{x\sqrt{2x^2-1}}$。再讨论下限改为 $a\ge1/\sqrt2$ 时的收敛性并求值，特别说明临界下限处的瑕点是否可积。
\end{exercise}
\begin{exercise}{646}{反常分部积分与对数积分}
计算 $\int_0^{+\infty}\frac{xe^{-x}}{(1+e^{-x})^2}\dd x$。先用比较法证明收敛，再选择使无穷端边界项为零的分部方式；最后计算将 $x$ 同时替换为 $ax$（$a>0$）后的积分。
\end{exercise}
\begin{exercise}{647}{奇函数与反常积分定义}
判断 $\int_{-\infty}^{+\infty}\sin2x\,e^{|x|}\dd x$ 的敛散性，并与对称截断积分 $\lim_{R\to+\infty}\int_{-R}^R\sin2x\,e^{|x|}\dd x$ 比较。证明两者不能仅由奇偶性等同；若 $e^{|x|}$ 改为 $e^{-|x|}$，判断是否绝对收敛并求积分值。
\end{exercise}
\begin{exercise}{648}{比较判别法与发散证明}
判断 $\int_2^{+\infty}\frac{\dd x}{\ln x}$ 的敛散性，给出明确的比较函数与成立区间。再讨论 $\int_2^{+\infty}\frac{\dd x}{x(\ln x)^p}$ 随实数 $p$ 的敛散性，说明增加 $x$ 因子后的变化。
\end{exercise}
\begin{exercise}{649}{高斯型积分的收敛性}
不用高斯积分的准确值，证明 $\int_0^{+\infty}e^{-x^2}\dd x$ 收敛。对 $a\in\mathbb R$ 讨论 $\int_0^{+\infty}e^{-ax^2}\dd x$ 的敛散性，并在 $a>0$ 时给出它与 $a=1$ 时积分的关系。
\end{exercise}
\begin{exercise}{650}{指数型极限与反常积分方程}
求实数 $c$，使 $\lim_{x\to+\infty}\left(\frac{x+c}{x-c}\right)^x=\int_{-\infty}^cte^{2t}\dd t$。分别计算两边，说明极限中幂底对充分大的 $x$ 为正，并证明参数解唯一。
\end{exercise}
\begin{exercise}{651}{瑕点识别与可去间断点}
分别考察 $\int_0^1\frac{\sin x}{x}\dd x$、$\int_0^1x\ln x\dd x$、$\int_{-1}^1\frac{\dd x}{\sqrt{x^2}}$、$\int_0^1x^{-2}\dd x$。指出未定义点，区分可作连续延拓的情形与无界瑕点，判断敛散性，并计算收敛积分中 $\int_0^1x\ln x\dd x$ 的值。
\end{exercise}
\begin{exercise}{652}{对数端点与连续延拓}
判断 $x\ln^2x$ 在 $x=0$ 是否无界；通过连续延拓计算 $\int_0^1x\ln^2x\dd x$。推广讨论 $\int_0^1x^p\ln^2x\dd x$ 随实参数 $p$ 的敛散性，收敛时求值。
\end{exercise}
\begin{exercise}{653}{端点瑕积分与参数条件}
对 $a>0$，计算 $\int_0^a\frac{\dd x}{\sqrt{a^2-x^2}}$，用极限说明上端点的处理。再讨论 $\int_0^a(a^2-x^2)^{-p}\dd x$ 随实参数 $p$ 的敛散性。
\end{exercise}
\begin{exercise}{654}{多瑕点与对数换元}
判断 $\int_0^1\frac{\dd x}{x(\ln x)^2}$ 的敛散性，分别考察 $0$ 和 $1$；再判断 $\int_0^{+\infty}\frac{\dd x}{x(\ln x)^2}$ 是否存在，指出全部需要拆分的瑕点，不能只检查无穷远。
\end{exercise}
\begin{exercise}{655}{瑕积分的幂次判别}
设 $a<b$，判断 $\int_a^b\frac{\dd x}{(x-a)^2}$ 的敛散性。推广讨论 $\int_a^b\frac{\dd x}{(x-a)^p}$（$p\in\mathbb R$），收敛时求值，并明确临界指数。
\end{exercise}
\begin{exercise}{656}{倒数换元与对数反常积分}
证明 $\int_0^{+\infty}\frac{\ln x}{1+x^2}\dd x$ 绝对收敛，再用 $x=1/t$ 求其值。说明在收敛性尚未建立前，不能直接依换元得到 $I=-I$ 就下结论。
\end{exercise}
\begin{exercise}{657}{分段反常积分与双重条件}
设 $f(x)=(x-1)^{1-\alpha}$（$1<x<e$），$f(x)=1/[x\ln^{\alpha+1}x]$（$x\ge e$）。求使 $\int_1^{+\infty}f(x)\dd x$ 收敛的全部实数 $\alpha$，并在收敛范围内求积分值。说明两个端点分别给出什么参数条件。
\end{exercise}
\begin{exercise}{658}{反常积分运算的有效条件}
分别判断 $\int_{-\infty}^{+\infty}\frac{x}{1-x^2}\dd x$、$\int_{-\infty}^{+\infty}\sin x\dd x$、$\int_{-1}^1x^{-2}\dd x$ 是否存在，指出用奇偶性或直接代原函数端点的失效条件。再计算 $\int_0^{+\infty}x^{3/2}e^{-x}\dd x$，可使用 $\Gamma(s+1)=s\Gamma(s)$、$\Gamma(1/2)=\sqrt\pi$，并说明其收敛性。
\end{exercise}
\begin{exercise}{659}{双端瑕点与参数判敛}
对实数 $p,q$，讨论反常积分
\[
\int_0^1\frac{\dd x}{x^p(1-x)^q}
\]
的敛散性，并在 $p=q=\frac12$ 时计算其值。说明为何必须分别考察两个端点。
\end{exercise}
\begin{exercise}{660}{对数因子与临界指数}
讨论 $\int_e^{+\infty}\frac{\dd x}{x(\ln x)^p}$ 与 $\int_0^{1/e}\frac{\dd x}{x|\ln x|^p}$ 随实参数 $p$ 的敛散性；收敛时分别求值，并用换元解释二者的联系。
\end{exercise}

\clearpage
\section{第6章 定积分的应用}
\subsection{定积分在几何学上的应用}
\begin{exercise}{661}{平面面积与形心}
求 $y=\sqrt x$、$y=0$、$x=4$ 围成区域的面积，并求该均匀薄片形心的横坐标。分别写出关于 $x$ 和关于 $y$ 的面积积分，核验二者相同。
\end{exercise}
\begin{exercise}{662}{横向切片与多边界区域}
求曲线 $x=\sqrt y$、直线 $y=x+1$、$y=1$ 与 $x$ 轴围成区域的面积，即 $0\le y\le1$、$y-1\le x\le\sqrt y$。求该均匀区域形心的横坐标，并说明直接选取关于 $y$ 的积分能避免什么分段。
\end{exercise}
\begin{exercise}{663}{区域面积与旋转体积}
求 $y=x^2$ 与 $y=x$ 围成区域的面积，再求该区域分别绕 $x$ 轴、$y$ 轴旋转一周所得体积。说明积分区间内两条曲线的位置关系。
\end{exercise}
\begin{exercise}{664}{交点、面积与平移轴旋转}
求 $y=-x^2+4$ 与 $y=-2x+4$ 围成区域的面积，并求该区域绕直线 $y=4$ 旋转所得体积；列式时明确内、外半径。
\end{exercise}
\begin{exercise}{665}{面积反求参数与尺度关系}
曲线 $y=\sqrt{ax}$（$a>0$）、$x$ 轴及 $x=3$ 围成的区域面积为 $6$。求 $a$，并求该区域绕 $x$ 轴旋转的体积。若直线改为 $x=b>0$ 且面积仍为 $6$，求 $a$ 与 $b$ 的关系。
\end{exercise}
\begin{exercise}{666}{指数曲线与积分区域判定}
求 $y=2^x$、$x+2y=2$、$x=2$ 围成区域的面积。证明相关交点唯一，并求该均匀区域对 $y$ 轴的静矩 $\iint_Dx\dd A$，可将其写成单积分计算。
\end{exercise}
\begin{exercise}{667}{多个区域的面积相加}
求两部分面积之和：一部分由 $y=x^2-x$ 与 $x$ 轴围成；另一部分在 $x>1$ 内，由 $y=2x+6$、$y=x^2-x$ 及 $x$ 轴围成。分别给出交点、积分区间及边界函数，并说明第二部分为何需要分段。
\end{exercise}
\begin{exercise}{668}{切线与抛物线间的面积}
求 $y=x^2$ 在 $(2,4)$ 处的切线，与 $y=-x^2+4x+1$ 围成区域的面积。将后一个抛物线的常数项 $1$ 改为 $c$，讨论何时形成面积为正的有界区域，并求其面积。
\end{exercise}
\begin{exercise}{669}{横向切片与区域分割}
抛物线 $y^2=9-x$ 与直线 $x=2$ 围成的区域被 $y=-1$ 分成两部分。求上方较大部分的面积，再求下方部分的面积，并核验二者之和为完整区域面积。
\end{exercise}
\begin{exercise}{670}{三角函数面积与参数上限}
求 $y=x\cos2x$（$0\le x\le\pi/4$）与 $x$ 轴围成的面积。令 $A(a)$ 为 $y=x\cos2x$ 在 $0\le x\le a$ 与 $x$ 轴间的面积，求 $0\le a\le\pi/2$ 时的分段表达式。
\end{exercise}
\begin{exercise}{671}{分支选择与对称面积}
由 $y=e^x$（$y\le1$）、$y=e^{-x}$（$y\le1$）、$x=\pm1$、$y=0$ 围成区域。求其面积及绕 $y$ 轴旋转所得体积，说明两条指数曲线各自对应的 $x$ 范围。
\end{exercise}
\begin{exercise}{672}{对数曲线与分段边界}
求 $y=x\ln x$、$y=x$、$y=0$ 围成区域的面积，并求绕 $x$ 轴旋转所得体积。列出所有边界交点，说明在 $x=1$ 处分段的原因。
\end{exercise}
\begin{exercise}{673}{三角曲线之间的面积}
求 $y=\sin x$、$y=\cos x$、$x=0$、$x=2\pi$ 围成区域的面积。再求将右端点改为 $2n\pi$（$n$ 为正整数）后的面积，并用周期性证明结果。
\end{exercise}
\begin{exercise}{674}{圆与抛物线分割的面积}
曲线 $y=x^2$、$x^2+y^2=2$ 与 $y=0$ 围成若干区域。求位于 $0\le y\le1$ 内、在抛物线外侧且在圆内的左右两块封闭区域的总面积，即 $\sqrt y\le|x|\le\sqrt{2-y^2}$。再求两区域绕 $y$ 轴旋转所得立体的体积，注意左右两块旋转后的重叠。
\end{exercise}
\begin{exercise}{675}{面积比与参数范围}
设 $0<b<4$。令 $D$ 为 $y=4x-x^2$ 与 $y=b$ 围成的抛物线弓形，令 $E$ 为 $y$ 轴、直线 $y=b$ 与抛物线靠近原点的一支围成的区域。已知 $|E|=|D|/2$，求 $b$，并证明在规定范围内解唯一。
\end{exercise}
\begin{exercise}{676}{椭圆分割与变量选择}
直线 $y=x$ 将椭圆 $x^2+3y^2-6y=0$ 内部区域分成两部分。求直线下方部分的面积，并求该部分占椭圆总面积的比例。说明选择关于 $y$ 的积分时需要在哪些交点分段。
\end{exercise}
\begin{exercise}{677}{相切约束与面积最优化}
抛物线 $y=px^2+qx$（$p<0,q>0$）在第一象限内与 $x+y=5$ 相切。设抛物线与 $x$ 轴围成的区域面积为 $S$。由相切条件求 $p,q$ 的约束及允许范围，再求使 $S$ 最大的 $p,q$ 与最大面积，证明最大值确能取到。
\end{exercise}
\begin{exercise}{678}{绝对值面积与唯一极小值}
在 $y=x^2$（$0\le x\le1$）上取 $P$，过 $P$ 作水平直线。求该直线与抛物线在 $x=0$、$x=1$ 之间所夹两部分的面积之和的最小值及 $P$ 坐标。把面积写为一个含绝对值的积分，并证明极小点唯一。
\end{exercise}
\begin{exercise}{679}{公切线与封闭区域}
求 $y=x^2$ 与 $y=-x^2+4x-4$ 的全部公切线及相应切点。对每条公切线，求其与两条抛物线围成区域的面积，并说明两种面积是否相同。
\end{exercise}
\begin{exercise}{680}{抛物线族与面积最小化}
抛物线 $y=ax^2+bx+c$ 经过 $(0,0)$、$(1,2)$，且 $a<0$。求使它与 $y=-x^2+2x$ 围成区域面积最小的 $a,b,c$；单独讨论 $a=-1$ 时是否存在有界区域，并说明最小值的取到情况。
\end{exercise}
\begin{exercise}{681}{极坐标面积与边界参数}
求极坐标曲线 $r=a\cos\theta+b$（$b\ge a>0$）围成区域的面积。若面积固定为 $S_0$，求允许的 $a$ 的范围及 $b$ 的表达式，包含 $b=a$ 的边界情形。
\end{exercise}
\begin{exercise}{682}{极坐标交集面积}
求 $r=a\sin\theta$ 与 $r=a(\cos\theta+\sin\theta)$（$a>0$）围成的两个圆盘的公共部分面积。明确公共区域的角度范围及径向上界切换点，并求该面积与 $a^2$ 的比值。
\end{exercise}
\begin{exercise}{683}{摆线参数积分与尺度关系}
摆线一拱为 $x=a(t-\sin t)$、$y=a(1-\cos t)$，$0\le t\le2\pi$，$a>0$。求它与 $x$ 轴围成的面积，并求绕 $x$ 轴旋转所得体积，说明参数端点处 $x'(t)=0$ 为何不妨碍积分。
\end{exercise}
\begin{exercise}{684}{参数面积函数与变化率最值}
点 $P(t)=(\cos t,2\sin^2t)$，$0\le t\le\pi/2$。设 $S(t)$ 为线段 $OP$、该曲线从 $P(t)$ 到 $(1,0)$ 的弧与 $x$ 轴围成的面积。求 $S(t)$ 及 $S'(t)$，确定 $S'(t)$ 取得最大值时 $P$ 的坐标，并核验两个参数端点。
\end{exercise}
\begin{exercise}{685}{双曲线截面与旋转体积}
求双曲线 $x^2/4-y^2/9=1$ 的两支与 $y=\pm3$ 围成区域绕 $y$ 轴旋转所得体积。改为 $y=\pm h$（$h>0$）后求体积表达式，说明两支旋转产生的立体为何只能计数一次。
\end{exercise}
\begin{exercise}{686}{旋转体的内外半径}
求 $xy=3$ 与 $x+y=4$ 在第一象限围成区域绕 $x$ 轴旋转所得体积，并用对称性判断绕 $y$ 轴所得体积是否相同。证明整个积分区间内外半径的次序。
\end{exercise}
\begin{exercise}{687}{无界区域的有限旋转体积}
设 $D=\{(x,y):x\ge0,\ 0\le y\le e^{-x}\}$。求 $D$ 分别绕 $x$ 轴与 $y$ 轴旋转的体积，先给出截断后再取极限的定义。再求曲线改为 $e^{-ax}$（$a>0$）时两体积的比值。
\end{exercise}
\begin{exercise}{688}{负半轴区域与柱壳半径}
求 $y=e^{-x^2}$、$x=0$、$x=-1$、$y=0$ 围成区域绕 $y$ 轴旋转所得体积。说明柱壳半径应取什么表达式；再将左端点改为 $-a$（$a>0$），求体积及其极限。
\end{exercise}
\begin{exercise}{689}{圆盘旋转与截面重叠}
圆盘 $x^2+(y-2)^2\le a^2$（$a>0$）绕 $x$ 轴旋转。分别在 $0<a\le2$ 与 $a>2$ 时求所得立体体积。先对固定 $x$ 判断截面是圆环还是圆盘；在 $a>2$ 时明确切换位置，不能将重叠部分重复计入。
\end{exercise}
\begin{exercise}{690}{切线确定与双轴旋转}
过原点作 $y=\ln x$ 的切线，求切点与切线方程。求切线、曲线及 $x$ 轴围成区域分别绕 $x$ 轴、$y$ 轴旋转所得体积，明确每个积分区间的边界函数。
\end{exercise}
\begin{exercise}{691}{无穷延展旋转体积}
求 $y=xe^{-x}$（$x\ge0$）与 $x$ 轴间区域绕 $x$ 轴旋转所得体积。证明体积收敛；将曲线改为 $xe^{-ax}$，讨论实参数 $a$ 使体积有限的条件并在有限时求值。
\end{exercise}
\begin{exercise}{692}{跨越旋转轴的截面合并}
求 $y=\sqrt{1-(x-2)^2}$ 与 $y=x^2-4x+3$ 围成区域分别绕 $x$ 轴、$y$ 轴旋转所得体积 $V_x,V_y$。在求 $V_x$ 时比较上下边界到轴的距离，说明为何不能直接使用两函数平方之差。
\end{exercise}
\begin{exercise}{693}{参数曲线与平移旋转轴}
星形线 $x=\cos^3t$、$y=\sin^3t$（$0\le t\le2\pi$）围成区域绕 $y=1$ 旋转。求体积，并说明可否用区域形心的运动路程乘区域面积核验；若旋转轴改为 $y=c\ge1$，求体积关于 $c$ 的关系。
\end{exercise}
\begin{exercise}{694}{面积反求切点与旋转体积}
在 $y=\sqrt x$ 的某点 $A$ 作切线，该切线与曲线及 $y$ 轴围成区域的面积为 $1/12$。求 $A$ 与切线方程；求该区域分别绕 $x$ 轴、$y$ 轴旋转所得体积，并证明满足面积条件的切点唯一。
\end{exercise}
\begin{exercise}{695}{累积体积与切线面积最大值}
设 $V(\xi)$ 为 $y=e^{-x}$、两坐标轴与 $x=\xi>0$ 围成区域绕 $x$ 轴旋转的体积。求 $V$ 及满足 $V(a)=\frac12\lim_{\xi\to+\infty}V(\xi)$ 的 $a$。再在该曲线 $x\ge0$ 的部分上求一点，使其切线与两坐标轴围成的三角形面积最大，求此面积并证明为全局最大值。
\end{exercise}
\begin{exercise}{696}{摆线的双轴旋转曲面面积}
摆线一拱为 $x=a(t-\sin t)$、$y=a(1-\cos t)$，$0\le t\le2\pi$，$a>0$。求弧长微元及该弧分别绕 $x$ 轴、$y$ 轴旋转所得曲面的面积。计算中说明绝对值 $|\sin(t/2)|$ 在此区间如何处理。
\end{exercise}
\begin{exercise}{697}{旋转曲面面积与分支计数}
求抛物线 $y^2=2x+3$ 的上支（$0\le x\le3$）绕 $x$ 轴旋转所得曲面的面积。若将上下两支一起旋转，所得曲面面积是否加倍？说明理由，并求仅上支的弧长积分表达式。
\end{exercise}
\begin{exercise}{698}{直角坐标弧长与参数上限}
求 $y=x\sqrt x$（$0\le x\le4$）的弧长。将右端点改为 $a>0$，求弧长 $L(a)$，并证明 $L(a)>a$。
\end{exercise}
\begin{exercise}{699}{对数三角曲线弧长}
求 $y=\ln\cos x$（$0\le x\le\pi/3$）的弧长。对 $0<a<\pi/2$ 求 $[0,a]$ 上的弧长，判断其在 $a\to\pi/2^-$ 时是否有限，给出理由。
\end{exercise}
\begin{exercise}{700}{变限积分定义曲线的弧长}
求曲线 $y=\int_0^x\tan t\dd t$（$0\le x\le\pi/4$）的弧长。比较它与 $y=\ln\cos x$ 在同一区间上的弧长，利用导数说明两者相同的原因。
\end{exercise}
\begin{exercise}{701}{渐开线参数弧长}
求圆的渐开线 $x=a(\cos t+t\sin t)$、$y=a(\sin t-t\cos t)$，$0\le t\le2\pi$、$a>0$ 的弧长。求使前段弧长恰为全长一半的参数值，并说明曲线初始点速度为零不影响结果。
\end{exercise}
\begin{exercise}{702}{极坐标弧长与指数参数}
求对数螺线 $r=e^\theta$（$0\le\theta\le\pi$）的弧长。推广求 $r=e^{a\theta}$ 在同一区间上的弧长，单独处理 $a=0$ 并验证表达式在该点连续。
\end{exercise}
\begin{exercise}{703}{根式曲线的参数化与弧长}
求 $\sqrt x+\sqrt y=1$ 的完整实曲线弧长。先给出不在端点产生无穷导数的参数方程，明确参数范围；并用曲线关于 $y=x$ 的对称性核验积分区间。
\end{exercise}
\begin{exercise}{704}{封闭区域的分段边界周长}
求 $y=x^2$、$y=2-x$ 与 $y$ 轴围成的、位于 $x\ge0$ 的区域边界周长。分别求曲线段、斜线段和竖直线段的长度，说明选择正交点的依据。
\end{exercise}
\begin{exercise}{705}{曲率半径与弧长参数求导}
在抛物线 $y=\sqrt x$ 上，$M(x,\sqrt x)$ 满足 $x\ge1$，$\rho(x)$ 为该点曲率半径，$s(x)$ 为从 $(1,1)$ 到 $M$ 的弧长。由 $\rho=(1+(y')^2)^{3/2}/|y''|$ 求 $\rho$、$\dd s/\dd x$，计算 $3\rho\frac{\dd^2\rho}{\dd s^2}-(\frac{\dd\rho}{\dd s})^2$，判断结果是否依赖 $x$。
\end{exercise}
\begin{exercise}{706}{圆弓形区域与轴外旋转}
设 $D=\{(x,y):x^2+y^2\le2x,\ y\ge x\}$。求 $D$ 绕 $x=2$ 旋转所得体积。明确 $x$ 的范围、柱壳半径与高度，并求 $D$ 的面积用于核验区域选择。
\end{exercise}
\begin{exercise}{707}{绝对值边界与旋转体积}
求 $y=3-|x^2-1|$ 与 $x$ 轴围成的封闭区域绕 $y=3$ 旋转所得体积。先求区域的横向范围，列出需要分段的位置，并解释内半径的表达式。
\end{exercise}
\begin{exercise}{708}{椭圆底面与已知截面体积}
正椭圆柱的底面为 $x^2/a^2+y^2/b^2\le1$。过底面短轴 $x=0$ 作与底面成 $\alpha\in(0,\pi/2)$ 的平面，在 $x\ge0$ 一侧它与底面间形成楔形体，高度为 $x\tan\alpha$。求该楔形体的体积，写出垂直于 $x$ 轴的截面面积，并说明体积对 $\tan\alpha$ 的依赖。
\end{exercise}
\begin{exercise}{709}{切线区域的面积与平移轴旋转}
过原点作 $y=\ln x$ 的切线，设切线、曲线与 $x$ 轴围成区域为 $D$。求 $D$ 的面积及绕直线 $x=e$ 旋转所得体积，明确柱壳法所用的半径与分段高度。
\end{exercise}
\begin{exercise}{710}{面积、体积与参数约束}
对 $a>0$，区域 $D_a$ 由 $y=x^2$、$y=ax$ 围成。求其面积及绕 $x$ 轴旋转所得立体的体积；若区域面积等于 $\frac43$，求 $a$ 与相应体积。
\end{exercise}
\begin{exercise}{711}{有限体积与无穷侧面积}
设 $D=\{(x,y):x\ge1,\ 0\le y\le1/x\}$。求 $D$ 绕 $x$ 轴旋转所得体积，并证明曲线 $y=1/x$（$x\ge1$）绕 $x$ 轴旋转所得侧面积为无穷大。两次计算均先截断到 $x=R$ 再取极限。
\end{exercise}
\subsection{定积分在物理学上的应用}
\begin{exercise}{712}{线性阻力与位移积分}
单位质量质点在水平直线上运动，仅受大小等于速率的阻力，初速度为 $v_0>0$。求速率降至 $v_0/3$ 的时间及此前路程，并求无限长时间内的总路程。根据牛顿第二定律写出速度方程，再通过积分求位移。
\end{exercise}
\begin{exercise}{713}{漏失载荷与变力做功}
井深 $30\,\mathrm m$，抓斗重 $400\,\mathrm N$，缆绳每米重 $50\,\mathrm N$，初始污泥重 $2000\,\mathrm N$。抓斗从井底以 $3\,\mathrm{m/s}$ 匀速提升，污泥以 $20\,\mathrm{N/s}$ 漏失。忽略抓斗高度及井口上方绳长，求提升至井口克服重力所做的功。分别写出提升高度为 $s$ 时污泥重量、悬垂绳重，核验污泥在全过程未漏尽。
\end{exercise}
\begin{exercise}{714}{矩形抛物线闸门的静水压力}
竖直闸门上部为宽 $2\,\mathrm m$、高 $h>0$ 的矩形，下部为深 $1\,\mathrm m$ 的抛物线弓形。以矩形下边为 $z=0$、向下为正，下部横截面为 $|x|\le\sqrt{1-z}$（$0\le z\le1$）；水面与矩形上边重合。求两部分的单侧水压力，若矩形部分与下部压力之比为 $5:4$，求 $h$ 并证明正解唯一。水密度为 $\rho$，重力加速度为 $g$。
\end{exercise}
\begin{exercise}{715}{变力做功与几何级数}
桩入土深度为 $x$ 时阻力为 $kx$（$k>0$）。气锤首次击打将桩打入 $a>0$，以后每次做功均为前一次的 $r$ 倍（$0<r<1$）。求击打 $3$ 次后的总深度及无限次击打可达到的极限深度；再求首次达到或超过给定深度 $d>0$ 所需的最少击打次数，说明哪些 $d$ 无法在有限次内达到。
\end{exercise}
\begin{exercise}{716}{横放椭圆柱油罐的质量}
油罐为长 $l$ 的椭圆柱，横放后端面为 $x^2/a^2+y^2/b^2\le1$，其中 $a>b>0$，长轴水平、短轴竖直。油面距罐底 $3b/2$，油密度为 $\rho$。求油的质量；将油深改为 $h\in[0,2b]$，写出质量 $M(h)$ 的积分并求 $M'(h)$，说明加注同一体积时油面升高速度在哪个油深最小。
\end{exercise}
\begin{exercise}{717}{分段容器体积与抽水功}
容器在高度 $y$ 处的水平截面为圆盘，半径平方为 $1-y^2$（$-1\le y\le1/2$）、$2y-y^2$（$1/2<y\le2$），即由两圆弧绕 $y$ 轴旋转得到，顶部抽出口在 $y=2$。求容积；若盛满水，将全部水抽至该出口至少需做多少功？水密度为 $10^3\,\mathrm{kg/m^3}$，重力加速度为 $g$，长度单位为米。写明每层水的提升距离，并利用截面对 $y=1/2$ 的对称性核验结果。
\end{exercise}
\begin{exercise}{718}{弧长与均匀薄片形心}
曲线 $L$ 为 $y=x^2/4-(\ln x)/2$，$1\le x\le e$。求 $L$ 的弧长。设 $D$ 由 $L$、$x=1$、$x=e$ 与 $x$ 轴围成，求其面积及形心横坐标，并先证明曲线在所给区间位于 $x$ 轴上方。
\end{exercise}
\begin{exercise}{719}{变力做功与动能定理}
质量为 $m$ 的物体沿水平直线运动，阻力大小为 $kv^2$，其中 $k>0$、$v$ 为速率。物体仅受此阻力，从速率 $v_0>0$ 减至 $v_0/2$。利用 $mv\dd v=-kv^2\dd s$ 求此过程中的位移、阻力做功，并核验动能定理。
\end{exercise}
\begin{exercise}{720}{液体压力与水深参数}
竖直薄板为底边水平的等腰三角形，高 $h$、底边长 $b$，顶点在底边上方。薄板顶点距水面 $d\ge0$，水的密度为 $\rho$。求单侧静水总压力 $P(d)$；若将薄板翻转使底边在上方且最高点仍在深度 $d$ 处，求压力变化量并解释其符号。
\end{exercise}

\clearpage
\section{第7章 微分方程}
\subsection{微分方程的基本概念}
\begin{exercise}{721}{微分方程辨识与阶数}
设 $y=y(x)$，分别判断 $x^2+x+1=0$、$y^2+2x+1=0$、$x^2y^2+y=3$、$x^2+y'+y=2$ 是否为微分方程；对符合者给出阶数及是否线性的判断，说明未知函数的代数幂次为何不能用来判定阶数。
\end{exercise}
\begin{exercise}{722}{导数阶数与非线性}
判断 $y''-2x(y')^3+y^2=0$ 的阶数及线性性质。给定 $y(0)=a,y'(0)=b$，求该解应满足的 $y''(0)$ 与 $y'''(0)$；解释 $(y')^3$ 的指数为何不能视为导数阶数。
\end{exercise}
\begin{exercise}{723}{微分形式与标准形式}
将 $x\dd y+(x+y)\dd x=0$ 在 $x>0$ 上写为标准形式并判断阶数。求通解及满足 $y(1)=0$ 的解，说明化为标准形式时对 $x=0$ 作了什么限制。
\end{exercise}
\begin{exercise}{724}{特解验证与通解区分}
验证 $y=5x^2$ 是 $xy'-2y=0$ 的解，说明它是通解还是特解。求该方程在 $x>0$ 上的通解及满足 $y(1)=5$ 的解，并解释为何讨论通解时需要指定区间。
\end{exercise}
\begin{exercise}{725}{解族验证与初值确定}
验证 $y=C_1\cos x+C_2\sin x$ 是 $y''+y=0$ 的解，并证明两个常数独立。求满足 $y(0)=2$、$y'(0)=3$ 的特解及其最大值。
\end{exercise}
\begin{exercise}{726}{任意常数个数与通解判定}
验证 $y=x^2/2+x+C$ 满足 $y''-y'=-x$。判断该一参数解族是否为通解；求完整通解，并给出一个不属于上述一参数族的解。说明仅验证代入成立为何还不能判定为通解。
\end{exercise}
\begin{exercise}{727}{隐函数解与正则分支}
在 $x>0$、$y>0$、$y\ne1$ 的正则分支上，验证 $y=\ln(xy)$ 所确定的函数满足
\[
x(y-1)y''+x(y')^2+(y-2)y'=0.
\]
说明该隐式关系给出通解还是特解，并解释为何必须排除 $y=1$ 处直接视作 $y(x)$ 求导。
\end{exercise}
\begin{exercise}{728}{解族、奇解与初值唯一性}
验证 $y=(x-C)^2$ 在 $x\ge C$ 上满足 $y'=2\sqrt y$，并验证 $y\equiv0$ 也是解。对初值 $y(0)=0$，构造至少两个在 $[0,+\infty)$ 上不同的解，说明直接分离变量可能遗漏什么解。
\end{exercise}
\begin{exercise}{729}{消去常数与微分方程阶数}
消去曲线族 $y=C_1e^x+C_2e^{-x}+x$ 中的任意常数，建立二阶微分方程。再求该方程满足 $y(0)=1$、$y'(0)=0$ 的特解，并核验两个初始条件。
\end{exercise}
\subsection{可分离变量的微分方程}
\begin{exercise}{730}{可分离变量方程辨识}
分别判断 $y'+xy=0$、$y^2\dd x+(2x+1)\dd y=0$、$y'=xye^{x+y}$、$(x-y)\dd x+(x+y)\dd y=0$ 是否可直接分离变量。对符合者写出分离形式，并检查分离时被除因子为零可能给出的常值解；说明最后一个方程的右端为何不能表示为单纯的 $x$ 函数与 $y$ 函数之积。
\end{exercise}
\begin{exercise}{731}{变量分离与奇点两侧解}
求 $y'=y(1-x)/x$ 在 $x>0$ 上的通解，包括零解；求满足 $y(1)=e^{-1}$ 的解。讨论该特解能否延拓为在全实轴上满足 $xy'=(1-x)y$ 的可导解。
\end{exercise}
\begin{exercise}{732}{变量分离与解的定义区间}
求 $y'=\sin x\cos x/y^2$ 的隐式通解。求满足 $y(0)=1$ 的解及其在实轴上的有界性；说明为什么通解中必须限制 $y\ne0$，而不能把零函数补作解。
\end{exercise}
\begin{exercise}{733}{分离变量与初值问题}
求 $(1+e^x)\dd y-ye^x\dd x=0$ 满足 $y(0)=2$ 的特解。推广至 $y(0)=a$，判断何种 $a$ 使解恒正、恒负或恒为零，并求 $y(x)/(1+e^x)$。
\end{exercise}
\begin{exercise}{734}{变量分离与平方根分支}
求 $(xy^2+x)\dd x+(y-x^2y)\dd y=0$、$y(0)=1$ 的特解。分离变量时说明需要避开的因子，确定包含 $0$ 的最大可导解区间，并解释不能将正、负平方根在零点任意拼接。
\end{exercise}
\begin{exercise}{735}{隐式解与初值分支}
求 $e^{x-y}\dd x+y\dd y=0$ 满足 $y(0)=1$ 的特解，以隐式形式表示即可。确定包含 $0$ 且使该解保持 $y>0$ 的最大区间，说明区间端点出现的原因。
\end{exercise}
\begin{exercise}{736}{切线斜率与曲线恢复}
曲线通过 $(0,-1)$，其任一点 $(x,y)$ 处切线斜率为 $y^2/(1+x^2)$。求曲线方程及包含 $0$ 的最大存在区间，并求曲线在 $x\to+\infty$ 时的水平渐近线（若该方向属于解的定义域）。
\end{exercise}
\begin{exercise}{737}{积分方程与链式求导}
设连续函数 $f$ 满足 $f(x)=\int_0^{3x}f(t/3)\dd t+2$。证明 $f$ 可导并求 $f$，再直接核验积分等式；若将外部常数 $2$ 改为 $a$，说明解如何改变。
\end{exercise}
\begin{exercise}{738}{含导数积分方程与非线性降阶}
设 $f\in C^1[0,+\infty)$，满足 $f(x)=2\int_0^x(x-t)f(t)f'(t)\dd t+x-1$。先恢复 $f(0)$，再化为一阶初值问题求 $f$；证明所得表达式在整个规定区间上有效。
\end{exercise}
\begin{exercise}{739}{卷积积分方程与初值恢复}
设连续函数 $f$ 满足 $\int_0^xf(t)\dd t=\int_0^xtf(x-t)\dd t+x$。通过换元将等式化为可求导形式，求 $f$，并说明初值从何而来。
\end{exercise}
\begin{exercise}{740}{非局部条件与常数确定}
设非零连续函数 $f$ 满足 $f(x)=\int_0^xf(t)\dd t+2\int_0^1tf^2(t)\dd t$。求 $f$，并证明非零条件下解唯一。若允许零函数，列出全部连续解并代回核验。
\end{exercise}
\begin{exercise}{741}{线性方程与积分条件}
设 $f'=xf/(1+x^2)$、$f(0)=a$，且 $\int_0^{\sqrt3}xf(x)\dd x=7$。求 $a$ 及 $f$，说明积分条件能否唯一确定 $a$；再求 $f$ 在实轴上的最小值。
\end{exercise}
\begin{exercise}{742}{非线性方程的平移代换}
求 $y'=y^2+2(\sin x-1)y+\sin^2x-2\sin x-\cos x+1$ 的通解。采用 $u=y+\sin x-1$ 化简，并补全分离变量可能遗漏的解；求满足 $y(0)=1$ 的解。
\end{exercise}
\begin{exercise}{743}{乘积代换与隐式通解}
求 $y(xy+1)\dd x+x(1+xy+x^2y^2)\dd y=0$ 的隐式通解。采用 $u=xy$，在 $xy\ne0$ 的区间上化为可分离形式并积分；另检查 $y\equiv0$，并通过全微分核验所得隐式关系。
\end{exercise}
\begin{exercise}{744}{平衡解与最大存在区间}
求初值问题 $y'=y^2$、$y(0)=a$ 的解，包括 $a=0$ 的情形；分别确定 $a>0$ 与 $a<0$ 时包含 $0$ 的最大存在区间，并讨论 $x\to+\infty$ 时解是否有意义。
\end{exercise}
\begin{exercise}{745}{受限增长与阈值时间}
种群满足 $N'=rN(1-N/K)$，其中 $r,K>0$、$N(0)=N_0\in(0,K)$。求 $N(t)$，证明其单调增加且趋于 $K$；当 $N_0<K/2$ 时，求首次达到 $K/2$ 的时间。
\end{exercise}
\subsection{齐次方程}
\begin{exercise}{746}{齐次方程辨识与求解}
分别判断 $y'+y=0$、$y^2\dd x+x\dd y=0$、$y'+xy=e^x$、$x^2y'+xy=y^2$ 是否为形如 $y'=\varphi(y/x)$ 的齐次方程。对符合者求 $x>0$ 上的通解，并补全常比例解。
\end{exercise}
\begin{exercise}{747}{根式齐次方程与漏解}
求 $y'=2\sqrt{y/x}+y/x$ 在 $x>0,y>0$ 上的通解，明确表达式成立的区间限制；另验证 $y\equiv0$ 为解。求满足 $y(1)=1$ 的正解及其保持为正的最大区间。
\end{exercise}
\begin{exercise}{748}{三角齐次方程与分支限制}
求 $y'=\sec(y/x)+y/x$ 在 $x>0$ 上的隐式通解。再求满足 $y(1)=0$ 的解在包含 $1$ 的区间上的显式表达式，并确定该区间端点，说明反三角函数分支的选择。
\end{exercise}
\begin{exercise}{749}{齐次初值问题与非零约束}
求 $y'=(x^2+y^2)/(xy)$、$y(1)=2$ 的特解。给出包含 $1$ 的最大存在区间，证明所选分支始终为正，并判断其能否以原方程解的形式跨越左端点。
\end{exercise}
\begin{exercise}{750}{切线斜率模型与有限区间}
曲线通过 $(1,0)$，切线斜率为 $1+y/x+(y/x)^2$。求曲线方程及包含 $x=1$ 的最大存在区间；说明方程的三角函数表达式为何不能跨过相邻极点继续当作同一解。
\end{exercise}
\begin{exercise}{751}{齐次代换与平衡比例解}
求 $y(x-y)\dd x+x^2\dd y=0$ 满足 $y(1)=1$ 的特解及包含 $1$ 的最大存在区间。求通解过程中，列出将 $y/x$ 视为常数时得到的解，并判断初值解是否属于其中之一。
\end{exercise}
\begin{exercise}{752}{对数齐次方程与正值定义域}
求 $xy'+y(\ln x-\ln y)=0$、$y(1)=e^3$ 的特解，明确对数要求的定义域。再将初值改为 $y(1)=a>0$，求解并讨论所得解在 $x>0$ 上的单调性。
\end{exercise}
\begin{exercise}{753}{交换因变量与指数齐次代换}
求 $(1+2e^{x/y})\dd x+2e^{x/y}(1-x/y)\dd y=0$ 的隐式通解，限定 $y\ne0$。利用 $u=x/y$ 或识别全微分完成推导，并通过微分核验结果。
\end{exercise}
\begin{exercise}{754}{齐次代换与三角隐式解}
求 $(2y\sin(x/y)+3x\cos(x/y))\dd y-3y\cos(x/y)\dd x=0$ 在 $y\ne0$ 上的通解。说明分离变量中除以 $\sin(x/y)$ 是否遗漏常比例解，列出所遗漏的情形。
\end{exercise}
\begin{exercise}{755}{齐次代换与不可跨越的奇点}
求微分方程 $xy'=y+x e^{y/x}$ 在 $x>0$ 上满足 $y(1)=0$ 的解，并确定包含 $1$ 的最大存在区间。证明所得解无法在该区间右端连续延拓为有限值。
\end{exercise}
\begin{exercise}{756}{齐次方程与代换遗漏解}
求 $xy'=y+y^2/x$ 在 $x>0$ 上的全部解，包括代换分离时可能遗漏的解；对初值 $y(1)=a$，确定解的表达式及包含 $1$ 的最大存在区间。
\end{exercise}
\subsection{一阶线性微分方程}
\begin{exercise}{757}{齐次线性方程与初值符号}
求 $y'+2xy=0$ 的通解。对初值 $y(0)=a$，求解在全实轴上的有界性、符号及 $x\to\pm\infty$ 时的极限，并讨论 $a=0$。
\end{exercise}
\begin{exercise}{758}{初值问题与系数奇点}
求 $y'=y\tan x$、$y(0)=1$ 的特解。确定包含 $0$ 的最大存在区间，并求该区间内的最小值，说明不能直接把同一表达式跨过 $\pi/2$ 作为初值解。
\end{exercise}
\begin{exercise}{759}{一阶线性方程辨识与积分形式解}
分别判断 $y'+y^2=0$、$(y')^2=\sqrt{xy}$、$y'+xy=e^xy+x^2$、$y'+\sin(xy)=1$ 是否为关于 $y$ 的一阶线性方程。将符合者写为标准形式，给出积分因子及满足 $y(0)=1$ 的积分形式解，不要求化为初等函数。
\end{exercise}
\begin{exercise}{760}{一阶线性方程与平衡解}
求 $y'+2xy=4x$ 的通解。对初值 $y(0)=a$，讨论解与直线 $y=2$ 的位置关系及两端极限；求恰为常数解的 $a$。
\end{exercise}
\begin{exercise}{761}{积分因子与有界性}
求 $y'+y\cos x=e^{-\sin x}$ 的通解。判断是否存在在整个实轴上有界的解，以及是否存在周期解，分别给出依据。
\end{exercise}
\begin{exercise}{762}{奇异系数与原方程延拓}
求 $(x+1)y'-y=e^x(x+1)^2$、$y(0)=2$ 的特解。先在 $x>-1$ 上求解，再检验所得函数能否延拓为全实轴上的原方程解；指出将方程除以 $x+1$ 会失去什么信息。
\end{exercise}
\begin{exercise}{763}{积分因子与对数外力}
求 $y'+y\cos x-e^{-\sin x}\ln x=0$、$y(1)=0$ 的特解。明确最大定义区间，判断 $x=1$ 处是否为极值点，并给出判断依据。
\end{exercise}
\begin{exercise}{764}{交换自因变量的线性化}
求 $(y^2-6x)y'+2y=0$ 的隐式通解。先将 $x$ 视为 $y$ 的函数并在线性化时注明 $y\ne0$，再单独检查 $y\equiv0$ 是否为原方程的解。
\end{exercise}
\begin{exercise}{765}{关于逆函数的一阶线性方程}
求 $y^3\dd x+(2xy^2-1)\dd y=0$ 的隐式通解。说明将 $x$ 看作 $y$ 的函数时的定义域限制，并补全可能遗漏的常值解；给出通过 $(0,1)$ 的局部分支。
\end{exercise}
\begin{exercise}{766}{隐式初值解与垂直切线}
求 $y\dd x+(x-3y^2)\dd y=0$、$y(1)=1$ 的特解，以隐式形式表示。求初始点附近 $y'(x)$，并确定该分支首次出现垂直切线时的点。
\end{exercise}
\begin{exercise}{767}{逆函数线性化与局部可解性}
求 $y'=1/(x+y)$ 的隐式通解。以 $x$ 为 $y$ 的函数求解，说明在什么点可局部恢复为 $y(x)$；再求通过 $(0,1)$ 的隐式特解及该点导数。
\end{exercise}
\begin{exercise}{768}{线性方程与全微分核验}
求 $y\dd x+(x+y)\dd y=0$ 的隐式通解，并检查 $y\equiv0$。对通过 $(1,1)$ 的分支，求 $y'(1)$，说明隐式解在该点附近能表示为函数的依据。
\end{exercise}
\begin{exercise}{769}{切线斜率与积分因子}
曲线 $y=f(x)$ 通过 $(0,1)$，任一点切线斜率为 $\sec x-f(x)\tan x$。求曲线方程及包含 $0$ 的最大存在区间，计算 $f'(0)$ 并判断该点是否为极值点。
\end{exercise}
\begin{exercise}{770}{积分方程与线性初值问题}
设连续函数 $f$ 满足 $f(x)=\int_0^xf(t)\dd t-x^2$。求 $f$，说明连续性为何足以保证所需可导性；再代回核验，并求 $f$ 在原点处最低阶非零泰勒项。
\end{exercise}
\begin{exercise}{771}{卷积换元与二次恢复}
设连续函数 $f$ 满足 $\int_0^xf(x-t)\dd t=\int_0^x(x-t)f(t)\dd t+e^{-x}-1$。通过换元和求导求 $f$，明确求导后获得的初值，并核验原积分等式。
\end{exercise}
\begin{exercise}{772}{积分因子与有界解}
求方程 $y'+2y=1+e^{-x}$ 的通解；判断是否存在在整个实轴上有界的解，证明你的判断。再求所有解在 $x\to+\infty$ 时的极限。
\end{exercise}
\begin{exercise}{773}{分段外力与连续拼接}
某系统满足 $y'+y=u(x)$、$y(0)=0$，其中 $u(x)=1$（$0\le x<1$）、$u(x)=0$（$x>1$）。分别求两个区间上的解，使 $y$ 在 $x=1$ 连续；判断 $y$ 在 $x=1$ 是否可导，并说明理由。
\end{exercise}
\subsection{可降阶的高阶微分方程}
\begin{exercise}{774}{连续积分与初值数目}
求 $y^{(4)}=x$ 的通解，说明四个积分常数独立。再求满足 $y(0)=y'(0)=y''(0)=0$、$y'''(0)=1$ 的特解，并核验四个条件。
\end{exercise}
\begin{exercise}{775}{连续积分与三角指数外力}
求 $y''=e^{3x}+\sin2x$ 的通解，再求满足 $y(0)=y'(0)=0$ 的解。分别对所得表达式求二阶导数，核验两种外力的积分系数与符号。
\end{exercise}
\begin{exercise}{776}{连续积分与指数多项式}
求 $y''=xe^x$ 的通解及满足 $y(0)=y'(0)=0$ 的特解。求该特解与 $xe^x$ 在 $x\to+\infty$ 时的比值极限，并说明积分常数是否影响该极限。
\end{exercise}
\begin{exercise}{777}{反三角原函数与初值性质}
求 $y''=1/(1+x^2)$ 的通解。再求 $y(0)=y'(0)=0$ 的解，证明它是偶函数、恒非负且严格凸，并确定唯一最小值点。
\end{exercise}
\begin{exercise}{778}{三阶积分与三个初值}
求 $y'''=e^{2x}$ 满足 $y(1)=y'(1)=y''(1)=0$ 的特解。将结果写成 $e^{2x}$ 减去其在 $x=1$ 的二次泰勒多项式的形式并给出系数，再证明解在 $x>1$ 时为正。
\end{exercise}
\begin{exercise}{779}{切线条件与曲线恢复}
曲线经过 $(0,1)$，在该点与 $y=x/2+1$ 相切，并满足 $y''=x$。求曲线方程，证明两项几何条件分别对应两个独立初值，并判断 $(0,1)$ 是否为拐点。
\end{exercise}
\begin{exercise}{780}{不显含因变量的降阶}
求 $y''=y'+x$ 的通解。确定初始斜率 $y'(0)$ 应取何值，才能使解不含 $e^x$ 项；在此条件下，再由 $y(0)=1$ 求唯一解并核验。
\end{exercise}
\begin{exercise}{781}{导数代换与三角分支}
求 $y''+(y')^2+1=0$ 的通解。再求满足 $y(0)=y'(0)=0$ 的特解及包含 $0$ 的最大存在区间，说明积分得到的对数绝对值与区间选择之间的关系。
\end{exercise}
\begin{exercise}{782}{奇异系数与端点有界性}
求 $xy''+3y'=0$ 在 $x>0$ 上的通解及满足 $y(1)=0,y'(1)=2$ 的特解。确定通解中在 $x\to0^+$ 时有界的全部解，说明化为标准形式时不能包含 $x=0$。
\end{exercise}
\begin{exercise}{783}{非线性速度方程与有限奇点}
求 $y''-2(y')^2=0$、$y(0)=0,y'(0)=-1$ 的特解。令 $p=y'$ 降阶并求包含 $0$ 的最大存在区间，判断能否在左端点连续延拓为有限值，并验证原初值。
\end{exercise}
\begin{exercise}{784}{速度线性方程与单调性}
求 $(1+x^2)y''=2xy'$、$y(0)=1,y'(0)=3$ 的特解。给出降阶后速度方程的积分因子，判断解在全实轴上的单调性及零点个数，并用方程核验。
\end{exercise}
\begin{exercise}{785}{不显含自变量与常值解}
求 $yy''+2(y')^2=0$ 在 $y\ne0$ 区间上的通解，并检查常值解是否包含其中。采用 $p=y'$ 降阶，或利用 $y^3$ 的二阶导数求解；说明非常值解为何不能以可导方式跨越其零点。
\end{exercise}
\begin{exercise}{786}{乘积求导与正值初始分支}
求 $yy''+(y')^2=0$、$y(0)=1,y'(0)=1/2$ 的特解。由 $(yy')'=0$ 降阶，确定包含 $0$ 的最大存在区间，并说明平方后求解时如何选取符号。
\end{exercise}
\begin{exercise}{787}{不显含自变量的降阶与能量式}
求 $y^3y''=1$ 在 $y\ne0$ 区间上的隐式通解，先令 $p=y'$ 推导一阶积分。再求 $y(0)=1$、$y'(0)=0$ 的特解，确定其在全实轴上的正负与最小值。
\end{exercise}
\begin{exercise}{788}{不显含因变量的降阶与初值}
求 $y''+(y')^2=1$、$y(0)=y'(0)=0$ 的特解。令 $p=y'$ 降阶后，分别验证求得的速度和原函数；讨论解在全实轴上的奇偶性与凹凸性。
\end{exercise}
\begin{exercise}{789}{不显含因变量的降阶}
求初值问题 $xy''+y'=x$、$y(1)=0$、$y'(1)=1$ 在 $x>0$ 上的解，判断该解能否以二阶可导方式延拓到 $x=0$，并给出依据。
\end{exercise}
\begin{exercise}{790}{不显含自变量与能量积分}
设 $y''=-y^3$、$y(0)=a>0$、$y'(0)=0$。通过降阶证明 $\frac12(y')^2+\frac14y^4=\frac14a^4$；在解首次降到 $0$ 之前确定 $y'$ 的符号，并用定积分表示首次到达 $0$ 的时间，证明该时间有限。
\end{exercise}
\subsection{高阶线性微分方程}
\begin{exercise}{791}{函数组线性相关性}
在实轴上分别判断 $\{x^2,2x^2\}$、$\{\sin x,\cos x\}$、$\{\sin2x,\sin x\cos x\}$、$\{e^{2x},3e^{2x}\}$ 的线性相关性。对无关者给出证明；说明将定义区间缩小为任意非空开区间是否改变结论。
\end{exercise}
\begin{exercise}{792}{已知解验证与基本解组}
对方程 $(x-1)y''-xy'+y=0$，分别判断 $\{x,2x\}$、$\{e^x,8e^x\}$、$\{x,x^2\}$、$\{x,e^x\}$ 能否构成 $x>1$ 上的基本解组。求通解，并说明为何需要把“两个解”与“两个线性无关的解”分开检验。
\end{exercise}
\begin{exercise}{793}{一阶非齐次方程的解结构}
设 $p,f$ 在区间 $I$ 上连续，$y_1$ 是 $y'+py=f$ 的一个解，$y_2\ne0$ 是对应齐次方程的解。证明全部解恰为 $y=y_1+Cy_2$，并用初值 $y(x_0)=a$ 表示 $C$。说明若误把 $C$ 乘在 $y_1$ 前会产生什么条件限制。
\end{exercise}
\begin{exercise}{794}{一阶齐次解与冗余常数}
设连续系数方程 $y'+p(x)y=0$ 在同一区间上有两个不同解 $y_1,y_2$。证明 $y=C(y_1-y_2)$ 为通解。判断 $C_1y_1+C_2y_2$ 中两常数是否独立，并讨论其中一个已知解为零函数的情形。
\end{exercise}
\begin{exercise}{795}{朗斯基行列式与通解判据}
设 $f_1,f_2$ 是 $y''+py'+qy=0$ 的两个非零解，$p,q$ 为常数。记 $W=f_1f_2'-f_2f_1'$。证明若某点 $W\ne0$，则 $C_1f_1+C_2f_2$ 是通解；推导 $W'=-pW$，并说明为什么“两个不同的非零解”本身还不足够。
\end{exercise}
\begin{exercise}{796}{三角基本解组与退化参数}
验证 $\cos ax,\sin ax$ 是 $y''+a^2y=0$（$a\ne0$）的线性无关解并写出通解。若 $a=0$，说明该函数组为何失效并另写通解；求两种情况下满足 $y(0)=0,y'(0)=1$ 的解。
\end{exercise}
\begin{exercise}{797}{非齐次通解的完整性验证}
验证 $y=C_1e^x+C_2e^{2x}+\frac1{12}e^{5x}$ 是 $y''-3y'+2y=e^{5x}$ 的通解。除代入验证外，还须证明齐次部分线性无关；再求满足 $y(0)=y'(0)=0$ 的特解。
\end{exercise}
\begin{exercise}{798}{三个非齐次解构造通解}
设 $p,q,f$ 在区间 $I$ 上连续，$y_1,y_2,y_3$ 是 $y''+py'+qy=f$ 的三个线性无关的解。证明
\[
y=2C_1y_1+(C_2-C_1)y_2+(1-C_1-C_2)y_3
\]
为通解，并说明两个齐次方向的独立性。再证明此假设蕴含 $f\not\equiv0$。
\end{exercise}
\begin{exercise}{799}{解的叠加与初始条件}
设 $y_1,y_2$ 是方程 $y''+p(x)y'+q(x)y=f(x)$ 的两个解。证明 $\alpha y_1+(1-\alpha)y_2$ 仍为解，并说明一般线性组合成为解所需的条件。取 $f\not\equiv0$ 时，判断 $y_1+y_2$ 与 $y_1-y_2$ 分别满足什么方程。
\end{exercise}
\begin{exercise}{800}{已知解降阶与基本解组}
已知 $y_1=x$ 是 $x^2y''-2xy'+2y=0$ 在 $x>0$ 上的解。令 $y=xv(x)$ 求通解，证明所得两解线性无关，并求满足 $y(1)=0$、$y'(1)=1$ 的解。
\end{exercise}
\subsection{常系数齐次线性微分方程}
\begin{exercise}{801}{候选特解验证与有界解}
分别判断 $e^x$、$e^{2x}$、$1+e^{-2x}$、$xe^{-2x}$ 是否为 $y''+2y'=0$ 的解，给出代入依据。求此方程通解，并确定全实轴上的全部有界解。
\end{exercise}
\begin{exercise}{802}{由解族恢复方程}
求以 $y=C_1e^{6x}+C_2e^{-8x}$ 为通解的首项系数为 $1$ 的二阶常系数齐次方程，证明系数唯一。确定通解中在 $x\to+\infty$ 时趋于零的全部解。
\end{exercise}
\begin{exercise}{803}{共轭特征根与方程恢复}
求以 $e^{2x}\sin3x,e^{2x}\cos3x$ 为基本解组的首项系数为 $1$ 的二阶常系数齐次方程。求满足 $y(0)=0,y'(0)=1$ 的特解，并判断其在 $[0,+\infty)$ 上是否有界。
\end{exercise}
\begin{exercise}{804}{已知指数解与系数信息充分性}
已知 $e^{2x}$ 是 $y''+py'+qy=0$ 的解，求所有可能的 $(p,q)$，说明能否唯一确定系数。若再知 $xe^{2x}$ 也是解，求 $p,q$ 及通解，并解释新增信息与特征根重数的关系。
\end{exercise}
\begin{exercise}{805}{不同实根与初值问题}
求 $y''+5y'+4y=0$ 的通解。再求满足 $y(0)=1,y'(0)=0$ 的解，证明该解在 $x\ge0$ 上恒正且不增加，并求其极限。
\end{exercise}
\begin{exercise}{806}{初值解与极值}
求 $y''+4y'+3y=0$、$y(0)=2,y'(0)=6$ 的特解。确定它在 $[0,+\infty)$ 上的零点及极值，并求最大值。
\end{exercise}
\begin{exercise}{807}{重复特征根与独立解}
求 $4y''-12y'+9y=0$ 的通解，证明所取两个解线性无关。再确定全部满足 $\lim_{x\to+\infty}y(x)=0$ 的解，并说明只消去指数项中的一个系数为何不够。
\end{exercise}
\begin{exercise}{808}{临界阻尼型初值问题}
求 $s''+2s'+s=0$（自变量为 $t$），$s(0)=2,s'(0)=-1$ 的解。证明它在 $t\ge0$ 上恒正且严格减少，求其极限，并说明极限值是否在有限时刻取到。
\end{exercise}
\begin{exercise}{809}{复特征根与零点间隔}
求 $y''+8y'+17y=0$ 的通解。证明任意非零解在 $x\ge0$ 上有无穷多个零点，求相邻零点间距，并判断解在 $x\to+\infty$ 时是否趋于零。
\end{exercise}
\begin{exercise}{810}{衰减振荡与初值确定}
求 $y''+10y'+29y=0$、$y(0)=1,y'(0)=-1$ 的特解。把三角部分化为单一余弦形式，确定其振幅包络，并求第一个正零点。
\end{exercise}
\begin{exercise}{811}{三阶常系数方程与重根}
求 $y'''+y''-y'-y=0$ 的通解，先将特征多项式分组因式分解。证明三个基本解线性无关，并确定全部在 $x\to+\infty$ 时趋于零的解。
\end{exercise}
\begin{exercise}{812}{四阶方程与实复根组合}
求 $y^{(4)}-y'''-y''-y'-2y=0$ 的通解，证明特征多项式可分解为 $(r^2+1)(r-2)(r+1)$。确定通解中在 $[0,+\infty)$ 上有界的解，以及其中趋于零的解。
\end{exercise}
\begin{exercise}{813}{反射导数关系与降阶恢复}
设 $f$ 在实轴上可导，满足 $f'(x)=f(-x)$、$f(0)=1$。证明 $f$ 二阶可导并求 $f$，再代回反射关系核验，说明只满足导出的二阶方程为何还不足以完成验证。
\end{exercise}
\begin{exercise}{814}{两个微分关系的相容性}
设 $f$ 同时满足 $f''+f'-2f=0$ 与 $f''+f=2e^x$。求 $f$，先分别给出两个方程解的形式，再确定公共解，并证明其唯一性。
\end{exercise}
\begin{exercise}{815}{重根、参数与长期行为}
求初值问题 $y''+2ay'+a^2y=0$、$y(0)=0$、$y'(0)=1$ 的解。讨论实参数 $a$ 使解在 $[0,+\infty)$ 上有界或趋于零的条件，并在 $a>0$ 时求其最大值。
\end{exercise}
\begin{exercise}{816}{阻尼分类与无零点条件}
考虑 $y''+2\zeta y'+y=0$，$y(0)=1$、$y'(0)=0$，其中 $\zeta\ge0$。分别求 $0\le\zeta<1$、$\zeta=1$、$\zeta>1$ 时的解，并确定解在 $x>0$ 上始终为正的参数范围。
\end{exercise}
\subsection{常系数非齐次线性微分方程}
\begin{exercise}{817}{多项式特解与通解}
求 $y''-4y'+3y=x$ 的一个一次多项式特解及通解。证明一次多项式特解唯一，并求满足 $y(0)=y'(0)=0$ 的特解。
\end{exercise}
\begin{exercise}{818}{多项式外力与长期行为}
求 $y''+5y'+4y=3-2x$ 的通解。证明所有解与同一个一次多项式之差在 $x\to+\infty$ 时趋于零，并写出该多项式。
\end{exercise}
\begin{exercise}{819}{零特征根与待定系数形式}
求 $2y''+5y'=5x^2-2x-1$ 的通解。先说明特解为何需要三次多项式；再求满足 $y(0)=y'(0)=0$ 的特解，并核验常数项可并入齐次解。
\end{exercise}
\begin{exercise}{820}{由齐次通解恢复系数与求特解}
已知 $y''+ay'+by=0$ 的通解为 $(C_1+C_2x)e^x$。求 $a,b$，再求 $y''+ay'+by=x$ 满足 $y(0)=2,y'(0)=0$ 的解，分别核验方程及初值。
\end{exercise}
\begin{exercise}{821}{混合外力与共振特解形式}
对 $y''-5y'+6y=3x+2e^{2x}$，说明为何可设特解为 $Ax+B+Cxe^{2x}$，求出系数并写出通解。解释指数项前乘 $x$ 的原因，以及单设 $Ce^{2x}$ 为什么无效。
\end{exercise}
\begin{exercise}{822}{指数多项式与单根共振}
求 $y''-3y'+2y=xe^{2x}$ 的通解。先写出最小的待定系数特解形式，再求满足 $y(0)=y'(0)=0$ 的特解，并用代换 $y=e^{2x}v$ 核验。
\end{exercise}
\begin{exercise}{823}{已知非齐次解反求常系数}
已知 $y=e^{-2x}+(x^2+2)e^x$ 是 $y''+ay'+by=(cx+d)e^x$ 的解，其中 $a,b,c,d$ 为常数。求全部系数并证明其唯一性，再写出该方程通解。
\end{exercise}
\begin{exercise}{824}{三个解恢复方程与外力}
某首项系数为 $1$ 的二阶常系数非齐次方程有三个解 $y_1=e^{2x}+xe^x$、$y_2=e^{-x}+xe^x$、$y_3=e^{2x}-e^{-x}+xe^x$。求方程及通解，并说明如何先从解的差恢复齐次部分，再恢复外力项。
\end{exercise}
\begin{exercise}{825}{复根共振与候选解验证}
对 $y''-2y'+5y=e^x\cos2x$，分别检验 $xe^x(\cos2x+\sin2x)$、$e^x(\cos2x+\sin2x)$、$\frac14xe^x\cos2x$、$\frac14xe^x\sin2x$ 是否为特解。选出成立者，写出通解，并说明共振所需的 $x$ 因子。
\end{exercise}
\begin{exercise}{826}{倍角分解与三角共振}
求 $y''+4y=\cos^2x$ 的一个特解及通解。先把外力分解为常数项与三角项，分别说明其待定系数形式；再求满足 $y(0)=y'(0)=0$ 的解。
\end{exercise}
\begin{exercise}{827}{多项式三角外力}
求 $y''+4y=x\cos x$ 的通解。说明仅设特解为 $Ax\cos x$ 是否足够，并写出包含必要正弦项的特解形式；再求满足 $y(0)=y'(0)=0$ 的特解。
\end{exercise}
\begin{exercise}{828}{外力叠加与指数共振}
求 $y''-y=e^x+4\cos x$ 的通解。分别处理两种外力并说明指数项的共振，再求满足 $y(0)=y'(0)=0$ 的特解，判断其在 $[0,+\infty)$ 上是否有界。
\end{exercise}
\begin{exercise}{829}{共振与参数分类}
对 $\omega>0$，求初值问题 $y''+y=\cos\omega x$、$y(0)=y'(0)=0$ 的解，分别讨论 $\omega\ne1$ 与 $\omega=1$。判断解有界的参数条件，并说明共振时无界项的来源。
\end{exercise}
\begin{exercise}{830}{指数多项式外力与特解形式}
求 $y''-2y'+y=e^x(x+1)$ 的通解及满足 $y(0)=y'(0)=0$ 的特解。先写出待定系数法所需的特解形式，并用 $y=e^xv$ 的代换核验计算结果。
\end{exercise}

\end{document}
